\documentclass[11pt,a4paper,oneside]{report}
% ======================================================================
%  Преамбюл на доклада (XeLaTeX)
% ======================================================================
\usepackage[a4paper,top=2.6cm,bottom=2.6cm,left=2.7cm,right=2.4cm,headheight=15pt]{geometry}
\usepackage{amsmath,amssymb,mathtools,bm}
\usepackage[no-math]{fontspec}
\usepackage{polyglossia}
\setdefaultlanguage{bulgarian}
\setotherlanguage{english}
\setmainfont{PT Serif}
\setsansfont{PT Sans}
\setmonofont{DejaVu Sans Mono}[Scale=0.84]
\newfontfamily\cyrillicfont{PT Serif}
\newfontfamily\cyrillicfontsf{PT Sans}
\newfontfamily\cyrillicfonttt{DejaVu Sans Mono}[Scale=0.84]

\usepackage{xcolor}
\definecolor{navy}{HTML}{1F2A44}
\definecolor{accent}{HTML}{2A78D6}
\definecolor{orange}{HTML}{EB6834}
\definecolor{attack}{HTML}{E34948}
\definecolor{muted}{HTML}{52514E}
\definecolor{rule}{HTML}{C9CDD6}
\definecolor{boxbg}{HTML}{F2F5FA}
\definecolor{rowalt}{HTML}{F5F7FA}

\usepackage{graphicx}
\usepackage{tikz}
\usetikzlibrary{arrows.meta,positioning,calc}
\usepackage{booktabs,tabularx,makecell,multirow,array,longtable}
\usepackage[table]{xcolor}
\usepackage{caption,subcaption}
\captionsetup{font={small,sf},labelfont={bf,color=navy},labelsep=period,justification=justified,singlelinecheck=true,skip=6pt}
\captionsetup[table]{position=top}
\usepackage{siunitx}
\sisetup{output-decimal-marker={,},group-separator={\,},group-minimum-digits=4}
\usepackage{enumitem}
\setlist{itemsep=2pt,topsep=4pt,leftmargin=1.6em}
\usepackage{microtype}\setlength{\emergencystretch}{2.5em}
\usepackage{setspace}
\setstretch{1.12}
\setlength{\parskip}{3pt}
\usepackage{float}
\usepackage{adjustbox}\newcommand{\fitbox}[1]{\adjustbox{max width=\linewidth}{#1}}
\usepackage{placeins}\usepackage{needspace}\usepackage{etoolbox}

% --- непрекъсната номерация на формули, фигури и таблици в целия документ ---
\counterwithout{equation}{chapter}
\counterwithout{figure}{chapter}
\counterwithout{table}{chapter}

% --- заглавия ------------------------------------------------------------
\usepackage{titlesec}
\titleformat{\chapter}[display]
  {\sffamily\color{navy}}
  {\normalsize\bfseries\color{accent}\MakeUppercase{\chaptertitlename}\ \thechapter}
  {4pt}
  {\Huge\bfseries}[\vspace{4pt}{\color{accent}\titlerule[1.2pt]}]
\titleformat{name=\chapter,numberless}[display]
  {\sffamily\color{navy}}{}{0pt}{\Huge\bfseries}[\vspace{4pt}{\color{accent}\titlerule[1.2pt]}]
\titlespacing*{\chapter}{0pt}{-10pt}{22pt}
\titleformat{\section}{\sffamily\Large\bfseries\color{navy}}{\color{accent}\thesection}{0.7em}{}
\titleformat{\subsection}{\sffamily\large\bfseries\color{navy}}{\color{accent}\thesubsection}{0.7em}{}
\titleformat{\subsubsection}{\sffamily\normalsize\bfseries\color{navy}}{\thesubsubsection}{0.6em}{}
\titlespacing*{\section}{0pt}{16pt}{7pt}
\titlespacing*{\subsection}{0pt}{12pt}{5pt}
\setcounter{secnumdepth}{3}
\setcounter{tocdepth}{2}

% --- колонтитули -----------------------------------------------------------
\usepackage{fancyhdr}
\pagestyle{fancy}
\fancyhf{}
\renewcommand{\headrulewidth}{0.4pt}
\renewcommand{\headrule}{{\color{rule}\hrule width\headwidth height\headrulewidth}}
\fancyhead[L]{\sffamily\footnotesize\color{muted}ML откриване на прониквания в трафика на център за данни}
\fancyhead[R]{\sffamily\footnotesize\color{muted}\nouppercase{\leftmark}}
\fancyfoot[C]{\sffamily\footnotesize\color{navy}\thepage}
\fancypagestyle{plain}{\fancyhf{}\renewcommand{\headrulewidth}{0pt}\fancyfoot[C]{\sffamily\footnotesize\color{navy}\thepage}}
\renewcommand{\chaptermark}[1]{\markboth{\thechapter.\ #1}{}}

% --- кутии -----------------------------------------------------------------
\usepackage[most]{tcolorbox}
\newtcolorbox{keybox}[1][Основни изводи]{enhanced,breakable,colback=boxbg,colframe=accent,boxrule=0pt,
  leftrule=3pt,arc=1pt,left=8pt,right=8pt,top=5pt,bottom=5pt,fonttitle=\sffamily\bfseries\small,
  coltitle=navy,colbacktitle=boxbg,title=#1,attach title to upper={\par\smallskip},fontupper=\small}
\newtcolorbox{notebox}[1][Бележка]{enhanced,breakable,colback=white,colframe=orange,boxrule=0.6pt,arc=1pt,
  left=8pt,right=8pt,top=4pt,bottom=4pt,fonttitle=\sffamily\bfseries\small,coltitle=orange,
  colbacktitle=white,title=#1,attach title to upper={\par\smallskip},fontupper=\small}
\newtcolorbox{interp}{enhanced,colback=white,colframe=white,boxrule=0pt,borderline west={2pt}{0pt}{rule},
  left=10pt,right=4pt,top=2pt,bottom=2pt,fontupper=\small}
\newcommand{\interpret}[1]{\begin{interp}\textsf{\textbf{\color{navy}Какво показва фигурата.}}\ #1\end{interp}}
\newcommand{\interprettab}[1]{\begin{interp}\textsf{\textbf{\color{navy}Как се чете таблицата.}}\ #1\end{interp}}

% --- код -------------------------------------------------------------------
\usepackage{listings}
\lstdefinestyle{py}{language=Python,basicstyle=\ttfamily\footnotesize,keywordstyle=\color{accent}\bfseries,
  commentstyle=\color{muted},stringstyle=\color{orange},showstringspaces=false,frame=l,
  framerule=2pt,rulecolor=\color{accent},backgroundcolor=\color{boxbg},xleftmargin=12pt,framexleftmargin=8pt,
  breaklines=true,columns=fullflexible,extendedchars=true}
\lstset{style=py}

% --- таблици ---------------------------------------------------------------
\newcolumntype{R}{>{\raggedleft\arraybackslash}X}
\newcolumntype{L}{>{\raggedright\arraybackslash}X}
\newcommand{\hd}[1]{\makecell{\sffamily\bfseries\color{white}#1}}
\newcommand{\headrow}{\rowcolor{navy}}
\renewcommand{\arraystretch}{1.18}

% --- числа от резултатите (генерирани автоматично) ------------------------
\newcommand{\val}[1]{\ifcsname val@#1\endcsname\csname val@#1\endcsname\else\textbf{\color{red}??#1??}\fi}

% --- математика ------------------------------------------------------------
\newcommand{\R}{\mathbb{R}}
\newcommand{\E}{\mathbb{E}}
\newcommand{\vect}[1]{\mathbf{#1}}
\DeclareMathOperator*{\argmax}{arg\,max}
\DeclareMathOperator*{\argmin}{arg\,min}
\DeclareMathOperator{\sign}{sign}
\newcommand{\txt}[1]{\text{\normalfont #1}}

\usepackage[hidelinks]{hyperref}
\hypersetup{colorlinks=true,linkcolor=navy,citecolor=accent,urlcolor=accent,
  pdftitle={Експериментално изследване на откриване на проникване в мрежов трафик на център за данни, базирано на машинно обучение},
  pdfauthor={Kristiyan Yanutsev}}
\usepackage[english,nameinlink,noabbrev]{cleveref}
\crefname{equation}{формула}{формули}
\crefname{figure}{фиг.}{фиг.}
\crefname{table}{табл.}{табл.}
\crefname{section}{раздел}{раздели}
\crefname{chapter}{глава}{глави}
\crefname{appendix}{приложение}{приложения}
\creflabelformat{equation}{(#2#1#3)}
\pretocmd{\section}{\Needspace{20\baselineskip}}{}{}\pretocmd{\subsection}{\Needspace{12\baselineskip}}{}{}

% Без курсив никъде в текста
\renewcommand{\emph}[1]{#1}
\renewcommand{\textit}[1]{#1}
\renewcommand{\textsl}[1]{#1}
\AtBeginDocument{\renewcommand{\itshape}{\upshape}\renewcommand{\slshape}{\upshape}}

% Автоматично генерирано от report/build_report.py - не редактирайте ръчно
\expandafter\def\csname val@RF-Bot-AUC\endcsname{1,0000}
\expandafter\def\csname val@RF-Bot-AUC-5\endcsname{1,00000}
\expandafter\def\csname val@RF-Bot-Accuracy\endcsname{99,97\%}
\expandafter\def\csname val@RF-Bot-Accuracy-100\endcsname{100}
\expandafter\def\csname val@RF-Bot-Accuracy-3\endcsname{99,967\%}
\expandafter\def\csname val@RF-Bot-Accuracy-hi\endcsname{99,98\%}
\expandafter\def\csname val@RF-Bot-Accuracy-hw\endcsname{0,01}
\expandafter\def\csname val@RF-Bot-Accuracy-lo\endcsname{99,95\%}
\expandafter\def\csname val@RF-Bot-F1\endcsname{99,97\%}
\expandafter\def\csname val@RF-Bot-F1-3\endcsname{99,967\%}
\expandafter\def\csname val@RF-Bot-F1-hi\endcsname{99,98\%}
\expandafter\def\csname val@RF-Bot-F1-hw\endcsname{0,01}
\expandafter\def\csname val@RF-Bot-F1-lo\endcsname{99,95\%}
\expandafter\def\csname val@RF-Bot-FN\endcsname{12}
\expandafter\def\csname val@RF-Bot-FNR\endcsname{0,04\%}
\expandafter\def\csname val@RF-Bot-FNR-3\endcsname{0,040\%}
\expandafter\def\csname val@RF-Bot-FP\endcsname{8}
\expandafter\def\csname val@RF-Bot-FPR\endcsname{0,03\%}
\expandafter\def\csname val@RF-Bot-FPR-3\endcsname{0,027\%}
\expandafter\def\csname val@RF-Bot-Precision\endcsname{99,97\%}
\expandafter\def\csname val@RF-Bot-Precision-100\endcsname{100}
\expandafter\def\csname val@RF-Bot-Precision-3\endcsname{99,973\%}
\expandafter\def\csname val@RF-Bot-Precision-hi\endcsname{99,99\%}
\expandafter\def\csname val@RF-Bot-Precision-hw\endcsname{0,02}
\expandafter\def\csname val@RF-Bot-Precision-lo\endcsname{99,95\%}
\expandafter\def\csname val@RF-Bot-ROC_AUC-hi\endcsname{100,00\%}
\expandafter\def\csname val@RF-Bot-ROC_AUC-hw\endcsname{0,00}
\expandafter\def\csname val@RF-Bot-ROC_AUC-lo\endcsname{100,00\%}
\expandafter\def\csname val@RF-Bot-Recall-100\endcsname{100}
\expandafter\def\csname val@RF-Bot-Sensitivity\endcsname{99,96\%}
\expandafter\def\csname val@RF-Bot-Sensitivity-3\endcsname{99,960\%}
\expandafter\def\csname val@RF-Bot-Sensitivity-hi\endcsname{99,98\%}
\expandafter\def\csname val@RF-Bot-Sensitivity-hw\endcsname{0,02}
\expandafter\def\csname val@RF-Bot-Sensitivity-lo\endcsname{99,94\%}
\expandafter\def\csname val@RF-Bot-Specificity\endcsname{99,97\%}
\expandafter\def\csname val@RF-Bot-Specificity-100\endcsname{100}
\expandafter\def\csname val@RF-Bot-Specificity-3\endcsname{99,973\%}
\expandafter\def\csname val@RF-Bot-Specificity-hi\endcsname{99,99\%}
\expandafter\def\csname val@RF-Bot-Specificity-hw\endcsname{0,02}
\expandafter\def\csname val@RF-Bot-Specificity-lo\endcsname{99,95\%}
\expandafter\def\csname val@RF-Bot-TN\endcsname{29\,992}
\expandafter\def\csname val@RF-Bot-TP\endcsname{29\,988}
\expandafter\def\csname val@RF-Bot-cv\endcsname{99,93\%}
\expandafter\def\csname val@RF-Bot-cv-std\endcsname{0,03\%}
\expandafter\def\csname val@RF-Bot-cv-time\endcsname{226}
\expandafter\def\csname val@RF-Bot-dF1\endcsname{0,03}
\expandafter\def\csname val@RF-Bot-fit\endcsname{6,3}
\expandafter\def\csname val@RF-Bot-full-F1\endcsname{99,938\%}
\expandafter\def\csname val@RF-Bot-full-fit\endcsname{4,3}
\expandafter\def\csname val@RF-Bot-full-pred-us\endcsname{6,0}
\expandafter\def\csname val@RF-Bot-gap\endcsname{0,016}
\expandafter\def\csname val@RF-Bot-nat-Accuracy\endcsname{100,0\%}
\expandafter\def\csname val@RF-Bot-nat-F1\endcsname{99,8\%}
\expandafter\def\csname val@RF-Bot-nat-FN\endcsname{14}
\expandafter\def\csname val@RF-Bot-nat-FP\endcsname{73}
\expandafter\def\csname val@RF-Bot-nat-Precision\endcsname{99,7\%}
\expandafter\def\csname val@RF-Bot-nat-Sensitivity\endcsname{99,9\%}
\expandafter\def\csname val@RF-Bot-nat-Specificity\endcsname{100,0\%}
\expandafter\def\csname val@RF-Bot-nat-TN\endcsname{172\,369}
\expandafter\def\csname val@RF-Bot-nat-TP\endcsname{27\,544}
\expandafter\def\csname val@RF-Bot-nat-fp-per-tp\endcsname{0,00}
\expandafter\def\csname val@RF-Bot-nat-n\endcsname{200\,000}
\expandafter\def\csname val@RF-Bot-noiat-F1\endcsname{99,958\%}
\expandafter\def\csname val@RF-Bot-noiat-dF1\endcsname{-0,01}
\expandafter\def\csname val@RF-Bot-ntrain\endcsname{140\,000}
\expandafter\def\csname val@RF-Bot-params\endcsname{max\_depth=25, max\_features=sqrt, min\_samples\_leaf=2, n\_estimators=200}
\expandafter\def\csname val@RF-Bot-pca-F1\endcsname{99,948\%}
\expandafter\def\csname val@RF-Bot-pca-dF1\endcsname{0,02}
\expandafter\def\csname val@RF-Bot-pred-us\endcsname{6,6}
\expandafter\def\csname val@RF-Bot-speedup\endcsname{0,91}
\expandafter\def\csname val@RF-Bot-throughput\endcsname{151\,108}
\expandafter\def\csname val@RF-Bot-train-acc\endcsname{99,983\%}
\expandafter\def\csname val@RF-BruteForce-AUC\endcsname{1,0000}
\expandafter\def\csname val@RF-BruteForce-AUC-5\endcsname{1,00000}
\expandafter\def\csname val@RF-BruteForce-Accuracy\endcsname{100,00\%}
\expandafter\def\csname val@RF-BruteForce-Accuracy-100\endcsname{100}
\expandafter\def\csname val@RF-BruteForce-Accuracy-3\endcsname{99,998\%}
\expandafter\def\csname val@RF-BruteForce-Accuracy-hi\endcsname{100,00\%}
\expandafter\def\csname val@RF-BruteForce-Accuracy-hw\endcsname{0,00}
\expandafter\def\csname val@RF-BruteForce-Accuracy-lo\endcsname{99,99\%}
\expandafter\def\csname val@RF-BruteForce-F1\endcsname{100,00\%}
\expandafter\def\csname val@RF-BruteForce-F1-3\endcsname{99,998\%}
\expandafter\def\csname val@RF-BruteForce-F1-hi\endcsname{100,00\%}
\expandafter\def\csname val@RF-BruteForce-F1-hw\endcsname{0,00}
\expandafter\def\csname val@RF-BruteForce-F1-lo\endcsname{99,99\%}
\expandafter\def\csname val@RF-BruteForce-FN\endcsname{1}
\expandafter\def\csname val@RF-BruteForce-FNR\endcsname{0,00\%}
\expandafter\def\csname val@RF-BruteForce-FNR-3\endcsname{0,004\%}
\expandafter\def\csname val@RF-BruteForce-FP\endcsname{0}
\expandafter\def\csname val@RF-BruteForce-FPR\endcsname{0,00\%}
\expandafter\def\csname val@RF-BruteForce-FPR-3\endcsname{0,000\%}
\expandafter\def\csname val@RF-BruteForce-Precision\endcsname{100,00\%}
\expandafter\def\csname val@RF-BruteForce-Precision-100\endcsname{100}
\expandafter\def\csname val@RF-BruteForce-Precision-3\endcsname{100,000\%}
\expandafter\def\csname val@RF-BruteForce-Precision-hi\endcsname{100,00\%}
\expandafter\def\csname val@RF-BruteForce-Precision-hw\endcsname{0,00}
\expandafter\def\csname val@RF-BruteForce-Precision-lo\endcsname{100,00\%}
\expandafter\def\csname val@RF-BruteForce-ROC_AUC-hi\endcsname{100,00\%}
\expandafter\def\csname val@RF-BruteForce-ROC_AUC-hw\endcsname{0,00}
\expandafter\def\csname val@RF-BruteForce-ROC_AUC-lo\endcsname{100,00\%}
\expandafter\def\csname val@RF-BruteForce-Recall-100\endcsname{100}
\expandafter\def\csname val@RF-BruteForce-Sensitivity\endcsname{100,00\%}
\expandafter\def\csname val@RF-BruteForce-Sensitivity-3\endcsname{99,996\%}
\expandafter\def\csname val@RF-BruteForce-Sensitivity-hi\endcsname{100,00\%}
\expandafter\def\csname val@RF-BruteForce-Sensitivity-hw\endcsname{0,01}
\expandafter\def\csname val@RF-BruteForce-Sensitivity-lo\endcsname{99,99\%}
\expandafter\def\csname val@RF-BruteForce-Specificity\endcsname{100,00\%}
\expandafter\def\csname val@RF-BruteForce-Specificity-100\endcsname{100}
\expandafter\def\csname val@RF-BruteForce-Specificity-3\endcsname{100,000\%}
\expandafter\def\csname val@RF-BruteForce-Specificity-hi\endcsname{100,00\%}
\expandafter\def\csname val@RF-BruteForce-Specificity-hw\endcsname{0,00}
\expandafter\def\csname val@RF-BruteForce-Specificity-lo\endcsname{100,00\%}
\expandafter\def\csname val@RF-BruteForce-TN\endcsname{28\,229}
\expandafter\def\csname val@RF-BruteForce-TP\endcsname{28\,227}
\expandafter\def\csname val@RF-BruteForce-cv\endcsname{99,99\%}
\expandafter\def\csname val@RF-BruteForce-cv-std\endcsname{0,01\%}
\expandafter\def\csname val@RF-BruteForce-cv-time\endcsname{189}
\expandafter\def\csname val@RF-BruteForce-dF1\endcsname{0,00}
\expandafter\def\csname val@RF-BruteForce-fit\endcsname{5,1}
\expandafter\def\csname val@RF-BruteForce-full-F1\endcsname{99,996\%}
\expandafter\def\csname val@RF-BruteForce-full-fit\endcsname{3,7}
\expandafter\def\csname val@RF-BruteForce-full-pred-us\endcsname{6,0}
\expandafter\def\csname val@RF-BruteForce-gap\endcsname{0,000}
\expandafter\def\csname val@RF-BruteForce-nat-Accuracy\endcsname{100,0\%}
\expandafter\def\csname val@RF-BruteForce-nat-F1\endcsname{100,0\%}
\expandafter\def\csname val@RF-BruteForce-nat-FN\endcsname{1}
\expandafter\def\csname val@RF-BruteForce-nat-FP\endcsname{4}
\expandafter\def\csname val@RF-BruteForce-nat-Precision\endcsname{100,0\%}
\expandafter\def\csname val@RF-BruteForce-nat-Sensitivity\endcsname{100,0\%}
\expandafter\def\csname val@RF-BruteForce-nat-Specificity\endcsname{100,0\%}
\expandafter\def\csname val@RF-BruteForce-nat-TN\endcsname{132\,693}
\expandafter\def\csname val@RF-BruteForce-nat-TP\endcsname{28\,227}
\expandafter\def\csname val@RF-BruteForce-nat-fp-per-tp\endcsname{0,00}
\expandafter\def\csname val@RF-BruteForce-nat-n\endcsname{160\,925}
\expandafter\def\csname val@RF-BruteForce-noiat-F1\endcsname{99,998\%}
\expandafter\def\csname val@RF-BruteForce-noiat-dF1\endcsname{0,00}
\expandafter\def\csname val@RF-BruteForce-ntrain\endcsname{131\,731}
\expandafter\def\csname val@RF-BruteForce-params\endcsname{max\_depth=25, max\_features=sqrt, min\_samples\_leaf=2, n\_estimators=200}
\expandafter\def\csname val@RF-BruteForce-pca-F1\endcsname{99,993\%}
\expandafter\def\csname val@RF-BruteForce-pca-dF1\endcsname{0,01}
\expandafter\def\csname val@RF-BruteForce-pred-us\endcsname{6,0}
\expandafter\def\csname val@RF-BruteForce-speedup\endcsname{1,00}
\expandafter\def\csname val@RF-BruteForce-throughput\endcsname{166\,228}
\expandafter\def\csname val@RF-BruteForce-train-acc\endcsname{99,998\%}
\expandafter\def\csname val@RF-DDoS-AUC\endcsname{1,0000}
\expandafter\def\csname val@RF-DDoS-AUC-5\endcsname{1,00000}
\expandafter\def\csname val@RF-DDoS-Accuracy\endcsname{100,00\%}
\expandafter\def\csname val@RF-DDoS-Accuracy-100\endcsname{100}
\expandafter\def\csname val@RF-DDoS-Accuracy-3\endcsname{100,000\%}
\expandafter\def\csname val@RF-DDoS-Accuracy-hi\endcsname{100,00\%}
\expandafter\def\csname val@RF-DDoS-Accuracy-hw\endcsname{0,00}
\expandafter\def\csname val@RF-DDoS-Accuracy-lo\endcsname{100,00\%}
\expandafter\def\csname val@RF-DDoS-F1\endcsname{100,00\%}
\expandafter\def\csname val@RF-DDoS-F1-3\endcsname{100,000\%}
\expandafter\def\csname val@RF-DDoS-F1-hi\endcsname{100,00\%}
\expandafter\def\csname val@RF-DDoS-F1-hw\endcsname{0,00}
\expandafter\def\csname val@RF-DDoS-F1-lo\endcsname{100,00\%}
\expandafter\def\csname val@RF-DDoS-FN\endcsname{0}
\expandafter\def\csname val@RF-DDoS-FNR\endcsname{0,00\%}
\expandafter\def\csname val@RF-DDoS-FNR-3\endcsname{0,000\%}
\expandafter\def\csname val@RF-DDoS-FP\endcsname{0}
\expandafter\def\csname val@RF-DDoS-FPR\endcsname{0,00\%}
\expandafter\def\csname val@RF-DDoS-FPR-3\endcsname{0,000\%}
\expandafter\def\csname val@RF-DDoS-Precision\endcsname{100,00\%}
\expandafter\def\csname val@RF-DDoS-Precision-100\endcsname{100}
\expandafter\def\csname val@RF-DDoS-Precision-3\endcsname{100,000\%}
\expandafter\def\csname val@RF-DDoS-Precision-hi\endcsname{100,00\%}
\expandafter\def\csname val@RF-DDoS-Precision-hw\endcsname{0,00}
\expandafter\def\csname val@RF-DDoS-Precision-lo\endcsname{100,00\%}
\expandafter\def\csname val@RF-DDoS-ROC_AUC-hi\endcsname{100,00\%}
\expandafter\def\csname val@RF-DDoS-ROC_AUC-hw\endcsname{0,00}
\expandafter\def\csname val@RF-DDoS-ROC_AUC-lo\endcsname{100,00\%}
\expandafter\def\csname val@RF-DDoS-Recall-100\endcsname{100}
\expandafter\def\csname val@RF-DDoS-Sensitivity\endcsname{100,00\%}
\expandafter\def\csname val@RF-DDoS-Sensitivity-3\endcsname{100,000\%}
\expandafter\def\csname val@RF-DDoS-Sensitivity-hi\endcsname{100,00\%}
\expandafter\def\csname val@RF-DDoS-Sensitivity-hw\endcsname{0,00}
\expandafter\def\csname val@RF-DDoS-Sensitivity-lo\endcsname{100,00\%}
\expandafter\def\csname val@RF-DDoS-Specificity\endcsname{100,00\%}
\expandafter\def\csname val@RF-DDoS-Specificity-100\endcsname{100}
\expandafter\def\csname val@RF-DDoS-Specificity-3\endcsname{100,000\%}
\expandafter\def\csname val@RF-DDoS-Specificity-hi\endcsname{100,00\%}
\expandafter\def\csname val@RF-DDoS-Specificity-hw\endcsname{0,00}
\expandafter\def\csname val@RF-DDoS-Specificity-lo\endcsname{100,00\%}
\expandafter\def\csname val@RF-DDoS-TN\endcsname{30\,000}
\expandafter\def\csname val@RF-DDoS-TP\endcsname{30\,000}
\expandafter\def\csname val@RF-DDoS-cv\endcsname{99,99\%}
\expandafter\def\csname val@RF-DDoS-cv-std\endcsname{0,01\%}
\expandafter\def\csname val@RF-DDoS-cv-time\endcsname{241}
\expandafter\def\csname val@RF-DDoS-dF1\endcsname{0,00}
\expandafter\def\csname val@RF-DDoS-fit\endcsname{9,8}
\expandafter\def\csname val@RF-DDoS-full-F1\endcsname{100,000\%}
\expandafter\def\csname val@RF-DDoS-full-fit\endcsname{15,8}
\expandafter\def\csname val@RF-DDoS-full-pred-us\endcsname{5,6}
\expandafter\def\csname val@RF-DDoS-gap\endcsname{0,000}
\expandafter\def\csname val@RF-DDoS-nat-Accuracy\endcsname{100,0\%}
\expandafter\def\csname val@RF-DDoS-nat-F1\endcsname{100,0\%}
\expandafter\def\csname val@RF-DDoS-nat-FN\endcsname{0}
\expandafter\def\csname val@RF-DDoS-nat-FP\endcsname{2}
\expandafter\def\csname val@RF-DDoS-nat-Precision\endcsname{100,0\%}
\expandafter\def\csname val@RF-DDoS-nat-Sensitivity\endcsname{100,0\%}
\expandafter\def\csname val@RF-DDoS-nat-Specificity\endcsname{100,0\%}
\expandafter\def\csname val@RF-DDoS-nat-TN\endcsname{144\,503}
\expandafter\def\csname val@RF-DDoS-nat-TP\endcsname{55\,495}
\expandafter\def\csname val@RF-DDoS-nat-fp-per-tp\endcsname{0,00}
\expandafter\def\csname val@RF-DDoS-nat-n\endcsname{200\,000}
\expandafter\def\csname val@RF-DDoS-noiat-F1\endcsname{100,000\%}
\expandafter\def\csname val@RF-DDoS-noiat-dF1\endcsname{0,00}
\expandafter\def\csname val@RF-DDoS-ntrain\endcsname{140\,000}
\expandafter\def\csname val@RF-DDoS-params\endcsname{max\_depth=25, max\_features=0.5, min\_samples\_leaf=2, n\_estimators=200}
\expandafter\def\csname val@RF-DDoS-pca-F1\endcsname{100,000\%}
\expandafter\def\csname val@RF-DDoS-pca-dF1\endcsname{0,00}
\expandafter\def\csname val@RF-DDoS-pred-us\endcsname{5,2}
\expandafter\def\csname val@RF-DDoS-speedup\endcsname{1,08}
\expandafter\def\csname val@RF-DDoS-throughput\endcsname{190\,821}
\expandafter\def\csname val@RF-DDoS-train-acc\endcsname{100,000\%}
\expandafter\def\csname val@SVM-Bot-AUC\endcsname{0,9998}
\expandafter\def\csname val@SVM-Bot-AUC-5\endcsname{0,99983}
\expandafter\def\csname val@SVM-Bot-Accuracy\endcsname{99,60\%}
\expandafter\def\csname val@SVM-Bot-Accuracy-100\endcsname{100}
\expandafter\def\csname val@SVM-Bot-Accuracy-3\endcsname{99,598\%}
\expandafter\def\csname val@SVM-Bot-Accuracy-hi\endcsname{99,65\%}
\expandafter\def\csname val@SVM-Bot-Accuracy-hw\endcsname{0,05}
\expandafter\def\csname val@SVM-Bot-Accuracy-lo\endcsname{99,55\%}
\expandafter\def\csname val@SVM-Bot-F1\endcsname{99,60\%}
\expandafter\def\csname val@SVM-Bot-F1-3\endcsname{99,597\%}
\expandafter\def\csname val@SVM-Bot-F1-hi\endcsname{99,65\%}
\expandafter\def\csname val@SVM-Bot-F1-hw\endcsname{0,05}
\expandafter\def\csname val@SVM-Bot-F1-lo\endcsname{99,55\%}
\expandafter\def\csname val@SVM-Bot-FN\endcsname{214}
\expandafter\def\csname val@SVM-Bot-FNR\endcsname{0,71\%}
\expandafter\def\csname val@SVM-Bot-FNR-3\endcsname{0,713\%}
\expandafter\def\csname val@SVM-Bot-FP\endcsname{27}
\expandafter\def\csname val@SVM-Bot-FPR\endcsname{0,09\%}
\expandafter\def\csname val@SVM-Bot-FPR-3\endcsname{0,090\%}
\expandafter\def\csname val@SVM-Bot-Precision\endcsname{99,91\%}
\expandafter\def\csname val@SVM-Bot-Precision-100\endcsname{100}
\expandafter\def\csname val@SVM-Bot-Precision-3\endcsname{99,909\%}
\expandafter\def\csname val@SVM-Bot-Precision-hi\endcsname{99,94\%}
\expandafter\def\csname val@SVM-Bot-Precision-hw\endcsname{0,03}
\expandafter\def\csname val@SVM-Bot-Precision-lo\endcsname{99,88\%}
\expandafter\def\csname val@SVM-Bot-ROC_AUC-hi\endcsname{99,99\%}
\expandafter\def\csname val@SVM-Bot-ROC_AUC-hw\endcsname{0,01}
\expandafter\def\csname val@SVM-Bot-ROC_AUC-lo\endcsname{99,97\%}
\expandafter\def\csname val@SVM-Bot-Recall-100\endcsname{99}
\expandafter\def\csname val@SVM-Bot-Sensitivity\endcsname{99,29\%}
\expandafter\def\csname val@SVM-Bot-Sensitivity-3\endcsname{99,287\%}
\expandafter\def\csname val@SVM-Bot-Sensitivity-hi\endcsname{99,38\%}
\expandafter\def\csname val@SVM-Bot-Sensitivity-hw\endcsname{0,09}
\expandafter\def\csname val@SVM-Bot-Sensitivity-lo\endcsname{99,19\%}
\expandafter\def\csname val@SVM-Bot-Specificity\endcsname{99,91\%}
\expandafter\def\csname val@SVM-Bot-Specificity-100\endcsname{100}
\expandafter\def\csname val@SVM-Bot-Specificity-3\endcsname{99,910\%}
\expandafter\def\csname val@SVM-Bot-Specificity-hi\endcsname{99,94\%}
\expandafter\def\csname val@SVM-Bot-Specificity-hw\endcsname{0,03}
\expandafter\def\csname val@SVM-Bot-Specificity-lo\endcsname{99,88\%}
\expandafter\def\csname val@SVM-Bot-TN\endcsname{29\,973}
\expandafter\def\csname val@SVM-Bot-TP\endcsname{29\,786}
\expandafter\def\csname val@SVM-Bot-cv\endcsname{99,47\%}
\expandafter\def\csname val@SVM-Bot-cv-std\endcsname{0,10\%}
\expandafter\def\csname val@SVM-Bot-cv-time\endcsname{15}
\expandafter\def\csname val@SVM-Bot-dF1\endcsname{-0,01}
\expandafter\def\csname val@SVM-Bot-fit\endcsname{2,2}
\expandafter\def\csname val@SVM-Bot-full-F1\endcsname{99,609\%}
\expandafter\def\csname val@SVM-Bot-full-fit\endcsname{2,6}
\expandafter\def\csname val@SVM-Bot-full-pred-us\endcsname{56,1}
\expandafter\def\csname val@SVM-Bot-gap\endcsname{-0,012}
\expandafter\def\csname val@SVM-Bot-nat-Accuracy\endcsname{99,8\%}
\expandafter\def\csname val@SVM-Bot-nat-F1\endcsname{99,4\%}
\expandafter\def\csname val@SVM-Bot-nat-FN\endcsname{201}
\expandafter\def\csname val@SVM-Bot-nat-FP\endcsname{156}
\expandafter\def\csname val@SVM-Bot-nat-Precision\endcsname{99,4\%}
\expandafter\def\csname val@SVM-Bot-nat-Sensitivity\endcsname{99,3\%}
\expandafter\def\csname val@SVM-Bot-nat-Specificity\endcsname{99,9\%}
\expandafter\def\csname val@SVM-Bot-nat-TN\endcsname{172\,286}
\expandafter\def\csname val@SVM-Bot-nat-TP\endcsname{27\,357}
\expandafter\def\csname val@SVM-Bot-nat-fp-per-tp\endcsname{0,01}
\expandafter\def\csname val@SVM-Bot-nat-n\endcsname{200\,000}
\expandafter\def\csname val@SVM-Bot-noiat-F1\endcsname{99,484\%}
\expandafter\def\csname val@SVM-Bot-noiat-dF1\endcsname{-0,11}
\expandafter\def\csname val@SVM-Bot-ntrain\endcsname{50\,000}
\expandafter\def\csname val@SVM-Bot-params\endcsname{C=1, gamma=scale}
\expandafter\def\csname val@SVM-Bot-pca-F1\endcsname{99,600\%}
\expandafter\def\csname val@SVM-Bot-pca-dF1\endcsname{-0,00}
\expandafter\def\csname val@SVM-Bot-pred-us\endcsname{24,1}
\expandafter\def\csname val@SVM-Bot-speedup\endcsname{2,33}
\expandafter\def\csname val@SVM-Bot-throughput\endcsname{41\,468}
\expandafter\def\csname val@SVM-Bot-train-acc\endcsname{99,586\%}
\expandafter\def\csname val@SVM-BruteForce-AUC\endcsname{1,0000}
\expandafter\def\csname val@SVM-BruteForce-AUC-5\endcsname{0,99995}
\expandafter\def\csname val@SVM-BruteForce-Accuracy\endcsname{99,99\%}
\expandafter\def\csname val@SVM-BruteForce-Accuracy-100\endcsname{100}
\expandafter\def\csname val@SVM-BruteForce-Accuracy-3\endcsname{99,986\%}
\expandafter\def\csname val@SVM-BruteForce-Accuracy-hi\endcsname{99,99\%}
\expandafter\def\csname val@SVM-BruteForce-Accuracy-hw\endcsname{0,01}
\expandafter\def\csname val@SVM-BruteForce-Accuracy-lo\endcsname{99,98\%}
\expandafter\def\csname val@SVM-BruteForce-F1\endcsname{99,99\%}
\expandafter\def\csname val@SVM-BruteForce-F1-3\endcsname{99,986\%}
\expandafter\def\csname val@SVM-BruteForce-F1-hi\endcsname{99,99\%}
\expandafter\def\csname val@SVM-BruteForce-F1-hw\endcsname{0,01}
\expandafter\def\csname val@SVM-BruteForce-F1-lo\endcsname{99,98\%}
\expandafter\def\csname val@SVM-BruteForce-FN\endcsname{7}
\expandafter\def\csname val@SVM-BruteForce-FNR\endcsname{0,02\%}
\expandafter\def\csname val@SVM-BruteForce-FNR-3\endcsname{0,025\%}
\expandafter\def\csname val@SVM-BruteForce-FP\endcsname{1}
\expandafter\def\csname val@SVM-BruteForce-FPR\endcsname{0,00\%}
\expandafter\def\csname val@SVM-BruteForce-FPR-3\endcsname{0,004\%}
\expandafter\def\csname val@SVM-BruteForce-Precision\endcsname{100,00\%}
\expandafter\def\csname val@SVM-BruteForce-Precision-100\endcsname{100}
\expandafter\def\csname val@SVM-BruteForce-Precision-3\endcsname{99,996\%}
\expandafter\def\csname val@SVM-BruteForce-Precision-hi\endcsname{100,00\%}
\expandafter\def\csname val@SVM-BruteForce-Precision-hw\endcsname{0,01}
\expandafter\def\csname val@SVM-BruteForce-Precision-lo\endcsname{99,99\%}
\expandafter\def\csname val@SVM-BruteForce-ROC_AUC-hi\endcsname{100,00\%}
\expandafter\def\csname val@SVM-BruteForce-ROC_AUC-hw\endcsname{0,01}
\expandafter\def\csname val@SVM-BruteForce-ROC_AUC-lo\endcsname{99,99\%}
\expandafter\def\csname val@SVM-BruteForce-Recall-100\endcsname{100}
\expandafter\def\csname val@SVM-BruteForce-Sensitivity\endcsname{99,98\%}
\expandafter\def\csname val@SVM-BruteForce-Sensitivity-3\endcsname{99,975\%}
\expandafter\def\csname val@SVM-BruteForce-Sensitivity-hi\endcsname{99,99\%}
\expandafter\def\csname val@SVM-BruteForce-Sensitivity-hw\endcsname{0,02}
\expandafter\def\csname val@SVM-BruteForce-Sensitivity-lo\endcsname{99,95\%}
\expandafter\def\csname val@SVM-BruteForce-Specificity\endcsname{100,00\%}
\expandafter\def\csname val@SVM-BruteForce-Specificity-100\endcsname{100}
\expandafter\def\csname val@SVM-BruteForce-Specificity-3\endcsname{99,996\%}
\expandafter\def\csname val@SVM-BruteForce-Specificity-hi\endcsname{100,00\%}
\expandafter\def\csname val@SVM-BruteForce-Specificity-hw\endcsname{0,01}
\expandafter\def\csname val@SVM-BruteForce-Specificity-lo\endcsname{99,99\%}
\expandafter\def\csname val@SVM-BruteForce-TN\endcsname{28\,228}
\expandafter\def\csname val@SVM-BruteForce-TP\endcsname{28\,221}
\expandafter\def\csname val@SVM-BruteForce-cv\endcsname{99,96\%}
\expandafter\def\csname val@SVM-BruteForce-cv-std\endcsname{0,03\%}
\expandafter\def\csname val@SVM-BruteForce-cv-time\endcsname{10}
\expandafter\def\csname val@SVM-BruteForce-dF1\endcsname{0,01}
\expandafter\def\csname val@SVM-BruteForce-fit\endcsname{1,2}
\expandafter\def\csname val@SVM-BruteForce-full-F1\endcsname{99,977\%}
\expandafter\def\csname val@SVM-BruteForce-full-fit\endcsname{0,4}
\expandafter\def\csname val@SVM-BruteForce-full-pred-us\endcsname{4,1}
\expandafter\def\csname val@SVM-BruteForce-gap\endcsname{-0,014}
\expandafter\def\csname val@SVM-BruteForce-nat-Accuracy\endcsname{100,0\%}
\expandafter\def\csname val@SVM-BruteForce-nat-F1\endcsname{100,0\%}
\expandafter\def\csname val@SVM-BruteForce-nat-FN\endcsname{7}
\expandafter\def\csname val@SVM-BruteForce-nat-FP\endcsname{12}
\expandafter\def\csname val@SVM-BruteForce-nat-Precision\endcsname{100,0\%}
\expandafter\def\csname val@SVM-BruteForce-nat-Sensitivity\endcsname{100,0\%}
\expandafter\def\csname val@SVM-BruteForce-nat-Specificity\endcsname{100,0\%}
\expandafter\def\csname val@SVM-BruteForce-nat-TN\endcsname{132\,685}
\expandafter\def\csname val@SVM-BruteForce-nat-TP\endcsname{28\,221}
\expandafter\def\csname val@SVM-BruteForce-nat-fp-per-tp\endcsname{0,00}
\expandafter\def\csname val@SVM-BruteForce-nat-n\endcsname{160\,925}
\expandafter\def\csname val@SVM-BruteForce-noiat-F1\endcsname{99,977\%}
\expandafter\def\csname val@SVM-BruteForce-noiat-dF1\endcsname{-0,01}
\expandafter\def\csname val@SVM-BruteForce-ntrain\endcsname{50\,000}
\expandafter\def\csname val@SVM-BruteForce-params\endcsname{C=10, gamma=scale}
\expandafter\def\csname val@SVM-BruteForce-pca-F1\endcsname{99,972\%}
\expandafter\def\csname val@SVM-BruteForce-pca-dF1\endcsname{0,01}
\expandafter\def\csname val@SVM-BruteForce-pred-us\endcsname{2,4}
\expandafter\def\csname val@SVM-BruteForce-speedup\endcsname{1,70}
\expandafter\def\csname val@SVM-BruteForce-throughput\endcsname{419\,217}
\expandafter\def\csname val@SVM-BruteForce-train-acc\endcsname{99,972\%}
\expandafter\def\csname val@SVM-DDoS-AUC\endcsname{1,0000}
\expandafter\def\csname val@SVM-DDoS-AUC-5\endcsname{0,99996}
\expandafter\def\csname val@SVM-DDoS-Accuracy\endcsname{99,95\%}
\expandafter\def\csname val@SVM-DDoS-Accuracy-100\endcsname{100}
\expandafter\def\csname val@SVM-DDoS-Accuracy-3\endcsname{99,950\%}
\expandafter\def\csname val@SVM-DDoS-Accuracy-hi\endcsname{99,97\%}
\expandafter\def\csname val@SVM-DDoS-Accuracy-hw\endcsname{0,02}
\expandafter\def\csname val@SVM-DDoS-Accuracy-lo\endcsname{99,93\%}
\expandafter\def\csname val@SVM-DDoS-F1\endcsname{99,95\%}
\expandafter\def\csname val@SVM-DDoS-F1-3\endcsname{99,950\%}
\expandafter\def\csname val@SVM-DDoS-F1-hi\endcsname{99,97\%}
\expandafter\def\csname val@SVM-DDoS-F1-hw\endcsname{0,02}
\expandafter\def\csname val@SVM-DDoS-F1-lo\endcsname{99,93\%}
\expandafter\def\csname val@SVM-DDoS-FN\endcsname{12}
\expandafter\def\csname val@SVM-DDoS-FNR\endcsname{0,04\%}
\expandafter\def\csname val@SVM-DDoS-FNR-3\endcsname{0,040\%}
\expandafter\def\csname val@SVM-DDoS-FP\endcsname{18}
\expandafter\def\csname val@SVM-DDoS-FPR\endcsname{0,06\%}
\expandafter\def\csname val@SVM-DDoS-FPR-3\endcsname{0,060\%}
\expandafter\def\csname val@SVM-DDoS-Precision\endcsname{99,94\%}
\expandafter\def\csname val@SVM-DDoS-Precision-100\endcsname{100}
\expandafter\def\csname val@SVM-DDoS-Precision-3\endcsname{99,940\%}
\expandafter\def\csname val@SVM-DDoS-Precision-hi\endcsname{99,96\%}
\expandafter\def\csname val@SVM-DDoS-Precision-hw\endcsname{0,03}
\expandafter\def\csname val@SVM-DDoS-Precision-lo\endcsname{99,91\%}
\expandafter\def\csname val@SVM-DDoS-ROC_AUC-hi\endcsname{100,00\%}
\expandafter\def\csname val@SVM-DDoS-ROC_AUC-hw\endcsname{0,01}
\expandafter\def\csname val@SVM-DDoS-ROC_AUC-lo\endcsname{99,99\%}
\expandafter\def\csname val@SVM-DDoS-Recall-100\endcsname{100}
\expandafter\def\csname val@SVM-DDoS-Sensitivity\endcsname{99,96\%}
\expandafter\def\csname val@SVM-DDoS-Sensitivity-3\endcsname{99,960\%}
\expandafter\def\csname val@SVM-DDoS-Sensitivity-hi\endcsname{99,98\%}
\expandafter\def\csname val@SVM-DDoS-Sensitivity-hw\endcsname{0,02}
\expandafter\def\csname val@SVM-DDoS-Sensitivity-lo\endcsname{99,94\%}
\expandafter\def\csname val@SVM-DDoS-Specificity\endcsname{99,94\%}
\expandafter\def\csname val@SVM-DDoS-Specificity-100\endcsname{100}
\expandafter\def\csname val@SVM-DDoS-Specificity-3\endcsname{99,940\%}
\expandafter\def\csname val@SVM-DDoS-Specificity-hi\endcsname{99,96\%}
\expandafter\def\csname val@SVM-DDoS-Specificity-hw\endcsname{0,03}
\expandafter\def\csname val@SVM-DDoS-Specificity-lo\endcsname{99,91\%}
\expandafter\def\csname val@SVM-DDoS-TN\endcsname{29\,982}
\expandafter\def\csname val@SVM-DDoS-TP\endcsname{29\,988}
\expandafter\def\csname val@SVM-DDoS-cv\endcsname{99,96\%}
\expandafter\def\csname val@SVM-DDoS-cv-std\endcsname{0,03\%}
\expandafter\def\csname val@SVM-DDoS-cv-time\endcsname{12}
\expandafter\def\csname val@SVM-DDoS-dF1\endcsname{-0,02}
\expandafter\def\csname val@SVM-DDoS-fit\endcsname{1,8}
\expandafter\def\csname val@SVM-DDoS-full-F1\endcsname{99,972\%}
\expandafter\def\csname val@SVM-DDoS-full-fit\endcsname{3,1}
\expandafter\def\csname val@SVM-DDoS-full-pred-us\endcsname{60,5}
\expandafter\def\csname val@SVM-DDoS-gap\endcsname{0,004}
\expandafter\def\csname val@SVM-DDoS-nat-Accuracy\endcsname{99,9\%}
\expandafter\def\csname val@SVM-DDoS-nat-F1\endcsname{99,9\%}
\expandafter\def\csname val@SVM-DDoS-nat-FN\endcsname{17}
\expandafter\def\csname val@SVM-DDoS-nat-FP\endcsname{116}
\expandafter\def\csname val@SVM-DDoS-nat-Precision\endcsname{99,8\%}
\expandafter\def\csname val@SVM-DDoS-nat-Sensitivity\endcsname{100,0\%}
\expandafter\def\csname val@SVM-DDoS-nat-Specificity\endcsname{99,9\%}
\expandafter\def\csname val@SVM-DDoS-nat-TN\endcsname{144\,389}
\expandafter\def\csname val@SVM-DDoS-nat-TP\endcsname{55\,478}
\expandafter\def\csname val@SVM-DDoS-nat-fp-per-tp\endcsname{0,00}
\expandafter\def\csname val@SVM-DDoS-nat-n\endcsname{200\,000}
\expandafter\def\csname val@SVM-DDoS-noiat-F1\endcsname{99,933\%}
\expandafter\def\csname val@SVM-DDoS-noiat-dF1\endcsname{-0,02}
\expandafter\def\csname val@SVM-DDoS-ntrain\endcsname{50\,000}
\expandafter\def\csname val@SVM-DDoS-params\endcsname{C=10, gamma=0.1}
\expandafter\def\csname val@SVM-DDoS-pca-F1\endcsname{99,955\%}
\expandafter\def\csname val@SVM-DDoS-pca-dF1\endcsname{-0,01}
\expandafter\def\csname val@SVM-DDoS-pred-us\endcsname{10,9}
\expandafter\def\csname val@SVM-DDoS-speedup\endcsname{5,58}
\expandafter\def\csname val@SVM-DDoS-throughput\endcsname{92\,118}
\expandafter\def\csname val@SVM-DDoS-train-acc\endcsname{99,954\%}
\expandafter\def\csname val@XGB-Bot-AUC\endcsname{1,0000}
\expandafter\def\csname val@XGB-Bot-AUC-5\endcsname{1,00000}
\expandafter\def\csname val@XGB-Bot-Accuracy\endcsname{99,98\%}
\expandafter\def\csname val@XGB-Bot-Accuracy-100\endcsname{100}
\expandafter\def\csname val@XGB-Bot-Accuracy-3\endcsname{99,982\%}
\expandafter\def\csname val@XGB-Bot-Accuracy-hi\endcsname{99,99\%}
\expandafter\def\csname val@XGB-Bot-Accuracy-hw\endcsname{0,01}
\expandafter\def\csname val@XGB-Bot-Accuracy-lo\endcsname{99,97\%}
\expandafter\def\csname val@XGB-Bot-F1\endcsname{99,98\%}
\expandafter\def\csname val@XGB-Bot-F1-3\endcsname{99,982\%}
\expandafter\def\csname val@XGB-Bot-F1-hi\endcsname{99,99\%}
\expandafter\def\csname val@XGB-Bot-F1-hw\endcsname{0,01}
\expandafter\def\csname val@XGB-Bot-F1-lo\endcsname{99,97\%}
\expandafter\def\csname val@XGB-Bot-FN\endcsname{6}
\expandafter\def\csname val@XGB-Bot-FNR\endcsname{0,02\%}
\expandafter\def\csname val@XGB-Bot-FNR-3\endcsname{0,020\%}
\expandafter\def\csname val@XGB-Bot-FP\endcsname{5}
\expandafter\def\csname val@XGB-Bot-FPR\endcsname{0,02\%}
\expandafter\def\csname val@XGB-Bot-FPR-3\endcsname{0,017\%}
\expandafter\def\csname val@XGB-Bot-Precision\endcsname{99,98\%}
\expandafter\def\csname val@XGB-Bot-Precision-100\endcsname{100}
\expandafter\def\csname val@XGB-Bot-Precision-3\endcsname{99,983\%}
\expandafter\def\csname val@XGB-Bot-Precision-hi\endcsname{100,00\%}
\expandafter\def\csname val@XGB-Bot-Precision-hw\endcsname{0,01}
\expandafter\def\csname val@XGB-Bot-Precision-lo\endcsname{99,97\%}
\expandafter\def\csname val@XGB-Bot-ROC_AUC-hi\endcsname{100,00\%}
\expandafter\def\csname val@XGB-Bot-ROC_AUC-hw\endcsname{0,00}
\expandafter\def\csname val@XGB-Bot-ROC_AUC-lo\endcsname{100,00\%}
\expandafter\def\csname val@XGB-Bot-Recall-100\endcsname{100}
\expandafter\def\csname val@XGB-Bot-Sensitivity\endcsname{99,98\%}
\expandafter\def\csname val@XGB-Bot-Sensitivity-3\endcsname{99,980\%}
\expandafter\def\csname val@XGB-Bot-Sensitivity-hi\endcsname{99,99\%}
\expandafter\def\csname val@XGB-Bot-Sensitivity-hw\endcsname{0,02}
\expandafter\def\csname val@XGB-Bot-Sensitivity-lo\endcsname{99,96\%}
\expandafter\def\csname val@XGB-Bot-Specificity\endcsname{99,98\%}
\expandafter\def\csname val@XGB-Bot-Specificity-100\endcsname{100}
\expandafter\def\csname val@XGB-Bot-Specificity-3\endcsname{99,983\%}
\expandafter\def\csname val@XGB-Bot-Specificity-hi\endcsname{100,00\%}
\expandafter\def\csname val@XGB-Bot-Specificity-hw\endcsname{0,02}
\expandafter\def\csname val@XGB-Bot-Specificity-lo\endcsname{99,97\%}
\expandafter\def\csname val@XGB-Bot-TN\endcsname{29\,995}
\expandafter\def\csname val@XGB-Bot-TP\endcsname{29\,994}
\expandafter\def\csname val@XGB-Bot-cv\endcsname{99,96\%}
\expandafter\def\csname val@XGB-Bot-cv-std\endcsname{0,02\%}
\expandafter\def\csname val@XGB-Bot-cv-time\endcsname{70}
\expandafter\def\csname val@XGB-Bot-dF1\endcsname{0,00}
\expandafter\def\csname val@XGB-Bot-fit\endcsname{5,1}
\expandafter\def\csname val@XGB-Bot-full-F1\endcsname{99,978\%}
\expandafter\def\csname val@XGB-Bot-full-fit\endcsname{4,1}
\expandafter\def\csname val@XGB-Bot-full-pred-us\endcsname{1,5}
\expandafter\def\csname val@XGB-Bot-gap\endcsname{0,010}
\expandafter\def\csname val@XGB-Bot-nat-Accuracy\endcsname{100,0\%}
\expandafter\def\csname val@XGB-Bot-nat-F1\endcsname{99,9\%}
\expandafter\def\csname val@XGB-Bot-nat-FN\endcsname{7}
\expandafter\def\csname val@XGB-Bot-nat-FP\endcsname{36}
\expandafter\def\csname val@XGB-Bot-nat-Precision\endcsname{99,9\%}
\expandafter\def\csname val@XGB-Bot-nat-Sensitivity\endcsname{100,0\%}
\expandafter\def\csname val@XGB-Bot-nat-Specificity\endcsname{100,0\%}
\expandafter\def\csname val@XGB-Bot-nat-TN\endcsname{172\,406}
\expandafter\def\csname val@XGB-Bot-nat-TP\endcsname{27\,551}
\expandafter\def\csname val@XGB-Bot-nat-fp-per-tp\endcsname{0,00}
\expandafter\def\csname val@XGB-Bot-nat-n\endcsname{200\,000}
\expandafter\def\csname val@XGB-Bot-noiat-F1\endcsname{99,965\%}
\expandafter\def\csname val@XGB-Bot-noiat-dF1\endcsname{-0,02}
\expandafter\def\csname val@XGB-Bot-ntrain\endcsname{140\,000}
\expandafter\def\csname val@XGB-Bot-params\endcsname{learning\_rate=0.1, max\_depth=6, min\_child\_weight=5, n\_estimators=800}
\expandafter\def\csname val@XGB-Bot-pca-F1\endcsname{99,942\%}
\expandafter\def\csname val@XGB-Bot-pca-dF1\endcsname{0,04}
\expandafter\def\csname val@XGB-Bot-pred-us\endcsname{1,4}
\expandafter\def\csname val@XGB-Bot-speedup\endcsname{1,05}
\expandafter\def\csname val@XGB-Bot-throughput\endcsname{704\,173}
\expandafter\def\csname val@XGB-Bot-train-acc\endcsname{99,991\%}
\expandafter\def\csname val@XGB-BruteForce-AUC\endcsname{1,0000}
\expandafter\def\csname val@XGB-BruteForce-AUC-5\endcsname{1,00000}
\expandafter\def\csname val@XGB-BruteForce-Accuracy\endcsname{100,00\%}
\expandafter\def\csname val@XGB-BruteForce-Accuracy-100\endcsname{100}
\expandafter\def\csname val@XGB-BruteForce-Accuracy-3\endcsname{99,998\%}
\expandafter\def\csname val@XGB-BruteForce-Accuracy-hi\endcsname{100,00\%}
\expandafter\def\csname val@XGB-BruteForce-Accuracy-hw\endcsname{0,00}
\expandafter\def\csname val@XGB-BruteForce-Accuracy-lo\endcsname{99,99\%}
\expandafter\def\csname val@XGB-BruteForce-F1\endcsname{100,00\%}
\expandafter\def\csname val@XGB-BruteForce-F1-3\endcsname{99,998\%}
\expandafter\def\csname val@XGB-BruteForce-F1-hi\endcsname{100,00\%}
\expandafter\def\csname val@XGB-BruteForce-F1-hw\endcsname{0,00}
\expandafter\def\csname val@XGB-BruteForce-F1-lo\endcsname{99,99\%}
\expandafter\def\csname val@XGB-BruteForce-FN\endcsname{1}
\expandafter\def\csname val@XGB-BruteForce-FNR\endcsname{0,00\%}
\expandafter\def\csname val@XGB-BruteForce-FNR-3\endcsname{0,004\%}
\expandafter\def\csname val@XGB-BruteForce-FP\endcsname{0}
\expandafter\def\csname val@XGB-BruteForce-FPR\endcsname{0,00\%}
\expandafter\def\csname val@XGB-BruteForce-FPR-3\endcsname{0,000\%}
\expandafter\def\csname val@XGB-BruteForce-Precision\endcsname{100,00\%}
\expandafter\def\csname val@XGB-BruteForce-Precision-100\endcsname{100}
\expandafter\def\csname val@XGB-BruteForce-Precision-3\endcsname{100,000\%}
\expandafter\def\csname val@XGB-BruteForce-Precision-hi\endcsname{100,00\%}
\expandafter\def\csname val@XGB-BruteForce-Precision-hw\endcsname{0,00}
\expandafter\def\csname val@XGB-BruteForce-Precision-lo\endcsname{100,00\%}
\expandafter\def\csname val@XGB-BruteForce-ROC_AUC-hi\endcsname{100,00\%}
\expandafter\def\csname val@XGB-BruteForce-ROC_AUC-hw\endcsname{0,00}
\expandafter\def\csname val@XGB-BruteForce-ROC_AUC-lo\endcsname{100,00\%}
\expandafter\def\csname val@XGB-BruteForce-Recall-100\endcsname{100}
\expandafter\def\csname val@XGB-BruteForce-Sensitivity\endcsname{100,00\%}
\expandafter\def\csname val@XGB-BruteForce-Sensitivity-3\endcsname{99,996\%}
\expandafter\def\csname val@XGB-BruteForce-Sensitivity-hi\endcsname{100,00\%}
\expandafter\def\csname val@XGB-BruteForce-Sensitivity-hw\endcsname{0,01}
\expandafter\def\csname val@XGB-BruteForce-Sensitivity-lo\endcsname{99,99\%}
\expandafter\def\csname val@XGB-BruteForce-Specificity\endcsname{100,00\%}
\expandafter\def\csname val@XGB-BruteForce-Specificity-100\endcsname{100}
\expandafter\def\csname val@XGB-BruteForce-Specificity-3\endcsname{100,000\%}
\expandafter\def\csname val@XGB-BruteForce-Specificity-hi\endcsname{100,00\%}
\expandafter\def\csname val@XGB-BruteForce-Specificity-hw\endcsname{0,00}
\expandafter\def\csname val@XGB-BruteForce-Specificity-lo\endcsname{100,00\%}
\expandafter\def\csname val@XGB-BruteForce-TN\endcsname{28\,229}
\expandafter\def\csname val@XGB-BruteForce-TP\endcsname{28\,227}
\expandafter\def\csname val@XGB-BruteForce-cv\endcsname{99,99\%}
\expandafter\def\csname val@XGB-BruteForce-cv-std\endcsname{0,01\%}
\expandafter\def\csname val@XGB-BruteForce-cv-time\endcsname{53}
\expandafter\def\csname val@XGB-BruteForce-dF1\endcsname{-0,00}
\expandafter\def\csname val@XGB-BruteForce-fit\endcsname{3,8}
\expandafter\def\csname val@XGB-BruteForce-full-F1\endcsname{100,000\%}
\expandafter\def\csname val@XGB-BruteForce-full-fit\endcsname{2,4}
\expandafter\def\csname val@XGB-BruteForce-full-pred-us\endcsname{0,8}
\expandafter\def\csname val@XGB-BruteForce-gap\endcsname{-0,001}
\expandafter\def\csname val@XGB-BruteForce-nat-Accuracy\endcsname{100,0\%}
\expandafter\def\csname val@XGB-BruteForce-nat-F1\endcsname{100,0\%}
\expandafter\def\csname val@XGB-BruteForce-nat-FN\endcsname{1}
\expandafter\def\csname val@XGB-BruteForce-nat-FP\endcsname{5}
\expandafter\def\csname val@XGB-BruteForce-nat-Precision\endcsname{100,0\%}
\expandafter\def\csname val@XGB-BruteForce-nat-Sensitivity\endcsname{100,0\%}
\expandafter\def\csname val@XGB-BruteForce-nat-Specificity\endcsname{100,0\%}
\expandafter\def\csname val@XGB-BruteForce-nat-TN\endcsname{132\,692}
\expandafter\def\csname val@XGB-BruteForce-nat-TP\endcsname{28\,227}
\expandafter\def\csname val@XGB-BruteForce-nat-fp-per-tp\endcsname{0,00}
\expandafter\def\csname val@XGB-BruteForce-nat-n\endcsname{160\,925}
\expandafter\def\csname val@XGB-BruteForce-noiat-F1\endcsname{99,998\%}
\expandafter\def\csname val@XGB-BruteForce-noiat-dF1\endcsname{0,00}
\expandafter\def\csname val@XGB-BruteForce-ntrain\endcsname{131\,731}
\expandafter\def\csname val@XGB-BruteForce-params\endcsname{learning\_rate=0.03, max\_depth=6, min\_child\_weight=5, n\_estimators=300}
\expandafter\def\csname val@XGB-BruteForce-pca-F1\endcsname{99,991\%}
\expandafter\def\csname val@XGB-BruteForce-pca-dF1\endcsname{0,01}
\expandafter\def\csname val@XGB-BruteForce-pred-us\endcsname{0,9}
\expandafter\def\csname val@XGB-BruteForce-speedup\endcsname{0,86}
\expandafter\def\csname val@XGB-BruteForce-throughput\endcsname{1\,098\,465}
\expandafter\def\csname val@XGB-BruteForce-train-acc\endcsname{99,997\%}
\expandafter\def\csname val@XGB-DDoS-AUC\endcsname{1,0000}
\expandafter\def\csname val@XGB-DDoS-AUC-5\endcsname{1,00000}
\expandafter\def\csname val@XGB-DDoS-Accuracy\endcsname{100,00\%}
\expandafter\def\csname val@XGB-DDoS-Accuracy-100\endcsname{100}
\expandafter\def\csname val@XGB-DDoS-Accuracy-3\endcsname{100,000\%}
\expandafter\def\csname val@XGB-DDoS-Accuracy-hi\endcsname{100,00\%}
\expandafter\def\csname val@XGB-DDoS-Accuracy-hw\endcsname{0,00}
\expandafter\def\csname val@XGB-DDoS-Accuracy-lo\endcsname{100,00\%}
\expandafter\def\csname val@XGB-DDoS-F1\endcsname{100,00\%}
\expandafter\def\csname val@XGB-DDoS-F1-3\endcsname{100,000\%}
\expandafter\def\csname val@XGB-DDoS-F1-hi\endcsname{100,00\%}
\expandafter\def\csname val@XGB-DDoS-F1-hw\endcsname{0,00}
\expandafter\def\csname val@XGB-DDoS-F1-lo\endcsname{100,00\%}
\expandafter\def\csname val@XGB-DDoS-FN\endcsname{0}
\expandafter\def\csname val@XGB-DDoS-FNR\endcsname{0,00\%}
\expandafter\def\csname val@XGB-DDoS-FNR-3\endcsname{0,000\%}
\expandafter\def\csname val@XGB-DDoS-FP\endcsname{0}
\expandafter\def\csname val@XGB-DDoS-FPR\endcsname{0,00\%}
\expandafter\def\csname val@XGB-DDoS-FPR-3\endcsname{0,000\%}
\expandafter\def\csname val@XGB-DDoS-Precision\endcsname{100,00\%}
\expandafter\def\csname val@XGB-DDoS-Precision-100\endcsname{100}
\expandafter\def\csname val@XGB-DDoS-Precision-3\endcsname{100,000\%}
\expandafter\def\csname val@XGB-DDoS-Precision-hi\endcsname{100,00\%}
\expandafter\def\csname val@XGB-DDoS-Precision-hw\endcsname{0,00}
\expandafter\def\csname val@XGB-DDoS-Precision-lo\endcsname{100,00\%}
\expandafter\def\csname val@XGB-DDoS-ROC_AUC-hi\endcsname{100,00\%}
\expandafter\def\csname val@XGB-DDoS-ROC_AUC-hw\endcsname{0,00}
\expandafter\def\csname val@XGB-DDoS-ROC_AUC-lo\endcsname{100,00\%}
\expandafter\def\csname val@XGB-DDoS-Recall-100\endcsname{100}
\expandafter\def\csname val@XGB-DDoS-Sensitivity\endcsname{100,00\%}
\expandafter\def\csname val@XGB-DDoS-Sensitivity-3\endcsname{100,000\%}
\expandafter\def\csname val@XGB-DDoS-Sensitivity-hi\endcsname{100,00\%}
\expandafter\def\csname val@XGB-DDoS-Sensitivity-hw\endcsname{0,00}
\expandafter\def\csname val@XGB-DDoS-Sensitivity-lo\endcsname{100,00\%}
\expandafter\def\csname val@XGB-DDoS-Specificity\endcsname{100,00\%}
\expandafter\def\csname val@XGB-DDoS-Specificity-100\endcsname{100}
\expandafter\def\csname val@XGB-DDoS-Specificity-3\endcsname{100,000\%}
\expandafter\def\csname val@XGB-DDoS-Specificity-hi\endcsname{100,00\%}
\expandafter\def\csname val@XGB-DDoS-Specificity-hw\endcsname{0,00}
\expandafter\def\csname val@XGB-DDoS-Specificity-lo\endcsname{100,00\%}
\expandafter\def\csname val@XGB-DDoS-TN\endcsname{30\,000}
\expandafter\def\csname val@XGB-DDoS-TP\endcsname{30\,000}
\expandafter\def\csname val@XGB-DDoS-cv\endcsname{100,00\%}
\expandafter\def\csname val@XGB-DDoS-cv-std\endcsname{0,00\%}
\expandafter\def\csname val@XGB-DDoS-cv-time\endcsname{61}
\expandafter\def\csname val@XGB-DDoS-dF1\endcsname{0,00}
\expandafter\def\csname val@XGB-DDoS-fit\endcsname{4,2}
\expandafter\def\csname val@XGB-DDoS-full-F1\endcsname{100,000\%}
\expandafter\def\csname val@XGB-DDoS-full-fit\endcsname{1,9}
\expandafter\def\csname val@XGB-DDoS-full-pred-us\endcsname{1,3}
\expandafter\def\csname val@XGB-DDoS-gap\endcsname{0,000}
\expandafter\def\csname val@XGB-DDoS-nat-Accuracy\endcsname{100,0\%}
\expandafter\def\csname val@XGB-DDoS-nat-F1\endcsname{100,0\%}
\expandafter\def\csname val@XGB-DDoS-nat-FN\endcsname{0}
\expandafter\def\csname val@XGB-DDoS-nat-FP\endcsname{0}
\expandafter\def\csname val@XGB-DDoS-nat-Precision\endcsname{100,0\%}
\expandafter\def\csname val@XGB-DDoS-nat-Sensitivity\endcsname{100,0\%}
\expandafter\def\csname val@XGB-DDoS-nat-Specificity\endcsname{100,0\%}
\expandafter\def\csname val@XGB-DDoS-nat-TN\endcsname{144\,505}
\expandafter\def\csname val@XGB-DDoS-nat-TP\endcsname{55\,495}
\expandafter\def\csname val@XGB-DDoS-nat-fp-per-tp\endcsname{0,00}
\expandafter\def\csname val@XGB-DDoS-nat-n\endcsname{200\,000}
\expandafter\def\csname val@XGB-DDoS-noiat-F1\endcsname{100,000\%}
\expandafter\def\csname val@XGB-DDoS-noiat-dF1\endcsname{0,00}
\expandafter\def\csname val@XGB-DDoS-ntrain\endcsname{140\,000}
\expandafter\def\csname val@XGB-DDoS-params\endcsname{learning\_rate=0.03, max\_depth=6, min\_child\_weight=1, n\_estimators=300}
\expandafter\def\csname val@XGB-DDoS-pca-F1\endcsname{99,998\%}
\expandafter\def\csname val@XGB-DDoS-pca-dF1\endcsname{0,00}
\expandafter\def\csname val@XGB-DDoS-pred-us\endcsname{0,9}
\expandafter\def\csname val@XGB-DDoS-speedup\endcsname{1,37}
\expandafter\def\csname val@XGB-DDoS-throughput\endcsname{1\,058\,630}
\expandafter\def\csname val@XGB-DDoS-train-acc\endcsname{100,000\%}
\expandafter\def\csname val@ap-RF-Bot\endcsname{1,000}
\expandafter\def\csname val@ap-RF-BruteForce\endcsname{1,000}
\expandafter\def\csname val@ap-RF-DDoS\endcsname{1,000}
\expandafter\def\csname val@ap-SVM-Bot\endcsname{0,997}
\expandafter\def\csname val@ap-SVM-BruteForce\endcsname{0,999}
\expandafter\def\csname val@ap-SVM-DDoS\endcsname{1,000}
\expandafter\def\csname val@ap-XGB-Bot\endcsname{1,000}
\expandafter\def\csname val@ap-XGB-BruteForce\endcsname{1,000}
\expandafter\def\csname val@ap-XGB-DDoS\endcsname{1,000}
\expandafter\def\csname val@attack-avail-Bot\endcsname{144\,287}
\expandafter\def\csname val@attack-avail-BruteForce\endcsname{94\,094}
\expandafter\def\csname val@attack-avail-DDoS\endcsname{200\,590}
\expandafter\def\csname val@attack-orig-Bot\endcsname{286\,191}
\expandafter\def\csname val@attack-orig-BruteForce\endcsname{380\,949}
\expandafter\def\csname val@attack-orig-DDoS\endcsname{687\,742}
\expandafter\def\csname val@benign-2102-rows\endcsname{360\,774}
\expandafter\def\csname val@benign-2102-tcp\endcsname{99,5\%}
\expandafter\def\csname val@benign-2102-win65535\endcsname{99,4\%}
\expandafter\def\csname val@benign-avail\endcsname{1\,472\,333}
\expandafter\def\csname val@benign-orig\endcsname{1\,790\,843}
\expandafter\def\csname val@benign-other-tcp\endcsname{73,8\%}
\expandafter\def\csname val@ca-RF-Bot-on-Bot\endcsname{100,0\%}
\expandafter\def\csname val@ca-RF-Bot-on-Bot-fa\endcsname{0,03\%}
\expandafter\def\csname val@ca-RF-Bot-on-BruteForce\endcsname{0,0\%}
\expandafter\def\csname val@ca-RF-Bot-on-BruteForce-fa\endcsname{0,08\%}
\expandafter\def\csname val@ca-RF-Bot-on-DDoS\endcsname{0,0\%}
\expandafter\def\csname val@ca-RF-Bot-on-DDoS-fa\endcsname{0,05\%}
\expandafter\def\csname val@ca-RF-BruteForce-on-Bot\endcsname{0,0\%}
\expandafter\def\csname val@ca-RF-BruteForce-on-Bot-fa\endcsname{0,00\%}
\expandafter\def\csname val@ca-RF-BruteForce-on-BruteForce\endcsname{100,0\%}
\expandafter\def\csname val@ca-RF-BruteForce-on-BruteForce-fa\endcsname{0,00\%}
\expandafter\def\csname val@ca-RF-BruteForce-on-DDoS\endcsname{0,0\%}
\expandafter\def\csname val@ca-RF-BruteForce-on-DDoS-fa\endcsname{0,01\%}
\expandafter\def\csname val@ca-RF-DDoS-on-Bot\endcsname{0,0\%}
\expandafter\def\csname val@ca-RF-DDoS-on-Bot-fa\endcsname{0,00\%}
\expandafter\def\csname val@ca-RF-DDoS-on-BruteForce\endcsname{0,0\%}
\expandafter\def\csname val@ca-RF-DDoS-on-BruteForce-fa\endcsname{0,00\%}
\expandafter\def\csname val@ca-RF-DDoS-on-DDoS\endcsname{100,0\%}
\expandafter\def\csname val@ca-RF-DDoS-on-DDoS-fa\endcsname{0,00\%}
\expandafter\def\csname val@ca-SVM-Bot-on-Bot\endcsname{99,4\%}
\expandafter\def\csname val@ca-SVM-Bot-on-Bot-fa\endcsname{0,13\%}
\expandafter\def\csname val@ca-SVM-Bot-on-BruteForce\endcsname{0,0\%}
\expandafter\def\csname val@ca-SVM-Bot-on-BruteForce-fa\endcsname{0,05\%}
\expandafter\def\csname val@ca-SVM-Bot-on-DDoS\endcsname{0,0\%}
\expandafter\def\csname val@ca-SVM-Bot-on-DDoS-fa\endcsname{0,08\%}
\expandafter\def\csname val@ca-SVM-BruteForce-on-Bot\endcsname{0,0\%}
\expandafter\def\csname val@ca-SVM-BruteForce-on-Bot-fa\endcsname{0,00\%}
\expandafter\def\csname val@ca-SVM-BruteForce-on-BruteForce\endcsname{100,0\%}
\expandafter\def\csname val@ca-SVM-BruteForce-on-BruteForce-fa\endcsname{0,01\%}
\expandafter\def\csname val@ca-SVM-BruteForce-on-DDoS\endcsname{0,0\%}
\expandafter\def\csname val@ca-SVM-BruteForce-on-DDoS-fa\endcsname{0,00\%}
\expandafter\def\csname val@ca-SVM-DDoS-on-Bot\endcsname{0,0\%}
\expandafter\def\csname val@ca-SVM-DDoS-on-Bot-fa\endcsname{0,05\%}
\expandafter\def\csname val@ca-SVM-DDoS-on-BruteForce\endcsname{0,0\%}
\expandafter\def\csname val@ca-SVM-DDoS-on-BruteForce-fa\endcsname{0,05\%}
\expandafter\def\csname val@ca-SVM-DDoS-on-DDoS\endcsname{100,0\%}
\expandafter\def\csname val@ca-SVM-DDoS-on-DDoS-fa\endcsname{0,06\%}
\expandafter\def\csname val@ca-XGB-Bot-on-Bot\endcsname{100,0\%}
\expandafter\def\csname val@ca-XGB-Bot-on-Bot-fa\endcsname{0,02\%}
\expandafter\def\csname val@ca-XGB-Bot-on-BruteForce\endcsname{0,0\%}
\expandafter\def\csname val@ca-XGB-Bot-on-BruteForce-fa\endcsname{0,05\%}
\expandafter\def\csname val@ca-XGB-Bot-on-DDoS\endcsname{0,0\%}
\expandafter\def\csname val@ca-XGB-Bot-on-DDoS-fa\endcsname{0,05\%}
\expandafter\def\csname val@ca-XGB-BruteForce-on-Bot\endcsname{0,0\%}
\expandafter\def\csname val@ca-XGB-BruteForce-on-Bot-fa\endcsname{0,00\%}
\expandafter\def\csname val@ca-XGB-BruteForce-on-BruteForce\endcsname{100,0\%}
\expandafter\def\csname val@ca-XGB-BruteForce-on-BruteForce-fa\endcsname{0,00\%}
\expandafter\def\csname val@ca-XGB-BruteForce-on-DDoS\endcsname{0,0\%}
\expandafter\def\csname val@ca-XGB-BruteForce-on-DDoS-fa\endcsname{0,00\%}
\expandafter\def\csname val@ca-XGB-DDoS-on-Bot\endcsname{0,0\%}
\expandafter\def\csname val@ca-XGB-DDoS-on-Bot-fa\endcsname{0,00\%}
\expandafter\def\csname val@ca-XGB-DDoS-on-BruteForce\endcsname{0,0\%}
\expandafter\def\csname val@ca-XGB-DDoS-on-BruteForce-fa\endcsname{0,00\%}
\expandafter\def\csname val@ca-XGB-DDoS-on-DDoS\endcsname{100,0\%}
\expandafter\def\csname val@ca-XGB-DDoS-on-DDoS-fa\endcsname{0,00\%}
\expandafter\def\csname val@dup-rows\endcsname{1\,226\,290}
\expandafter\def\csname val@dup-share\endcsname{39,0\%}
\expandafter\def\csname val@filtered-Bot\endcsname{67}
\expandafter\def\csname val@filtered-BruteForce\endcsname{67}
\expandafter\def\csname val@filtered-DDoS\endcsname{67}
\expandafter\def\csname val@gap-max\endcsname{0,016}
\expandafter\def\csname val@hard-attack\endcsname{Bot}
\expandafter\def\csname val@hyb-balanced-Bot-acc\endcsname{99,938\%}
\expandafter\def\csname val@hyb-balanced-Bot-fn\endcsname{6}
\expandafter\def\csname val@hyb-balanced-Bot-fp\endcsname{31}
\expandafter\def\csname val@hyb-balanced-Bot-prec\endcsname{99,90\%}
\expandafter\def\csname val@hyb-balanced-Bot-sens\endcsname{99,98\%}
\expandafter\def\csname val@hyb-balanced-BruteForce-acc\endcsname{99,996\%}
\expandafter\def\csname val@hyb-balanced-BruteForce-fn\endcsname{1}
\expandafter\def\csname val@hyb-balanced-BruteForce-fp\endcsname{1}
\expandafter\def\csname val@hyb-balanced-BruteForce-prec\endcsname{100,00\%}
\expandafter\def\csname val@hyb-balanced-BruteForce-sens\endcsname{100,00\%}
\expandafter\def\csname val@hyb-balanced-DDoS-acc\endcsname{99,970\%}
\expandafter\def\csname val@hyb-balanced-DDoS-fn\endcsname{0}
\expandafter\def\csname val@hyb-balanced-DDoS-fp\endcsname{18}
\expandafter\def\csname val@hyb-balanced-DDoS-prec\endcsname{99,94\%}
\expandafter\def\csname val@hyb-balanced-DDoS-sens\endcsname{100,00\%}
\expandafter\def\csname val@hyb-natural-Bot-acc\endcsname{99,902\%}
\expandafter\def\csname val@hyb-natural-Bot-fn\endcsname{6}
\expandafter\def\csname val@hyb-natural-Bot-fp\endcsname{190}
\expandafter\def\csname val@hyb-natural-Bot-prec\endcsname{99,32\%}
\expandafter\def\csname val@hyb-natural-Bot-sens\endcsname{99,98\%}
\expandafter\def\csname val@hyb-natural-BruteForce-acc\endcsname{99,992\%}
\expandafter\def\csname val@hyb-natural-BruteForce-fn\endcsname{1}
\expandafter\def\csname val@hyb-natural-BruteForce-fp\endcsname{12}
\expandafter\def\csname val@hyb-natural-BruteForce-prec\endcsname{99,96\%}
\expandafter\def\csname val@hyb-natural-BruteForce-sens\endcsname{100,00\%}
\expandafter\def\csname val@hyb-natural-DDoS-acc\endcsname{99,942\%}
\expandafter\def\csname val@hyb-natural-DDoS-fn\endcsname{0}
\expandafter\def\csname val@hyb-natural-DDoS-fp\endcsname{116}
\expandafter\def\csname val@hyb-natural-DDoS-prec\endcsname{99,79\%}
\expandafter\def\csname val@hyb-natural-DDoS-sens\endcsname{100,00\%}
\expandafter\def\csname val@ks-median-Bot\endcsname{0,51}
\expandafter\def\csname val@ks-median-BruteForce\endcsname{0,52}
\expandafter\def\csname val@ks-median-DDoS\endcsname{0,39}
\expandafter\def\csname val@ks-n05-Bot\endcsname{40}
\expandafter\def\csname val@ks-n05-BruteForce\endcsname{44}
\expandafter\def\csname val@ks-n05-DDoS\endcsname{16}
\expandafter\def\csname val@ks-top-Bot\endcsname{Bwd Pkt Len Mean}
\expandafter\def\csname val@ks-top-BruteForce\endcsname{Fwd Seg Size Min}
\expandafter\def\csname val@ks-top-DDoS\endcsname{Bwd Pkt Len Std}
\expandafter\def\csname val@ks-top5-Bot\endcsname{0,75}
\expandafter\def\csname val@ks-top5-BruteForce\endcsname{0,99}
\expandafter\def\csname val@ks-top5-DDoS\endcsname{0,72}
\expandafter\def\csname val@ks-top5names-Bot\endcsname{Bwd Pkt Len Mean, Bwd Seg Size Avg, Pkt Size Avg, Fwd IAT Tot, Fwd IAT Max}
\expandafter\def\csname val@ks-top5names-BruteForce\endcsname{Fwd Seg Size Min, Fwd Act Data Pkts, Fwd Header Len, Subflow Fwd Byts, TotLen Fwd Pkts}
\expandafter\def\csname val@ks-top5names-DDoS\endcsname{Bwd Pkt Len Std, Bwd Seg Size Avg, Bwd Pkt Len Mean, Fwd IAT Min, Init Fwd Win Byts}
\expandafter\def\csname val@ks-topD-Bot\endcsname{0,78}
\expandafter\def\csname val@ks-topD-BruteForce\endcsname{1,00}
\expandafter\def\csname val@ks-topD-DDoS\endcsname{0,74}
\expandafter\def\csname val@lc-RF-Bot-first\endcsname{99,65\%}
\expandafter\def\csname val@lc-RF-Bot-gain-last\endcsname{0,03}
\expandafter\def\csname val@lc-RF-Bot-last\endcsname{99,97\%}
\expandafter\def\csname val@lc-RF-Bot-nfirst\endcsname{1\,000}
\expandafter\def\csname val@lc-RF-Bot-nlast\endcsname{100\,000}
\expandafter\def\csname val@lc-RF-Bot-tlast\endcsname{4,1}
\expandafter\def\csname val@lc-RF-BruteForce-first\endcsname{99,87\%}
\expandafter\def\csname val@lc-RF-BruteForce-gain-last\endcsname{0,00}
\expandafter\def\csname val@lc-RF-BruteForce-last\endcsname{100,00\%}
\expandafter\def\csname val@lc-RF-BruteForce-nfirst\endcsname{1\,000}
\expandafter\def\csname val@lc-RF-BruteForce-nlast\endcsname{100\,000}
\expandafter\def\csname val@lc-RF-BruteForce-tlast\endcsname{3,6}
\expandafter\def\csname val@lc-RF-DDoS-first\endcsname{99,93\%}
\expandafter\def\csname val@lc-RF-DDoS-gain-last\endcsname{0,00}
\expandafter\def\csname val@lc-RF-DDoS-last\endcsname{100,00\%}
\expandafter\def\csname val@lc-RF-DDoS-nfirst\endcsname{1\,000}
\expandafter\def\csname val@lc-RF-DDoS-nlast\endcsname{100\,000}
\expandafter\def\csname val@lc-RF-DDoS-tlast\endcsname{6,6}
\expandafter\def\csname val@lc-SVM-Bot-first\endcsname{99,36\%}
\expandafter\def\csname val@lc-SVM-Bot-gain-last\endcsname{0,04}
\expandafter\def\csname val@lc-SVM-Bot-last\endcsname{99,57\%}
\expandafter\def\csname val@lc-SVM-Bot-nfirst\endcsname{1\,000}
\expandafter\def\csname val@lc-SVM-Bot-nlast\endcsname{50\,000}
\expandafter\def\csname val@lc-SVM-Bot-tlast\endcsname{2,2}
\expandafter\def\csname val@lc-SVM-BruteForce-first\endcsname{99,82\%}
\expandafter\def\csname val@lc-SVM-BruteForce-gain-last\endcsname{0,03}
\expandafter\def\csname val@lc-SVM-BruteForce-last\endcsname{99,98\%}
\expandafter\def\csname val@lc-SVM-BruteForce-nfirst\endcsname{1\,000}
\expandafter\def\csname val@lc-SVM-BruteForce-nlast\endcsname{50\,000}
\expandafter\def\csname val@lc-SVM-BruteForce-tlast\endcsname{1,1}
\expandafter\def\csname val@lc-SVM-DDoS-first\endcsname{99,60\%}
\expandafter\def\csname val@lc-SVM-DDoS-gain-last\endcsname{0,01}
\expandafter\def\csname val@lc-SVM-DDoS-last\endcsname{99,98\%}
\expandafter\def\csname val@lc-SVM-DDoS-nfirst\endcsname{1\,000}
\expandafter\def\csname val@lc-SVM-DDoS-nlast\endcsname{50\,000}
\expandafter\def\csname val@lc-SVM-DDoS-tlast\endcsname{1,6}
\expandafter\def\csname val@lc-XGB-Bot-first\endcsname{99,42\%}
\expandafter\def\csname val@lc-XGB-Bot-gain-last\endcsname{-0,01}
\expandafter\def\csname val@lc-XGB-Bot-last\endcsname{99,97\%}
\expandafter\def\csname val@lc-XGB-Bot-nfirst\endcsname{1\,000}
\expandafter\def\csname val@lc-XGB-Bot-nlast\endcsname{100\,000}
\expandafter\def\csname val@lc-XGB-Bot-tlast\endcsname{3,8}
\expandafter\def\csname val@lc-XGB-BruteForce-first\endcsname{99,84\%}
\expandafter\def\csname val@lc-XGB-BruteForce-gain-last\endcsname{0,00}
\expandafter\def\csname val@lc-XGB-BruteForce-last\endcsname{100,00\%}
\expandafter\def\csname val@lc-XGB-BruteForce-nfirst\endcsname{1\,000}
\expandafter\def\csname val@lc-XGB-BruteForce-nlast\endcsname{100\,000}
\expandafter\def\csname val@lc-XGB-BruteForce-tlast\endcsname{2,4}
\expandafter\def\csname val@lc-XGB-DDoS-first\endcsname{99,87\%}
\expandafter\def\csname val@lc-XGB-DDoS-gain-last\endcsname{0,00}
\expandafter\def\csname val@lc-XGB-DDoS-last\endcsname{100,00\%}
\expandafter\def\csname val@lc-XGB-DDoS-nfirst\endcsname{1\,000}
\expandafter\def\csname val@lc-XGB-DDoS-nlast\endcsname{100\,000}
\expandafter\def\csname val@lc-XGB-DDoS-tlast\endcsname{2,6}
\expandafter\def\csname val@lowvar-Bot\endcsname{10}
\expandafter\def\csname val@lowvar-BruteForce\endcsname{10}
\expandafter\def\csname val@lowvar-DDoS\endcsname{10}
\expandafter\def\csname val@lowvar-names\endcsname{Bwd PSH Flags, Fwd URG Flags, Bwd URG Flags, CWE Flag Count, Fwd Byts/b Avg, Fwd Pkts/b Avg, Fwd Blk Rate Avg, Bwd Byts/b Avg, Bwd Pkts/b Avg, Bwd Blk Rate Avg}
\expandafter\def\csname val@mcn-Bot-RF-SVM-accuracy_diff\endcsname{0,368}
\expandafter\def\csname val@mcn-Bot-RF-SVM-accuracy_diff-ci\endcsname{[0,322; 0,417]}
\expandafter\def\csname val@mcn-Bot-RF-SVM-b\endcsname{228}
\expandafter\def\csname val@mcn-Bot-RF-SVM-both-wrong\endcsname{13}
\expandafter\def\csname val@mcn-Bot-RF-SVM-c\endcsname{7}
\expandafter\def\csname val@mcn-Bot-RF-SVM-chi2\endcsname{206,0}
\expandafter\def\csname val@mcn-Bot-RF-SVM-f1_diff\endcsname{0,370}
\expandafter\def\csname val@mcn-Bot-RF-SVM-f1_diff-ci\endcsname{[0,323; 0,419]}
\expandafter\def\csname val@mcn-Bot-RF-SVM-p\endcsname{$1{,}0\cdot10^{-46}$}
\expandafter\def\csname val@mcn-Bot-RF-SVM-precision_diff\endcsname{0,064}
\expandafter\def\csname val@mcn-Bot-RF-SVM-precision_diff-ci\endcsname{[0,030; 0,104]}
\expandafter\def\csname val@mcn-Bot-RF-SVM-sensitivity_diff\endcsname{0,673}
\expandafter\def\csname val@mcn-Bot-RF-SVM-sensitivity_diff-ci\endcsname{[0,583; 0,762]}
\expandafter\def\csname val@mcn-Bot-RF-SVM-sig\endcsname{значима}
\expandafter\def\csname val@mcn-Bot-RF-XGB-accuracy_diff\endcsname{-0,015}
\expandafter\def\csname val@mcn-Bot-RF-XGB-accuracy_diff-ci\endcsname{[-0,027; -0,003]}
\expandafter\def\csname val@mcn-Bot-RF-XGB-b\endcsname{2}
\expandafter\def\csname val@mcn-Bot-RF-XGB-both-wrong\endcsname{9}
\expandafter\def\csname val@mcn-Bot-RF-XGB-c\endcsname{11}
\expandafter\def\csname val@mcn-Bot-RF-XGB-chi2\endcsname{2,0}
\expandafter\def\csname val@mcn-Bot-RF-XGB-f1_diff\endcsname{-0,015}
\expandafter\def\csname val@mcn-Bot-RF-XGB-f1_diff-ci\endcsname{[-0,027; -0,003]}
\expandafter\def\csname val@mcn-Bot-RF-XGB-p\endcsname{$0{,}0225$}
\expandafter\def\csname val@mcn-Bot-RF-XGB-precision_diff\endcsname{-0,010}
\expandafter\def\csname val@mcn-Bot-RF-XGB-precision_diff-ci\endcsname{[-0,027; 0,007]}
\expandafter\def\csname val@mcn-Bot-RF-XGB-sensitivity_diff\endcsname{-0,020}
\expandafter\def\csname val@mcn-Bot-RF-XGB-sensitivity_diff-ci\endcsname{[-0,037; -0,007]}
\expandafter\def\csname val@mcn-Bot-RF-XGB-sig\endcsname{значима}
\expandafter\def\csname val@mcn-Bot-SVM-XGB-accuracy_diff\endcsname{-0,383}
\expandafter\def\csname val@mcn-Bot-SVM-XGB-accuracy_diff-ci\endcsname{[-0,432; -0,332]}
\expandafter\def\csname val@mcn-Bot-SVM-XGB-b\endcsname{4}
\expandafter\def\csname val@mcn-Bot-SVM-XGB-both-wrong\endcsname{7}
\expandafter\def\csname val@mcn-Bot-SVM-XGB-c\endcsname{234}
\expandafter\def\csname val@mcn-Bot-SVM-XGB-chi2\endcsname{220,3}
\expandafter\def\csname val@mcn-Bot-SVM-XGB-f1_diff\endcsname{-0,385}
\expandafter\def\csname val@mcn-Bot-SVM-XGB-f1_diff-ci\endcsname{[-0,433; -0,333]}
\expandafter\def\csname val@mcn-Bot-SVM-XGB-p\endcsname{$7{,}6\cdot10^{-50}$}
\expandafter\def\csname val@mcn-Bot-SVM-XGB-precision_diff\endcsname{-0,074}
\expandafter\def\csname val@mcn-Bot-SVM-XGB-precision_diff-ci\endcsname{[-0,111; -0,040]}
\expandafter\def\csname val@mcn-Bot-SVM-XGB-sensitivity_diff\endcsname{-0,693}
\expandafter\def\csname val@mcn-Bot-SVM-XGB-sensitivity_diff-ci\endcsname{[-0,784; -0,598]}
\expandafter\def\csname val@mcn-Bot-SVM-XGB-sig\endcsname{значима}
\expandafter\def\csname val@mcn-BruteForce-RF-SVM-accuracy_diff\endcsname{0,012}
\expandafter\def\csname val@mcn-BruteForce-RF-SVM-accuracy_diff-ci\endcsname{[0,005; 0,023]}
\expandafter\def\csname val@mcn-BruteForce-RF-SVM-b\endcsname{7}
\expandafter\def\csname val@mcn-BruteForce-RF-SVM-both-wrong\endcsname{1}
\expandafter\def\csname val@mcn-BruteForce-RF-SVM-c\endcsname{0}
\expandafter\def\csname val@mcn-BruteForce-RF-SVM-chi2\endcsname{0,0}
\expandafter\def\csname val@mcn-BruteForce-RF-SVM-f1_diff\endcsname{0,012}
\expandafter\def\csname val@mcn-BruteForce-RF-SVM-f1_diff-ci\endcsname{[0,005; 0,023]}
\expandafter\def\csname val@mcn-BruteForce-RF-SVM-p\endcsname{$0{,}0156$}
\expandafter\def\csname val@mcn-BruteForce-RF-SVM-precision_diff\endcsname{0,004}
\expandafter\def\csname val@mcn-BruteForce-RF-SVM-precision_diff-ci\endcsname{[0,000; 0,011]}
\expandafter\def\csname val@mcn-BruteForce-RF-SVM-sensitivity_diff\endcsname{0,021}
\expandafter\def\csname val@mcn-BruteForce-RF-SVM-sensitivity_diff-ci\endcsname{[0,007; 0,039]}
\expandafter\def\csname val@mcn-BruteForce-RF-SVM-sig\endcsname{значима}
\expandafter\def\csname val@mcn-BruteForce-RF-XGB-accuracy_diff\endcsname{0,000}
\expandafter\def\csname val@mcn-BruteForce-RF-XGB-accuracy_diff-ci\endcsname{[0,000; 0,000]}
\expandafter\def\csname val@mcn-BruteForce-RF-XGB-b\endcsname{0}
\expandafter\def\csname val@mcn-BruteForce-RF-XGB-both-wrong\endcsname{1}
\expandafter\def\csname val@mcn-BruteForce-RF-XGB-c\endcsname{0}
\expandafter\def\csname val@mcn-BruteForce-RF-XGB-chi2\endcsname{0,0}
\expandafter\def\csname val@mcn-BruteForce-RF-XGB-f1_diff\endcsname{0,000}
\expandafter\def\csname val@mcn-BruteForce-RF-XGB-f1_diff-ci\endcsname{[0,000; 0,000]}
\expandafter\def\csname val@mcn-BruteForce-RF-XGB-p\endcsname{$1{,}0000$}
\expandafter\def\csname val@mcn-BruteForce-RF-XGB-precision_diff\endcsname{0,000}
\expandafter\def\csname val@mcn-BruteForce-RF-XGB-precision_diff-ci\endcsname{[0,000; 0,000]}
\expandafter\def\csname val@mcn-BruteForce-RF-XGB-sensitivity_diff\endcsname{0,000}
\expandafter\def\csname val@mcn-BruteForce-RF-XGB-sensitivity_diff-ci\endcsname{[0,000; 0,000]}
\expandafter\def\csname val@mcn-BruteForce-RF-XGB-sig\endcsname{незначима}
\expandafter\def\csname val@mcn-BruteForce-SVM-XGB-accuracy_diff\endcsname{-0,012}
\expandafter\def\csname val@mcn-BruteForce-SVM-XGB-accuracy_diff-ci\endcsname{[-0,023; -0,005]}
\expandafter\def\csname val@mcn-BruteForce-SVM-XGB-b\endcsname{0}
\expandafter\def\csname val@mcn-BruteForce-SVM-XGB-both-wrong\endcsname{1}
\expandafter\def\csname val@mcn-BruteForce-SVM-XGB-c\endcsname{7}
\expandafter\def\csname val@mcn-BruteForce-SVM-XGB-chi2\endcsname{0,0}
\expandafter\def\csname val@mcn-BruteForce-SVM-XGB-f1_diff\endcsname{-0,012}
\expandafter\def\csname val@mcn-BruteForce-SVM-XGB-f1_diff-ci\endcsname{[-0,023; -0,005]}
\expandafter\def\csname val@mcn-BruteForce-SVM-XGB-p\endcsname{$0{,}0156$}
\expandafter\def\csname val@mcn-BruteForce-SVM-XGB-precision_diff\endcsname{-0,004}
\expandafter\def\csname val@mcn-BruteForce-SVM-XGB-precision_diff-ci\endcsname{[-0,011; 0,000]}
\expandafter\def\csname val@mcn-BruteForce-SVM-XGB-sensitivity_diff\endcsname{-0,021}
\expandafter\def\csname val@mcn-BruteForce-SVM-XGB-sensitivity_diff-ci\endcsname{[-0,039; -0,007]}
\expandafter\def\csname val@mcn-BruteForce-SVM-XGB-sig\endcsname{значима}
\expandafter\def\csname val@mcn-DDoS-RF-SVM-accuracy_diff\endcsname{0,050}
\expandafter\def\csname val@mcn-DDoS-RF-SVM-accuracy_diff-ci\endcsname{[0,033; 0,068]}
\expandafter\def\csname val@mcn-DDoS-RF-SVM-b\endcsname{30}
\expandafter\def\csname val@mcn-DDoS-RF-SVM-both-wrong\endcsname{0}
\expandafter\def\csname val@mcn-DDoS-RF-SVM-c\endcsname{0}
\expandafter\def\csname val@mcn-DDoS-RF-SVM-chi2\endcsname{28,0}
\expandafter\def\csname val@mcn-DDoS-RF-SVM-f1_diff\endcsname{0,050}
\expandafter\def\csname val@mcn-DDoS-RF-SVM-f1_diff-ci\endcsname{[0,033; 0,068]}
\expandafter\def\csname val@mcn-DDoS-RF-SVM-p\endcsname{$1{,}2\cdot10^{-7}$}
\expandafter\def\csname val@mcn-DDoS-RF-SVM-precision_diff\endcsname{0,060}
\expandafter\def\csname val@mcn-DDoS-RF-SVM-precision_diff-ci\endcsname{[0,036; 0,090]}
\expandafter\def\csname val@mcn-DDoS-RF-SVM-sensitivity_diff\endcsname{0,040}
\expandafter\def\csname val@mcn-DDoS-RF-SVM-sensitivity_diff-ci\endcsname{[0,020; 0,063]}
\expandafter\def\csname val@mcn-DDoS-RF-SVM-sig\endcsname{значима}
\expandafter\def\csname val@mcn-DDoS-RF-XGB-accuracy_diff\endcsname{0,000}
\expandafter\def\csname val@mcn-DDoS-RF-XGB-accuracy_diff-ci\endcsname{[0,000; 0,000]}
\expandafter\def\csname val@mcn-DDoS-RF-XGB-b\endcsname{0}
\expandafter\def\csname val@mcn-DDoS-RF-XGB-both-wrong\endcsname{0}
\expandafter\def\csname val@mcn-DDoS-RF-XGB-c\endcsname{0}
\expandafter\def\csname val@mcn-DDoS-RF-XGB-chi2\endcsname{0,0}
\expandafter\def\csname val@mcn-DDoS-RF-XGB-f1_diff\endcsname{0,000}
\expandafter\def\csname val@mcn-DDoS-RF-XGB-f1_diff-ci\endcsname{[0,000; 0,000]}
\expandafter\def\csname val@mcn-DDoS-RF-XGB-p\endcsname{$1{,}0000$}
\expandafter\def\csname val@mcn-DDoS-RF-XGB-precision_diff\endcsname{0,000}
\expandafter\def\csname val@mcn-DDoS-RF-XGB-precision_diff-ci\endcsname{[0,000; 0,000]}
\expandafter\def\csname val@mcn-DDoS-RF-XGB-sensitivity_diff\endcsname{0,000}
\expandafter\def\csname val@mcn-DDoS-RF-XGB-sensitivity_diff-ci\endcsname{[0,000; 0,000]}
\expandafter\def\csname val@mcn-DDoS-RF-XGB-sig\endcsname{незначима}
\expandafter\def\csname val@mcn-DDoS-SVM-XGB-accuracy_diff\endcsname{-0,050}
\expandafter\def\csname val@mcn-DDoS-SVM-XGB-accuracy_diff-ci\endcsname{[-0,068; -0,033]}
\expandafter\def\csname val@mcn-DDoS-SVM-XGB-b\endcsname{0}
\expandafter\def\csname val@mcn-DDoS-SVM-XGB-both-wrong\endcsname{0}
\expandafter\def\csname val@mcn-DDoS-SVM-XGB-c\endcsname{30}
\expandafter\def\csname val@mcn-DDoS-SVM-XGB-chi2\endcsname{28,0}
\expandafter\def\csname val@mcn-DDoS-SVM-XGB-f1_diff\endcsname{-0,050}
\expandafter\def\csname val@mcn-DDoS-SVM-XGB-f1_diff-ci\endcsname{[-0,068; -0,033]}
\expandafter\def\csname val@mcn-DDoS-SVM-XGB-p\endcsname{$1{,}2\cdot10^{-7}$}
\expandafter\def\csname val@mcn-DDoS-SVM-XGB-precision_diff\endcsname{-0,060}
\expandafter\def\csname val@mcn-DDoS-SVM-XGB-precision_diff-ci\endcsname{[-0,090; -0,036]}
\expandafter\def\csname val@mcn-DDoS-SVM-XGB-sensitivity_diff\endcsname{-0,040}
\expandafter\def\csname val@mcn-DDoS-SVM-XGB-sensitivity_diff-ci\endcsname{[-0,063; -0,020]}
\expandafter\def\csname val@mcn-DDoS-SVM-XGB-sig\endcsname{значима}
\expandafter\def\csname val@mcn-n\endcsname{9}
\expandafter\def\csname val@mcn-n-sig\endcsname{7}
\expandafter\def\csname val@mdi-Bot-top3\endcsname{Bwd Seg Size Avg, Bwd Pkt Len Mean, ECE Flag Cnt}
\expandafter\def\csname val@mdi-Bot-top5share\endcsname{35\%}
\expandafter\def\csname val@mdi-BruteForce-top3\endcsname{Subflow Fwd Byts, Fwd Seg Size Min, Fwd Header Len}
\expandafter\def\csname val@mdi-BruteForce-top5share\endcsname{47\%}
\expandafter\def\csname val@mdi-DDoS-top3\endcsname{Bwd Pkt Len Std, Init Fwd Win Byts, Bwd Pkt Len Mean}
\expandafter\def\csname val@mdi-DDoS-top5share\endcsname{67\%}
\expandafter\def\csname val@n-features\endcsname{77}
\expandafter\def\csname val@n-files\endcsname{3}
\expandafter\def\csname val@n-filtered\endcsname{67}
\expandafter\def\csname val@n-filtered-range\endcsname{67}
\expandafter\def\csname val@n-selected\endcsname{20}
\expandafter\def\csname val@nan-rows\endcsname{7\,874}
\expandafter\def\csname val@nat-share-Bot\endcsname{13,8\%}
\expandafter\def\csname val@nat-share-BruteForce\endcsname{17,5\%}
\expandafter\def\csname val@nat-share-DDoS\endcsname{27,7\%}
\expandafter\def\csname val@op-RF-Bot-50-alarms\endcsname{1886}
\expandafter\def\csname val@op-RF-Bot-50-recall\endcsname{100,0\%}
\expandafter\def\csname val@op-RF-Bot-80-alarms\endcsname{1645}
\expandafter\def\csname val@op-RF-Bot-80-recall\endcsname{100,0\%}
\expandafter\def\csname val@op-RF-Bot-90-alarms\endcsname{1530}
\expandafter\def\csname val@op-RF-Bot-90-recall\endcsname{100,0\%}
\expandafter\def\csname val@op-RF-Bot-95-alarms\endcsname{1450}
\expandafter\def\csname val@op-RF-Bot-95-recall\endcsname{100,0\%}
\expandafter\def\csname val@op-RF-Bot-99-alarms\endcsname{1392}
\expandafter\def\csname val@op-RF-Bot-99-recall\endcsname{100,0\%}
\expandafter\def\csname val@op-RF-BruteForce-50-alarms\endcsname{2027}
\expandafter\def\csname val@op-RF-BruteForce-50-recall\endcsname{100,0\%}
\expandafter\def\csname val@op-RF-BruteForce-80-alarms\endcsname{2027}
\expandafter\def\csname val@op-RF-BruteForce-80-recall\endcsname{100,0\%}
\expandafter\def\csname val@op-RF-BruteForce-90-alarms\endcsname{1940}
\expandafter\def\csname val@op-RF-BruteForce-90-recall\endcsname{100,0\%}
\expandafter\def\csname val@op-RF-BruteForce-95-alarms\endcsname{1842}
\expandafter\def\csname val@op-RF-BruteForce-95-recall\endcsname{100,0\%}
\expandafter\def\csname val@op-RF-BruteForce-99-alarms\endcsname{1772}
\expandafter\def\csname val@op-RF-BruteForce-99-recall\endcsname{100,0\%}
\expandafter\def\csname val@op-RF-DDoS-50-alarms\endcsname{2813}
\expandafter\def\csname val@op-RF-DDoS-50-recall\endcsname{100,0\%}
\expandafter\def\csname val@op-RF-DDoS-80-alarms\endcsname{2813}
\expandafter\def\csname val@op-RF-DDoS-80-recall\endcsname{100,0\%}
\expandafter\def\csname val@op-RF-DDoS-90-alarms\endcsname{2813}
\expandafter\def\csname val@op-RF-DDoS-90-recall\endcsname{100,0\%}
\expandafter\def\csname val@op-RF-DDoS-95-alarms\endcsname{2813}
\expandafter\def\csname val@op-RF-DDoS-95-recall\endcsname{100,0\%}
\expandafter\def\csname val@op-RF-DDoS-99-alarms\endcsname{2803}
\expandafter\def\csname val@op-RF-DDoS-99-recall\endcsname{100,0\%}
\expandafter\def\csname val@op-SVM-Bot-50-alarms\endcsname{2756}
\expandafter\def\csname val@op-SVM-Bot-50-recall\endcsname{100,0\%}
\expandafter\def\csname val@op-SVM-Bot-80-alarms\endcsname{1722}
\expandafter\def\csname val@op-SVM-Bot-80-recall\endcsname{100,0\%}
\expandafter\def\csname val@op-SVM-Bot-90-alarms\endcsname{1531}
\expandafter\def\csname val@op-SVM-Bot-90-recall\endcsname{100,0\%}
\expandafter\def\csname val@op-SVM-Bot-95-alarms\endcsname{1448}
\expandafter\def\csname val@op-SVM-Bot-95-recall\endcsname{99,8\%}
\expandafter\def\csname val@op-SVM-Bot-99-alarms\endcsname{1382}
\expandafter\def\csname val@op-SVM-Bot-99-recall\endcsname{99,3\%}
\expandafter\def\csname val@op-SVM-BruteForce-50-alarms\endcsname{3508}
\expandafter\def\csname val@op-SVM-BruteForce-50-recall\endcsname{100,0\%}
\expandafter\def\csname val@op-SVM-BruteForce-80-alarms\endcsname{2192}
\expandafter\def\csname val@op-SVM-BruteForce-80-recall\endcsname{100,0\%}
\expandafter\def\csname val@op-SVM-BruteForce-90-alarms\endcsname{1949}
\expandafter\def\csname val@op-SVM-BruteForce-90-recall\endcsname{100,0\%}
\expandafter\def\csname val@op-SVM-BruteForce-95-alarms\endcsname{1846}
\expandafter\def\csname val@op-SVM-BruteForce-95-recall\endcsname{100,0\%}
\expandafter\def\csname val@op-SVM-BruteForce-99-alarms\endcsname{1772}
\expandafter\def\csname val@op-SVM-BruteForce-99-recall\endcsname{100,0\%}
\expandafter\def\csname val@op-SVM-DDoS-50-alarms\endcsname{5550}
\expandafter\def\csname val@op-SVM-DDoS-50-recall\endcsname{100,0\%}
\expandafter\def\csname val@op-SVM-DDoS-80-alarms\endcsname{3468}
\expandafter\def\csname val@op-SVM-DDoS-80-recall\endcsname{100,0\%}
\expandafter\def\csname val@op-SVM-DDoS-90-alarms\endcsname{3083}
\expandafter\def\csname val@op-SVM-DDoS-90-recall\endcsname{100,0\%}
\expandafter\def\csname val@op-SVM-DDoS-95-alarms\endcsname{2921}
\expandafter\def\csname val@op-SVM-DDoS-95-recall\endcsname{100,0\%}
\expandafter\def\csname val@op-SVM-DDoS-99-alarms\endcsname{2803}
\expandafter\def\csname val@op-SVM-DDoS-99-recall\endcsname{100,0\%}
\expandafter\def\csname val@op-XGB-Bot-50-alarms\endcsname{2741}
\expandafter\def\csname val@op-XGB-Bot-50-recall\endcsname{100,0\%}
\expandafter\def\csname val@op-XGB-Bot-80-alarms\endcsname{1722}
\expandafter\def\csname val@op-XGB-Bot-80-recall\endcsname{100,0\%}
\expandafter\def\csname val@op-XGB-Bot-90-alarms\endcsname{1531}
\expandafter\def\csname val@op-XGB-Bot-90-recall\endcsname{100,0\%}
\expandafter\def\csname val@op-XGB-Bot-95-alarms\endcsname{1444}
\expandafter\def\csname val@op-XGB-Bot-95-recall\endcsname{100,0\%}
\expandafter\def\csname val@op-XGB-Bot-99-alarms\endcsname{1392}
\expandafter\def\csname val@op-XGB-Bot-99-recall\endcsname{100,0\%}
\expandafter\def\csname val@op-XGB-BruteForce-50-alarms\endcsname{3502}
\expandafter\def\csname val@op-XGB-BruteForce-50-recall\endcsname{100,0\%}
\expandafter\def\csname val@op-XGB-BruteForce-80-alarms\endcsname{2189}
\expandafter\def\csname val@op-XGB-BruteForce-80-recall\endcsname{100,0\%}
\expandafter\def\csname val@op-XGB-BruteForce-90-alarms\endcsname{1949}
\expandafter\def\csname val@op-XGB-BruteForce-90-recall\endcsname{100,0\%}
\expandafter\def\csname val@op-XGB-BruteForce-95-alarms\endcsname{1846}
\expandafter\def\csname val@op-XGB-BruteForce-95-recall\endcsname{100,0\%}
\expandafter\def\csname val@op-XGB-BruteForce-99-alarms\endcsname{1771}
\expandafter\def\csname val@op-XGB-BruteForce-99-recall\endcsname{100,0\%}
\expandafter\def\csname val@op-XGB-DDoS-50-alarms\endcsname{4363}
\expandafter\def\csname val@op-XGB-DDoS-50-recall\endcsname{100,0\%}
\expandafter\def\csname val@op-XGB-DDoS-80-alarms\endcsname{3468}
\expandafter\def\csname val@op-XGB-DDoS-80-recall\endcsname{100,0\%}
\expandafter\def\csname val@op-XGB-DDoS-90-alarms\endcsname{3024}
\expandafter\def\csname val@op-XGB-DDoS-90-recall\endcsname{100,0\%}
\expandafter\def\csname val@op-XGB-DDoS-95-alarms\endcsname{2914}
\expandafter\def\csname val@op-XGB-DDoS-95-recall\endcsname{100,0\%}
\expandafter\def\csname val@op-XGB-DDoS-99-alarms\endcsname{2803}
\expandafter\def\csname val@op-XGB-DDoS-99-recall\endcsname{100,0\%}
\expandafter\def\csname val@pca-Bot\endcsname{23}
\expandafter\def\csname val@pca-BruteForce\endcsname{22}
\expandafter\def\csname val@pca-DDoS\endcsname{22}
\expandafter\def\csname val@per-class-Bot\endcsname{100\,000}
\expandafter\def\csname val@per-class-BruteForce\endcsname{94\,094}
\expandafter\def\csname val@per-class-DDoS\endcsname{100\,000}
\expandafter\def\csname val@pi-RF-Bot-top\endcsname{RST Flag Cnt}
\expandafter\def\csname val@pi-RF-Bot-top3\endcsname{RST Flag Cnt, Fwd IAT Min, Init Fwd Win Byts}
\expandafter\def\csname val@pi-RF-Bot-topv\endcsname{18,44}
\expandafter\def\csname val@pi-RF-BruteForce-top\endcsname{Fwd Seg Size Min}
\expandafter\def\csname val@pi-RF-BruteForce-top3\endcsname{Fwd Seg Size Min, Flow Duration, Tot Fwd Pkts}
\expandafter\def\csname val@pi-RF-BruteForce-topv\endcsname{0,60}
\expandafter\def\csname val@pi-RF-DDoS-top\endcsname{Init Fwd Win Byts}
\expandafter\def\csname val@pi-RF-DDoS-top3\endcsname{Init Fwd Win Byts, Bwd Pkt Len Std, Tot Fwd Pkts}
\expandafter\def\csname val@pi-RF-DDoS-topv\endcsname{29,21}
\expandafter\def\csname val@pi-SVM-Bot-top\endcsname{RST Flag Cnt}
\expandafter\def\csname val@pi-SVM-Bot-top3\endcsname{RST Flag Cnt, Pkt Len Max, Flow IAT Max}
\expandafter\def\csname val@pi-SVM-Bot-topv\endcsname{28,71}
\expandafter\def\csname val@pi-SVM-BruteForce-top\endcsname{Fwd Seg Size Min}
\expandafter\def\csname val@pi-SVM-BruteForce-top3\endcsname{Fwd Seg Size Min, Bwd Pkt Len Max, Bwd Pkt Len Std}
\expandafter\def\csname val@pi-SVM-BruteForce-topv\endcsname{33,85}
\expandafter\def\csname val@pi-SVM-DDoS-top\endcsname{Init Fwd Win Byts}
\expandafter\def\csname val@pi-SVM-DDoS-top3\endcsname{Init Fwd Win Byts, Bwd Pkt Len Std, Pkt Len Std}
\expandafter\def\csname val@pi-SVM-DDoS-topv\endcsname{30,01}
\expandafter\def\csname val@pi-XGB-Bot-top\endcsname{Bwd Pkt Len Mean}
\expandafter\def\csname val@pi-XGB-Bot-top3\endcsname{Bwd Pkt Len Mean, RST Flag Cnt, Fwd Pkt Len Mean}
\expandafter\def\csname val@pi-XGB-Bot-topv\endcsname{3,16}
\expandafter\def\csname val@pi-XGB-BruteForce-top\endcsname{Fwd Seg Size Min}
\expandafter\def\csname val@pi-XGB-BruteForce-top3\endcsname{Fwd Seg Size Min, Flow Duration, Tot Fwd Pkts}
\expandafter\def\csname val@pi-XGB-BruteForce-topv\endcsname{34,90}
\expandafter\def\csname val@pi-XGB-DDoS-top\endcsname{Init Fwd Win Byts}
\expandafter\def\csname val@pi-XGB-DDoS-top3\endcsname{Init Fwd Win Byts, Fwd Act Data Pkts, Bwd Pkt Len Std}
\expandafter\def\csname val@pi-XGB-DDoS-topv\endcsname{29,72}
\expandafter\def\csname val@ps-RF-Bot-full\endcsname{99,84\%}
\expandafter\def\csname val@ps-RF-BruteForce-full\endcsname{99,99\%}
\expandafter\def\csname val@ps-RF-DDoS-full\endcsname{99,93\%}
\expandafter\def\csname val@ps-SVM-Bot-full\endcsname{99,51\%}
\expandafter\def\csname val@ps-SVM-BruteForce-full\endcsname{99,93\%}
\expandafter\def\csname val@ps-SVM-DDoS-full\endcsname{99,98\%}
\expandafter\def\csname val@ps-XGB-Bot-full\endcsname{99,92\%}
\expandafter\def\csname val@ps-XGB-BruteForce-full\endcsname{100,00\%}
\expandafter\def\csname val@ps-XGB-DDoS-full\endcsname{100,00\%}
\expandafter\def\csname val@ps-pca-RF-Bot-best\endcsname{99,90\%}
\expandafter\def\csname val@ps-pca-RF-Bot-bestk\endcsname{30}
\expandafter\def\csname val@ps-pca-RF-Bot-k2\endcsname{99,76\%}
\expandafter\def\csname val@ps-pca-RF-Bot-k5\endcsname{99,80\%}
\expandafter\def\csname val@ps-pca-RF-BruteForce-best\endcsname{99,98\%}
\expandafter\def\csname val@ps-pca-RF-BruteForce-bestk\endcsname{40}
\expandafter\def\csname val@ps-pca-RF-BruteForce-k2\endcsname{99,97\%}
\expandafter\def\csname val@ps-pca-RF-BruteForce-k5\endcsname{99,97\%}
\expandafter\def\csname val@ps-pca-RF-DDoS-best\endcsname{99,96\%}
\expandafter\def\csname val@ps-pca-RF-DDoS-bestk\endcsname{10}
\expandafter\def\csname val@ps-pca-RF-DDoS-k2\endcsname{99,92\%}
\expandafter\def\csname val@ps-pca-RF-DDoS-k5\endcsname{99,95\%}
\expandafter\def\csname val@ps-pca-SVM-Bot-best\endcsname{99,48\%}
\expandafter\def\csname val@ps-pca-SVM-Bot-bestk\endcsname{30}
\expandafter\def\csname val@ps-pca-SVM-Bot-k2\endcsname{99,21\%}
\expandafter\def\csname val@ps-pca-SVM-Bot-k5\endcsname{99,27\%}
\expandafter\def\csname val@ps-pca-SVM-BruteForce-best\endcsname{99,95\%}
\expandafter\def\csname val@ps-pca-SVM-BruteForce-bestk\endcsname{2}
\expandafter\def\csname val@ps-pca-SVM-BruteForce-k2\endcsname{99,95\%}
\expandafter\def\csname val@ps-pca-SVM-BruteForce-k5\endcsname{99,94\%}
\expandafter\def\csname val@ps-pca-SVM-DDoS-best\endcsname{99,98\%}
\expandafter\def\csname val@ps-pca-SVM-DDoS-bestk\endcsname{40}
\expandafter\def\csname val@ps-pca-SVM-DDoS-k2\endcsname{99,42\%}
\expandafter\def\csname val@ps-pca-SVM-DDoS-k5\endcsname{99,73\%}
\expandafter\def\csname val@ps-pca-XGB-Bot-best\endcsname{99,90\%}
\expandafter\def\csname val@ps-pca-XGB-Bot-bestk\endcsname{40}
\expandafter\def\csname val@ps-pca-XGB-Bot-k2\endcsname{99,59\%}
\expandafter\def\csname val@ps-pca-XGB-Bot-k5\endcsname{99,74\%}
\expandafter\def\csname val@ps-pca-XGB-BruteForce-best\endcsname{99,94\%}
\expandafter\def\csname val@ps-pca-XGB-BruteForce-bestk\endcsname{40}
\expandafter\def\csname val@ps-pca-XGB-BruteForce-k2\endcsname{99,85\%}
\expandafter\def\csname val@ps-pca-XGB-BruteForce-k5\endcsname{99,89\%}
\expandafter\def\csname val@ps-pca-XGB-DDoS-best\endcsname{99,97\%}
\expandafter\def\csname val@ps-pca-XGB-DDoS-bestk\endcsname{5}
\expandafter\def\csname val@ps-pca-XGB-DDoS-k2\endcsname{99,93\%}
\expandafter\def\csname val@ps-pca-XGB-DDoS-k5\endcsname{99,97\%}
\expandafter\def\csname val@ps-select-RF-Bot-best\endcsname{99,93\%}
\expandafter\def\csname val@ps-select-RF-Bot-bestk\endcsname{20}
\expandafter\def\csname val@ps-select-RF-Bot-k2\endcsname{99,08\%}
\expandafter\def\csname val@ps-select-RF-Bot-k5\endcsname{99,44\%}
\expandafter\def\csname val@ps-select-RF-BruteForce-best\endcsname{100,00\%}
\expandafter\def\csname val@ps-select-RF-BruteForce-bestk\endcsname{30}
\expandafter\def\csname val@ps-select-RF-BruteForce-k2\endcsname{99,96\%}
\expandafter\def\csname val@ps-select-RF-BruteForce-k5\endcsname{99,97\%}
\expandafter\def\csname val@ps-select-RF-DDoS-best\endcsname{100,00\%}
\expandafter\def\csname val@ps-select-RF-DDoS-bestk\endcsname{5}
\expandafter\def\csname val@ps-select-RF-DDoS-k2\endcsname{94,53\%}
\expandafter\def\csname val@ps-select-RF-DDoS-k5\endcsname{100,00\%}
\expandafter\def\csname val@ps-select-SVM-Bot-best\endcsname{99,53\%}
\expandafter\def\csname val@ps-select-SVM-Bot-bestk\endcsname{20}
\expandafter\def\csname val@ps-select-SVM-Bot-k2\endcsname{97,88\%}
\expandafter\def\csname val@ps-select-SVM-Bot-k5\endcsname{99,34\%}
\expandafter\def\csname val@ps-select-SVM-BruteForce-best\endcsname{99,97\%}
\expandafter\def\csname val@ps-select-SVM-BruteForce-bestk\endcsname{5}
\expandafter\def\csname val@ps-select-SVM-BruteForce-k2\endcsname{99,86\%}
\expandafter\def\csname val@ps-select-SVM-BruteForce-k5\endcsname{99,97\%}
\expandafter\def\csname val@ps-select-SVM-DDoS-best\endcsname{99,97\%}
\expandafter\def\csname val@ps-select-SVM-DDoS-bestk\endcsname{20}
\expandafter\def\csname val@ps-select-SVM-DDoS-k2\endcsname{88,37\%}
\expandafter\def\csname val@ps-select-SVM-DDoS-k5\endcsname{99,86\%}
\expandafter\def\csname val@ps-select-XGB-Bot-best\endcsname{99,94\%}
\expandafter\def\csname val@ps-select-XGB-Bot-bestk\endcsname{20}
\expandafter\def\csname val@ps-select-XGB-Bot-k2\endcsname{99,02\%}
\expandafter\def\csname val@ps-select-XGB-Bot-k5\endcsname{99,43\%}
\expandafter\def\csname val@ps-select-XGB-BruteForce-best\endcsname{100,00\%}
\expandafter\def\csname val@ps-select-XGB-BruteForce-bestk\endcsname{15}
\expandafter\def\csname val@ps-select-XGB-BruteForce-k2\endcsname{99,96\%}
\expandafter\def\csname val@ps-select-XGB-BruteForce-k5\endcsname{99,96\%}
\expandafter\def\csname val@ps-select-XGB-DDoS-best\endcsname{100,00\%}
\expandafter\def\csname val@ps-select-XGB-DDoS-bestk\endcsname{30}
\expandafter\def\csname val@ps-select-XGB-DDoS-k2\endcsname{94,44\%}
\expandafter\def\csname val@ps-select-XGB-DDoS-k5\endcsname{99,99\%}
\expandafter\def\csname val@rob-1day-RF-Bot\endcsname{0,09\%}
\expandafter\def\csname val@rob-1day-RF-BruteForce\endcsname{0,01\%}
\expandafter\def\csname val@rob-1day-RF-DDoS\endcsname{16,49\%}
\expandafter\def\csname val@rob-1day-SVM-Bot\endcsname{0,03\%}
\expandafter\def\csname val@rob-1day-SVM-BruteForce\endcsname{0,02\%}
\expandafter\def\csname val@rob-1day-SVM-DDoS\endcsname{23,36\%}
\expandafter\def\csname val@rob-1day-XGB-Bot\endcsname{9,42\%}
\expandafter\def\csname val@rob-1day-XGB-BruteForce\endcsname{0,01\%}
\expandafter\def\csname val@rob-1day-XGB-DDoS\endcsname{32,73\%}
\expandafter\def\csname val@rob-1day-XGBorSVM-Bot\endcsname{9,43\%}
\expandafter\def\csname val@rob-1day-XGBorSVM-BruteForce\endcsname{0,02\%}
\expandafter\def\csname val@rob-1day-XGBorSVM-DDoS\endcsname{37,29\%}
\expandafter\def\csname val@rob-2day-RF-Bot\endcsname{94,80\%}
\expandafter\def\csname val@rob-2day-RF-BruteForce\endcsname{0,00\%}
\expandafter\def\csname val@rob-2day-RF-DDoS\endcsname{0,00\%}
\expandafter\def\csname val@rob-2day-SVM-Bot\endcsname{0,03\%}
\expandafter\def\csname val@rob-2day-SVM-BruteForce\endcsname{0,02\%}
\expandafter\def\csname val@rob-2day-SVM-DDoS\endcsname{0,12\%}
\expandafter\def\csname val@rob-2day-XGB-Bot\endcsname{21,20\%}
\expandafter\def\csname val@rob-2day-XGB-BruteForce\endcsname{0,01\%}
\expandafter\def\csname val@rob-2day-XGB-DDoS\endcsname{0,01\%}
\expandafter\def\csname val@rob-2day-XGBorSVM-Bot\endcsname{21,20\%}
\expandafter\def\csname val@rob-2day-XGBorSVM-BruteForce\endcsname{0,02\%}
\expandafter\def\csname val@rob-2day-XGBorSVM-DDoS\endcsname{0,12\%}
\expandafter\def\csname val@rob-chrono-min-DDoSBF\endcsname{99,55\%}
\expandafter\def\csname val@rob-chronological-all-RF-Bot-acc\endcsname{94,97\%}
\expandafter\def\csname val@rob-chronological-all-RF-Bot-fpr\endcsname{0,04\%}
\expandafter\def\csname val@rob-chronological-all-RF-Bot-sens\endcsname{90,0\%}
\expandafter\def\csname val@rob-chronological-all-RF-BruteForce-acc\endcsname{99,67\%}
\expandafter\def\csname val@rob-chronological-all-RF-BruteForce-fpr\endcsname{0,62\%}
\expandafter\def\csname val@rob-chronological-all-RF-BruteForce-sens\endcsname{100,0\%}
\expandafter\def\csname val@rob-chronological-all-RF-DDoS-acc\endcsname{99,55\%}
\expandafter\def\csname val@rob-chronological-all-RF-DDoS-fpr\endcsname{0,00\%}
\expandafter\def\csname val@rob-chronological-all-RF-DDoS-sens\endcsname{99,1\%}
\expandafter\def\csname val@rob-chronological-all-SVM-Bot-acc\endcsname{99,72\%}
\expandafter\def\csname val@rob-chronological-all-SVM-Bot-fpr\endcsname{0,12\%}
\expandafter\def\csname val@rob-chronological-all-SVM-Bot-sens\endcsname{99,6\%}
\expandafter\def\csname val@rob-chronological-all-SVM-BruteForce-acc\endcsname{99,65\%}
\expandafter\def\csname val@rob-chronological-all-SVM-BruteForce-fpr\endcsname{0,66\%}
\expandafter\def\csname val@rob-chronological-all-SVM-BruteForce-sens\endcsname{100,0\%}
\expandafter\def\csname val@rob-chronological-all-SVM-DDoS-acc\endcsname{99,91\%}
\expandafter\def\csname val@rob-chronological-all-SVM-DDoS-fpr\endcsname{0,00\%}
\expandafter\def\csname val@rob-chronological-all-SVM-DDoS-sens\endcsname{99,8\%}
\expandafter\def\csname val@rob-chronological-all-XGB-Bot-acc\endcsname{94,98\%}
\expandafter\def\csname val@rob-chronological-all-XGB-Bot-fpr\endcsname{0,05\%}
\expandafter\def\csname val@rob-chronological-all-XGB-Bot-sens\endcsname{90,0\%}
\expandafter\def\csname val@rob-chronological-all-XGB-BruteForce-acc\endcsname{99,68\%}
\expandafter\def\csname val@rob-chronological-all-XGB-BruteForce-fpr\endcsname{0,63\%}
\expandafter\def\csname val@rob-chronological-all-XGB-BruteForce-sens\endcsname{100,0\%}
\expandafter\def\csname val@rob-chronological-all-XGB-DDoS-acc\endcsname{100,00\%}
\expandafter\def\csname val@rob-chronological-all-XGB-DDoS-fpr\endcsname{0,00\%}
\expandafter\def\csname val@rob-chronological-all-XGB-DDoS-sens\endcsname{100,0\%}
\expandafter\def\csname val@rob-chronological-all-XGBorSVM-Bot-acc\endcsname{99,91\%}
\expandafter\def\csname val@rob-chronological-all-XGBorSVM-Bot-fpr\endcsname{0,17\%}
\expandafter\def\csname val@rob-chronological-all-XGBorSVM-Bot-sens\endcsname{100,0\%}
\expandafter\def\csname val@rob-chronological-all-XGBorSVM-BruteForce-acc\endcsname{99,67\%}
\expandafter\def\csname val@rob-chronological-all-XGBorSVM-BruteForce-fpr\endcsname{0,67\%}
\expandafter\def\csname val@rob-chronological-all-XGBorSVM-BruteForce-sens\endcsname{100,0\%}
\expandafter\def\csname val@rob-chronological-all-XGBorSVM-DDoS-acc\endcsname{100,00\%}
\expandafter\def\csname val@rob-chronological-all-XGBorSVM-DDoS-fpr\endcsname{0,00\%}
\expandafter\def\csname val@rob-chronological-all-XGBorSVM-DDoS-sens\endcsname{100,0\%}
\expandafter\def\csname val@rob-chronological-no_fp-RF-Bot-acc\endcsname{94,92\%}
\expandafter\def\csname val@rob-chronological-no_fp-RF-Bot-fpr\endcsname{0,12\%}
\expandafter\def\csname val@rob-chronological-no_fp-RF-Bot-sens\endcsname{90,0\%}
\expandafter\def\csname val@rob-chronological-no_fp-RF-BruteForce-acc\endcsname{99,81\%}
\expandafter\def\csname val@rob-chronological-no_fp-RF-BruteForce-fpr\endcsname{0,33\%}
\expandafter\def\csname val@rob-chronological-no_fp-RF-BruteForce-sens\endcsname{100,0\%}
\expandafter\def\csname val@rob-chronological-no_fp-RF-DDoS-acc\endcsname{99,56\%}
\expandafter\def\csname val@rob-chronological-no_fp-RF-DDoS-fpr\endcsname{0,00\%}
\expandafter\def\csname val@rob-chronological-no_fp-RF-DDoS-sens\endcsname{99,1\%}
\expandafter\def\csname val@rob-chronological-no_fp-SVM-Bot-acc\endcsname{99,83\%}
\expandafter\def\csname val@rob-chronological-no_fp-SVM-Bot-fpr\endcsname{0,11\%}
\expandafter\def\csname val@rob-chronological-no_fp-SVM-Bot-sens\endcsname{99,8\%}
\expandafter\def\csname val@rob-chronological-no_fp-SVM-BruteForce-acc\endcsname{99,78\%}
\expandafter\def\csname val@rob-chronological-no_fp-SVM-BruteForce-fpr\endcsname{0,39\%}
\expandafter\def\csname val@rob-chronological-no_fp-SVM-BruteForce-sens\endcsname{99,9\%}
\expandafter\def\csname val@rob-chronological-no_fp-SVM-DDoS-acc\endcsname{99,91\%}
\expandafter\def\csname val@rob-chronological-no_fp-SVM-DDoS-fpr\endcsname{0,00\%}
\expandafter\def\csname val@rob-chronological-no_fp-SVM-DDoS-sens\endcsname{99,8\%}
\expandafter\def\csname val@rob-chronological-no_fp-XGB-Bot-acc\endcsname{94,91\%}
\expandafter\def\csname val@rob-chronological-no_fp-XGB-Bot-fpr\endcsname{0,13\%}
\expandafter\def\csname val@rob-chronological-no_fp-XGB-Bot-sens\endcsname{89,9\%}
\expandafter\def\csname val@rob-chronological-no_fp-XGB-BruteForce-acc\endcsname{99,97\%}
\expandafter\def\csname val@rob-chronological-no_fp-XGB-BruteForce-fpr\endcsname{0,01\%}
\expandafter\def\csname val@rob-chronological-no_fp-XGB-BruteForce-sens\endcsname{100,0\%}
\expandafter\def\csname val@rob-chronological-no_fp-XGB-DDoS-acc\endcsname{100,00\%}
\expandafter\def\csname val@rob-chronological-no_fp-XGB-DDoS-fpr\endcsname{0,00\%}
\expandafter\def\csname val@rob-chronological-no_fp-XGB-DDoS-sens\endcsname{100,0\%}
\expandafter\def\csname val@rob-chronological-no_fp-XGBorSVM-Bot-acc\endcsname{99,84\%}
\expandafter\def\csname val@rob-chronological-no_fp-XGBorSVM-Bot-fpr\endcsname{0,24\%}
\expandafter\def\csname val@rob-chronological-no_fp-XGBorSVM-Bot-sens\endcsname{99,9\%}
\expandafter\def\csname val@rob-chronological-no_fp-XGBorSVM-BruteForce-acc\endcsname{99,78\%}
\expandafter\def\csname val@rob-chronological-no_fp-XGBorSVM-BruteForce-fpr\endcsname{0,40\%}
\expandafter\def\csname val@rob-chronological-no_fp-XGBorSVM-BruteForce-sens\endcsname{100,0\%}
\expandafter\def\csname val@rob-chronological-no_fp-XGBorSVM-DDoS-acc\endcsname{100,00\%}
\expandafter\def\csname val@rob-chronological-no_fp-XGBorSVM-DDoS-fpr\endcsname{0,00\%}
\expandafter\def\csname val@rob-chronological-no_fp-XGBorSVM-DDoS-sens\endcsname{100,0\%}
\expandafter\def\csname val@rob-cut-Bot\endcsname{14:28}
\expandafter\def\csname val@rob-cut-BruteForce\endcsname{15:05}
\expandafter\def\csname val@rob-cut-DDoS\endcsname{14:15}
\expandafter\def\csname val@rob-fp-maxdrop\endcsname{0,11}
\expandafter\def\csname val@rob-random-all-RF-Bot-acc\endcsname{99,92\%}
\expandafter\def\csname val@rob-random-all-RF-Bot-fpr\endcsname{0,05\%}
\expandafter\def\csname val@rob-random-all-RF-Bot-sens\endcsname{99,9\%}
\expandafter\def\csname val@rob-random-all-RF-BruteForce-acc\endcsname{99,99\%}
\expandafter\def\csname val@rob-random-all-RF-BruteForce-fpr\endcsname{0,01\%}
\expandafter\def\csname val@rob-random-all-RF-BruteForce-sens\endcsname{100,0\%}
\expandafter\def\csname val@rob-random-all-RF-DDoS-acc\endcsname{100,00\%}
\expandafter\def\csname val@rob-random-all-RF-DDoS-fpr\endcsname{0,00\%}
\expandafter\def\csname val@rob-random-all-RF-DDoS-sens\endcsname{100,0\%}
\expandafter\def\csname val@rob-random-all-SVM-Bot-acc\endcsname{99,48\%}
\expandafter\def\csname val@rob-random-all-SVM-Bot-fpr\endcsname{0,16\%}
\expandafter\def\csname val@rob-random-all-SVM-Bot-sens\endcsname{99,1\%}
\expandafter\def\csname val@rob-random-all-SVM-BruteForce-acc\endcsname{99,98\%}
\expandafter\def\csname val@rob-random-all-SVM-BruteForce-fpr\endcsname{0,01\%}
\expandafter\def\csname val@rob-random-all-SVM-BruteForce-sens\endcsname{100,0\%}
\expandafter\def\csname val@rob-random-all-SVM-DDoS-acc\endcsname{99,99\%}
\expandafter\def\csname val@rob-random-all-SVM-DDoS-fpr\endcsname{0,00\%}
\expandafter\def\csname val@rob-random-all-SVM-DDoS-sens\endcsname{100,0\%}
\expandafter\def\csname val@rob-random-all-XGB-Bot-acc\endcsname{99,96\%}
\expandafter\def\csname val@rob-random-all-XGB-Bot-fpr\endcsname{0,03\%}
\expandafter\def\csname val@rob-random-all-XGB-Bot-sens\endcsname{99,9\%}
\expandafter\def\csname val@rob-random-all-XGB-BruteForce-acc\endcsname{99,99\%}
\expandafter\def\csname val@rob-random-all-XGB-BruteForce-fpr\endcsname{0,01\%}
\expandafter\def\csname val@rob-random-all-XGB-BruteForce-sens\endcsname{100,0\%}
\expandafter\def\csname val@rob-random-all-XGB-DDoS-acc\endcsname{100,00\%}
\expandafter\def\csname val@rob-random-all-XGB-DDoS-fpr\endcsname{0,00\%}
\expandafter\def\csname val@rob-random-all-XGB-DDoS-sens\endcsname{100,0\%}
\expandafter\def\csname val@rob-random-all-XGBorSVM-Bot-acc\endcsname{99,88\%}
\expandafter\def\csname val@rob-random-all-XGBorSVM-Bot-fpr\endcsname{0,18\%}
\expandafter\def\csname val@rob-random-all-XGBorSVM-Bot-sens\endcsname{99,9\%}
\expandafter\def\csname val@rob-random-all-XGBorSVM-BruteForce-acc\endcsname{99,99\%}
\expandafter\def\csname val@rob-random-all-XGBorSVM-BruteForce-fpr\endcsname{0,01\%}
\expandafter\def\csname val@rob-random-all-XGBorSVM-BruteForce-sens\endcsname{100,0\%}
\expandafter\def\csname val@rob-random-all-XGBorSVM-DDoS-acc\endcsname{100,00\%}
\expandafter\def\csname val@rob-random-all-XGBorSVM-DDoS-fpr\endcsname{0,00\%}
\expandafter\def\csname val@rob-random-all-XGBorSVM-DDoS-sens\endcsname{100,0\%}
\expandafter\def\csname val@rob-random-no_fp-RF-Bot-acc\endcsname{99,85\%}
\expandafter\def\csname val@rob-random-no_fp-RF-Bot-fpr\endcsname{0,17\%}
\expandafter\def\csname val@rob-random-no_fp-RF-Bot-sens\endcsname{99,9\%}
\expandafter\def\csname val@rob-random-no_fp-RF-BruteForce-acc\endcsname{100,00\%}
\expandafter\def\csname val@rob-random-no_fp-RF-BruteForce-fpr\endcsname{0,00\%}
\expandafter\def\csname val@rob-random-no_fp-RF-BruteForce-sens\endcsname{100,0\%}
\expandafter\def\csname val@rob-random-no_fp-RF-DDoS-acc\endcsname{100,00\%}
\expandafter\def\csname val@rob-random-no_fp-RF-DDoS-fpr\endcsname{0,00\%}
\expandafter\def\csname val@rob-random-no_fp-RF-DDoS-sens\endcsname{100,0\%}
\expandafter\def\csname val@rob-random-no_fp-SVM-Bot-acc\endcsname{99,42\%}
\expandafter\def\csname val@rob-random-no_fp-SVM-Bot-fpr\endcsname{0,06\%}
\expandafter\def\csname val@rob-random-no_fp-SVM-Bot-sens\endcsname{98,9\%}
\expandafter\def\csname val@rob-random-no_fp-SVM-BruteForce-acc\endcsname{99,88\%}
\expandafter\def\csname val@rob-random-no_fp-SVM-BruteForce-fpr\endcsname{0,09\%}
\expandafter\def\csname val@rob-random-no_fp-SVM-BruteForce-sens\endcsname{99,9\%}
\expandafter\def\csname val@rob-random-no_fp-SVM-DDoS-acc\endcsname{99,98\%}
\expandafter\def\csname val@rob-random-no_fp-SVM-DDoS-fpr\endcsname{0,01\%}
\expandafter\def\csname val@rob-random-no_fp-SVM-DDoS-sens\endcsname{100,0\%}
\expandafter\def\csname val@rob-random-no_fp-XGB-Bot-acc\endcsname{99,85\%}
\expandafter\def\csname val@rob-random-no_fp-XGB-Bot-fpr\endcsname{0,17\%}
\expandafter\def\csname val@rob-random-no_fp-XGB-Bot-sens\endcsname{99,9\%}
\expandafter\def\csname val@rob-random-no_fp-XGB-BruteForce-acc\endcsname{99,98\%}
\expandafter\def\csname val@rob-random-no_fp-XGB-BruteForce-fpr\endcsname{0,01\%}
\expandafter\def\csname val@rob-random-no_fp-XGB-BruteForce-sens\endcsname{100,0\%}
\expandafter\def\csname val@rob-random-no_fp-XGB-DDoS-acc\endcsname{100,00\%}
\expandafter\def\csname val@rob-random-no_fp-XGB-DDoS-fpr\endcsname{0,01\%}
\expandafter\def\csname val@rob-random-no_fp-XGB-DDoS-sens\endcsname{100,0\%}
\expandafter\def\csname val@rob-random-no_fp-XGBorSVM-Bot-acc\endcsname{99,82\%}
\expandafter\def\csname val@rob-random-no_fp-XGBorSVM-Bot-fpr\endcsname{0,23\%}
\expandafter\def\csname val@rob-random-no_fp-XGBorSVM-Bot-sens\endcsname{99,9\%}
\expandafter\def\csname val@rob-random-no_fp-XGBorSVM-BruteForce-acc\endcsname{99,93\%}
\expandafter\def\csname val@rob-random-no_fp-XGBorSVM-BruteForce-fpr\endcsname{0,10\%}
\expandafter\def\csname val@rob-random-no_fp-XGBorSVM-BruteForce-sens\endcsname{100,0\%}
\expandafter\def\csname val@rob-random-no_fp-XGBorSVM-DDoS-acc\endcsname{99,99\%}
\expandafter\def\csname val@rob-random-no_fp-XGBorSVM-DDoS-fpr\endcsname{0,01\%}
\expandafter\def\csname val@rob-random-no_fp-XGBorSVM-DDoS-sens\endcsname{100,0\%}
\expandafter\def\csname val@rop-RF-Bot-90-alarms\endcsname{1338}
\expandafter\def\csname val@rop-RF-Bot-90-fp-per-tp\endcsname{0,0}
\expandafter\def\csname val@rop-RF-Bot-90-precision\endcsname{100,0\%}
\expandafter\def\csname val@rop-RF-Bot-90-spec\endcsname{100,0\%}
\expandafter\def\csname val@rop-RF-Bot-95-alarms\endcsname{1338}
\expandafter\def\csname val@rop-RF-Bot-95-fp-per-tp\endcsname{0,0}
\expandafter\def\csname val@rop-RF-Bot-95-precision\endcsname{100,0\%}
\expandafter\def\csname val@rop-RF-Bot-95-spec\endcsname{100,0\%}
\expandafter\def\csname val@rop-RF-BruteForce-90-alarms\endcsname{1745}
\expandafter\def\csname val@rop-RF-BruteForce-90-fp-per-tp\endcsname{0,0}
\expandafter\def\csname val@rop-RF-BruteForce-90-precision\endcsname{100,0\%}
\expandafter\def\csname val@rop-RF-BruteForce-90-spec\endcsname{100,0\%}
\expandafter\def\csname val@rop-RF-BruteForce-95-alarms\endcsname{1745}
\expandafter\def\csname val@rop-RF-BruteForce-95-fp-per-tp\endcsname{0,0}
\expandafter\def\csname val@rop-RF-BruteForce-95-precision\endcsname{100,0\%}
\expandafter\def\csname val@rop-RF-BruteForce-95-spec\endcsname{100,0\%}
\expandafter\def\csname val@rop-RF-DDoS-90-alarms\endcsname{2760}
\expandafter\def\csname val@rop-RF-DDoS-90-fp-per-tp\endcsname{0,0}
\expandafter\def\csname val@rop-RF-DDoS-90-precision\endcsname{100,0\%}
\expandafter\def\csname val@rop-RF-DDoS-90-spec\endcsname{100,0\%}
\expandafter\def\csname val@rop-RF-DDoS-95-alarms\endcsname{2760}
\expandafter\def\csname val@rop-RF-DDoS-95-fp-per-tp\endcsname{0,0}
\expandafter\def\csname val@rop-RF-DDoS-95-precision\endcsname{100,0\%}
\expandafter\def\csname val@rop-RF-DDoS-95-spec\endcsname{100,0\%}
\expandafter\def\csname val@rop-SVM-Bot-90-alarms\endcsname{1241}
\expandafter\def\csname val@rop-SVM-Bot-90-fp-per-tp\endcsname{0,0}
\expandafter\def\csname val@rop-SVM-Bot-90-precision\endcsname{100,0\%}
\expandafter\def\csname val@rop-SVM-Bot-90-spec\endcsname{100,0\%}
\expandafter\def\csname val@rop-SVM-Bot-95-alarms\endcsname{1310}
\expandafter\def\csname val@rop-SVM-Bot-95-fp-per-tp\endcsname{0,0}
\expandafter\def\csname val@rop-SVM-Bot-95-precision\endcsname{100,0\%}
\expandafter\def\csname val@rop-SVM-Bot-95-spec\endcsname{100,0\%}
\expandafter\def\csname val@rop-SVM-BruteForce-90-alarms\endcsname{1579}
\expandafter\def\csname val@rop-SVM-BruteForce-90-fp-per-tp\endcsname{0,0}
\expandafter\def\csname val@rop-SVM-BruteForce-90-precision\endcsname{100,0\%}
\expandafter\def\csname val@rop-SVM-BruteForce-90-spec\endcsname{100,0\%}
\expandafter\def\csname val@rop-SVM-BruteForce-95-alarms\endcsname{1667}
\expandafter\def\csname val@rop-SVM-BruteForce-95-fp-per-tp\endcsname{0,0}
\expandafter\def\csname val@rop-SVM-BruteForce-95-precision\endcsname{100,0\%}
\expandafter\def\csname val@rop-SVM-BruteForce-95-spec\endcsname{100,0\%}
\expandafter\def\csname val@rop-SVM-DDoS-90-alarms\endcsname{2498}
\expandafter\def\csname val@rop-SVM-DDoS-90-fp-per-tp\endcsname{0,0}
\expandafter\def\csname val@rop-SVM-DDoS-90-precision\endcsname{100,0\%}
\expandafter\def\csname val@rop-SVM-DDoS-90-spec\endcsname{100,0\%}
\expandafter\def\csname val@rop-SVM-DDoS-95-alarms\endcsname{2636}
\expandafter\def\csname val@rop-SVM-DDoS-95-fp-per-tp\endcsname{0,0}
\expandafter\def\csname val@rop-SVM-DDoS-95-precision\endcsname{100,0\%}
\expandafter\def\csname val@rop-SVM-DDoS-95-spec\endcsname{100,0\%}
\expandafter\def\csname val@rop-XGB-Bot-90-alarms\endcsname{1240}
\expandafter\def\csname val@rop-XGB-Bot-90-fp-per-tp\endcsname{0,0}
\expandafter\def\csname val@rop-XGB-Bot-90-precision\endcsname{100,0\%}
\expandafter\def\csname val@rop-XGB-Bot-90-spec\endcsname{100,0\%}
\expandafter\def\csname val@rop-XGB-Bot-95-alarms\endcsname{1309}
\expandafter\def\csname val@rop-XGB-Bot-95-fp-per-tp\endcsname{0,0}
\expandafter\def\csname val@rop-XGB-Bot-95-precision\endcsname{100,0\%}
\expandafter\def\csname val@rop-XGB-Bot-95-spec\endcsname{100,0\%}
\expandafter\def\csname val@rop-XGB-BruteForce-90-alarms\endcsname{1587}
\expandafter\def\csname val@rop-XGB-BruteForce-90-fp-per-tp\endcsname{0,0}
\expandafter\def\csname val@rop-XGB-BruteForce-90-precision\endcsname{100,0\%}
\expandafter\def\csname val@rop-XGB-BruteForce-90-spec\endcsname{100,0\%}
\expandafter\def\csname val@rop-XGB-BruteForce-95-alarms\endcsname{1667}
\expandafter\def\csname val@rop-XGB-BruteForce-95-fp-per-tp\endcsname{0,0}
\expandafter\def\csname val@rop-XGB-BruteForce-95-precision\endcsname{100,0\%}
\expandafter\def\csname val@rop-XGB-BruteForce-95-spec\endcsname{100,0\%}
\expandafter\def\csname val@rop-XGB-DDoS-90-alarms\endcsname{2521}
\expandafter\def\csname val@rop-XGB-DDoS-90-fp-per-tp\endcsname{0,0}
\expandafter\def\csname val@rop-XGB-DDoS-90-precision\endcsname{100,0\%}
\expandafter\def\csname val@rop-XGB-DDoS-90-spec\endcsname{100,0\%}
\expandafter\def\csname val@rop-XGB-DDoS-95-alarms\endcsname{2641}
\expandafter\def\csname val@rop-XGB-DDoS-95-fp-per-tp\endcsname{0,0}
\expandafter\def\csname val@rop-XGB-DDoS-95-precision\endcsname{100,0\%}
\expandafter\def\csname val@rop-XGB-DDoS-95-spec\endcsname{100,0\%}
\expandafter\def\csname val@rows-clean\endcsname{1\,911\,561}
\expandafter\def\csname val@rows-read\endcsname{3\,145\,725}
\expandafter\def\csname val@rows-read-mln\endcsname{3,15}
\expandafter\def\csname val@sel-share-Bot\endcsname{99,5\%}
\expandafter\def\csname val@sel-share-BruteForce\endcsname{100,0\%}
\expandafter\def\csname val@sel-share-DDoS\endcsname{99,9\%}
\expandafter\def\csname val@sel-top-Bot\endcsname{Bwd Pkt Len Mean}
\expandafter\def\csname val@sel-top-BruteForce\endcsname{Fwd Seg Size Min}
\expandafter\def\csname val@sel-top-DDoS\endcsname{Fwd IAT Min}
\expandafter\def\csname val@sel-top3-Bot\endcsname{Bwd Pkt Len Mean, RST Flag Cnt, Fwd Pkt Len Mean}
\expandafter\def\csname val@sel-top3-BruteForce\endcsname{Fwd Seg Size Min, Flow Pkts/s, Init Bwd Win Byts}
\expandafter\def\csname val@sel-top3-DDoS\endcsname{Fwd IAT Min, Init Fwd Win Byts, Flow Duration}
\expandafter\def\csname val@sel-topv-Bot\endcsname{64,5\%}
\expandafter\def\csname val@sel-topv-BruteForce\endcsname{98,6\%}
\expandafter\def\csname val@sel-topv-DDoS\endcsname{61,1\%}
\expandafter\def\csname val@sub-DDOSattack-HOIC-n\endcsname{29\,754}
\expandafter\def\csname val@sub-DDOSattack-LOIC-UDP-n\endcsname{246}
\expandafter\def\csname val@sub-FTP-BruteForce-n\endcsname{18}
\expandafter\def\csname val@sub-FTP-SSH-RF-all\endcsname{0,2\%}
\expandafter\def\csname val@sub-FTP-SSH-RF-no_fp\endcsname{0,2\%}
\expandafter\def\csname val@sub-FTP-SSH-SVM-all\endcsname{0,2\%}
\expandafter\def\csname val@sub-FTP-SSH-SVM-no_fp\endcsname{0,2\%}
\expandafter\def\csname val@sub-FTP-SSH-XGB-all\endcsname{0,2\%}
\expandafter\def\csname val@sub-FTP-SSH-XGB-no_fp\endcsname{0,2\%}
\expandafter\def\csname val@sub-HOIC-LOIC-UDP-RF-all\endcsname{0,0\%}
\expandafter\def\csname val@sub-HOIC-LOIC-UDP-RF-no_fp\endcsname{0,0\%}
\expandafter\def\csname val@sub-HOIC-LOIC-UDP-SVM-all\endcsname{0,0\%}
\expandafter\def\csname val@sub-HOIC-LOIC-UDP-SVM-no_fp\endcsname{0,0\%}
\expandafter\def\csname val@sub-HOIC-LOIC-UDP-XGB-all\endcsname{0,0\%}
\expandafter\def\csname val@sub-HOIC-LOIC-UDP-XGB-no_fp\endcsname{0,0\%}
\expandafter\def\csname val@sub-LOIC-UDP-HOIC-RF-all\endcsname{0,0\%}
\expandafter\def\csname val@sub-LOIC-UDP-HOIC-RF-no_fp\endcsname{0,0\%}
\expandafter\def\csname val@sub-LOIC-UDP-HOIC-SVM-all\endcsname{17,6\%}
\expandafter\def\csname val@sub-LOIC-UDP-HOIC-SVM-no_fp\endcsname{17,6\%}
\expandafter\def\csname val@sub-LOIC-UDP-HOIC-XGB-all\endcsname{0,0\%}
\expandafter\def\csname val@sub-LOIC-UDP-HOIC-XGB-no_fp\endcsname{0,0\%}
\expandafter\def\csname val@sub-RF-DDOSattack-HOIC-missed\endcsname{0}
\expandafter\def\csname val@sub-RF-DDOSattack-HOIC-rate\endcsname{100,00\%}
\expandafter\def\csname val@sub-RF-DDOSattack-LOIC-UDP-missed\endcsname{0}
\expandafter\def\csname val@sub-RF-DDOSattack-LOIC-UDP-rate\endcsname{100,00\%}
\expandafter\def\csname val@sub-RF-FTP-BruteForce-missed\endcsname{0}
\expandafter\def\csname val@sub-RF-FTP-BruteForce-rate\endcsname{100,00\%}
\expandafter\def\csname val@sub-RF-SSH-Bruteforce-missed\endcsname{1}
\expandafter\def\csname val@sub-RF-SSH-Bruteforce-rate\endcsname{100,00\%}
\expandafter\def\csname val@sub-SSH-Bruteforce-n\endcsname{28\,210}
\expandafter\def\csname val@sub-SSH-FTP-RF-all\endcsname{100,0\%}
\expandafter\def\csname val@sub-SSH-FTP-RF-no_fp\endcsname{11,3\%}
\expandafter\def\csname val@sub-SSH-FTP-SVM-all\endcsname{1,9\%}
\expandafter\def\csname val@sub-SSH-FTP-SVM-no_fp\endcsname{1,9\%}
\expandafter\def\csname val@sub-SSH-FTP-XGB-all\endcsname{32,1\%}
\expandafter\def\csname val@sub-SSH-FTP-XGB-no_fp\endcsname{0,0\%}
\expandafter\def\csname val@sub-SVM-DDOSattack-HOIC-missed\endcsname{4}
\expandafter\def\csname val@sub-SVM-DDOSattack-HOIC-rate\endcsname{99,99\%}
\expandafter\def\csname val@sub-SVM-DDOSattack-LOIC-UDP-missed\endcsname{8}
\expandafter\def\csname val@sub-SVM-DDOSattack-LOIC-UDP-rate\endcsname{96,75\%}
\expandafter\def\csname val@sub-SVM-FTP-BruteForce-missed\endcsname{1}
\expandafter\def\csname val@sub-SVM-FTP-BruteForce-rate\endcsname{94,44\%}
\expandafter\def\csname val@sub-SVM-SSH-Bruteforce-missed\endcsname{6}
\expandafter\def\csname val@sub-SVM-SSH-Bruteforce-rate\endcsname{99,98\%}
\expandafter\def\csname val@sub-XGB-DDOSattack-HOIC-missed\endcsname{0}
\expandafter\def\csname val@sub-XGB-DDOSattack-HOIC-rate\endcsname{100,00\%}
\expandafter\def\csname val@sub-XGB-DDOSattack-LOIC-UDP-missed\endcsname{0}
\expandafter\def\csname val@sub-XGB-DDOSattack-LOIC-UDP-rate\endcsname{100,00\%}
\expandafter\def\csname val@sub-XGB-FTP-BruteForce-missed\endcsname{0}
\expandafter\def\csname val@sub-XGB-FTP-BruteForce-rate\endcsname{100,00\%}
\expandafter\def\csname val@sub-XGB-SSH-Bruteforce-missed\endcsname{1}
\expandafter\def\csname val@sub-XGB-SSH-Bruteforce-rate\endcsname{100,00\%}
\expandafter\def\csname val@task-conf-Bot\endcsname{496}
\expandafter\def\csname val@task-conf-BruteForce\endcsname{0}
\expandafter\def\csname val@task-conf-DDoS\endcsname{2}
\expandafter\def\csname val@task-confvec-Bot\endcsname{248}
\expandafter\def\csname val@task-confvec-BruteForce\endcsname{0}
\expandafter\def\csname val@task-confvec-DDoS\endcsname{1}
\expandafter\def\csname val@task-dup-Bot\endcsname{0}
\expandafter\def\csname val@task-dup-BruteForce\endcsname{7}
\expandafter\def\csname val@task-dup-DDoS\endcsname{0}
\expandafter\def\csname val@test-rows-Bot\endcsname{60\,000}
\expandafter\def\csname val@test-rows-BruteForce\endcsname{56\,457}
\expandafter\def\csname val@test-rows-DDoS\endcsname{60\,000}
\expandafter\def\csname val@train-rows-Bot\endcsname{140\,000}
\expandafter\def\csname val@train-rows-BruteForce\endcsname{131\,731}
\expandafter\def\csname val@train-rows-DDoS\endcsname{140\,000}

\graphicspath{{/home/user/ml-ids-cicids2018/results/figures/}}


\begin{document}
% ============================== ЗАГЛАВНА СТРАНИЦА ==============================
\begin{titlepage}
\newgeometry{top=2.2cm,bottom=2.2cm,left=2.6cm,right=2.6cm}
\pagecolor{white}
\begin{tikzpicture}[remember picture,overlay]
  \fill[navy] (current page.north west) rectangle ([yshift=-7.2cm]current page.north east);
  \fill[accent] ([yshift=-7.2cm]current page.north west) rectangle ([yshift=-7.45cm]current page.north east);
\end{tikzpicture}
\begin{center}
{\sffamily\color{white}
{\large\bfseries [Университет] \par}\vspace{2pt}
{\normalsize [Факултет] \textbullet\ Магистърска програма [наименование]\par}\vspace{14pt}
{\small Дисциплина\par}\vspace{2pt}
{\large Интелигентни системи за сигурност на екосистемите\\на големите данни и Интернет на нещата\par}\vspace{4pt}
{\small учебна 2026/27 г.\par}}
\vspace{3.4cm}

{\sffamily\small\color{accent}\bfseries КУРСОВ ПРОЕКТ \textbullet\ ЕКСПЕРИМЕНТАЛНО ИЗСЛЕДВАНЕ\par}\vspace{10pt}
{\sffamily\fontsize{22}{27}\selectfont\bfseries\color{navy}
Откриване на проникване в мрежов\\ трафик на център за данни чрез\\ машинно обучение\par}
\vspace{10pt}
{\large\color{muted} Random Forest, SVM и XGBoost върху реалния набор CSE-CIC-IDS2018,\\ с редукция на размерността и статистическа оценка\par}
\vspace{1.4cm}

\begin{tabular}{@{}p{4.2cm}p{4.2cm}p{3.6cm}@{}}
\toprule
\textsf{\textbf{Екип №}} & \textsf{\textbf{Студент}} & \textsf{\textbf{Факултетен номер}}\\
\midrule
{[\ \ ]} & Иван Терзиев & 961326004\\
    & Кристиян Януцев & 961326003\\
    & Тихомир Божков & 961326014\\
\bottomrule
\end{tabular}

\vfill
{\sffamily\small\color{muted} Преподавател: [ ] \hfill София, \today\par}
\end{center}
\restoregeometry
\end{titlepage}

\pagenumbering{roman}
% ============================== АБСТРАКТ ==============================
\chapter*{Абстракт}
\addcontentsline{toc}{chapter}{Абстракт}
\thispagestyle{plain}

Центровете за данни и облачните инфраструктури концентрират услугите на хиляди организации. Затова те са основна цел на атаките Denial of Service (DoS), на опитите за неоторизиран достъп с Brute Force и на мрежите Botnet. Настоящото изследване изгражда, обучава и експериментално сравнява девет модела за откриване на прониквания (IDS) в мрежовия трафик на център за данни. Използвани са три алгоритъма от различни семейства: \emph{Random Forest}, \emph{метод на опорните вектори с RBF ядро} (SVM) и \emph{градиентно усилване на решаващи дървета} (XGBoost). Те са приложени към три атаки с различен характер: Distributed Denial of Service (DDoS), Brute Force срещу FTP и SSH (BruteForce) и Botnet (Bot).

Изследването използва изключително реални записи от публичния набор CSE-CIC-IDS2018. Наборът е уловен в инфраструктура на Amazon Web Services с 420 машини и 30 сървъра. Изтеглени са \val{rows-read} записа на мрежови потоци с \val{n-features} признака. След почистване (премахване на \val{nan-rows} невалидни записа и \val{dup-rows} дубликата) остават \val{rows-clean}. Векторът на признаците се редуцира в два етапа: филтър по дисперсия (\val{n-features}\,$\to$\,\val{n-filtered}) и подбор на \val{n-selected}-те най-информативни признака по важност в помощен модел, изчислена само върху обучаващите данни. Анализът на главните компоненти (PCA) е контролен вариант. Хиперпараметрите се избират с 5-кратна кръстосана валидация и правилото „една стандартна грешка“, което предпочита по-простия модел и ограничава преобучението. Моделите се оценяват в отделна оперативна фаза върху 30\% невиждани данни по TP, TN, FP, FN, Accuracy, Precision, Recall (Sensitivity), F1 и Specificity.

Всички девет модела постигат Accuracy над 99,5\%. За DDoS тя е \val{RF-DDoS-Accuracy} (RF), \val{SVM-DDoS-Accuracy} (SVM) и \val{XGB-DDoS-Accuracy} (XGBoost). За BruteForce е \val{RF-BruteForce-Accuracy}, \val{SVM-BruteForce-Accuracy} и \val{XGB-BruteForce-Accuracy}, а за Bot е \val{RF-Bot-Accuracy}, \val{SVM-Bot-Accuracy} и \val{XGB-Bot-Accuracy}. Разликата между точността при обучение и при тест е най-много \val{gap-max} процентни пункта, а оценката от кръстосаната валидация съвпада с тестовата, т.е. моделите не са преобучени. Ансамблите от дървета превъзхождат SVM статистически значимо при Bot (тест на McNemar). При естественото съотношение на класовете (\val{nat-share-Bot}--\val{nat-share-DDoS} атаки) Precision остава над 99\%. Контролните експерименти показват, че резултатите не зависят от времевите признаци и от метода на редукция. Допълнителни проверки за приложимост показват, че точността почти не зависи от признаците-отпечатъци на атакуващия инструмент и се запазва при хронологично разделяне за DDoS и BruteForce. При Bot ансамблите пропускат около 10\% от по-късните атаки, а комбинираният детектор XGBoost или SVM ги открива почти всички. Предварителното изследване върху IoT набора CICIoT2023 показва обратното за Port Scan (Recon-PortScan): таванът там е около 85\%. Така качеството на една IDS се определя преди всичко от това дали признаците на потока носят отпечатъка на атаката.

\medskip
\noindent\textsf{\textbf{\color{navy}Ключови думи:}} откриване на прониквания, център за данни, облак, CSE-CIC-IDS2018, DDoS, Brute Force, Botnet, Random Forest, SVM, XGBoost, подбор на признаци, преобучение, bootstrap, тест на McNemar.


\tableofcontents
% \listoffigures
% \listoftables
\clearpage
\pagenumbering{arabic}

% ============================== УВОД ==============================
\chapter*{Увод}
\addcontentsline{toc}{chapter}{Увод}
\markboth{Увод}{}

\section*{Контекст}
Центровете за данни и публичните облаци обслужват уеб приложения, бази от данни, системи за обработка на големи данни и услугите за Интернет на нещата. Концентрацията на ресурси ги прави привлекателна цел. Атаките Distributed Denial of Service (DDoS) достигат терабитови мощности. Botnet Mirai през 2016~г. използва стотици хиляди заразени устройства за атаки срещу доставчици на DNS и хостинг услуги~\cite{antonakakis2017}. Едновременно с това сканирането на портове и автоматизираните опити за отгатване на пароли срещу публично достъпни услуги като SSH и FTP са постоянен фонов шум във всяка облачна инфраструктура. Сканирането е първата стъпка на почти всяка атака: с него нападателят открива кои услуги са достъпни~\cite{lyon2009}.

Защитата на центъра за данни не може да разчита само на агенти върху крайните точки. Много системи са виртуални машини или контейнери с кратък живот, а част от трафика минава между вътрешни услуги. Затова основна линия на защита е \emph{мрежовата система за откриване на прониквания} (Network Intrusion Detection System, NIDS), която наблюдава потоците и сигнализира при подозрително поведение. Сигнатурните NIDS, като Snort и Suricata, откриват само вече описани атаки. Подходите с машинно обучение (ML) се учат от примери и могат да разпознават закономерности, които не са описани ръчно~\cite{sommer2010,buczak2016}. Те обаче имат собствени рискове: зависят от качеството на данните за обучение, могат да постигат подвеждащо високи метрики при нереалистична оценка и могат да се преобучат.

\section*{Мотивация и изследователски въпроси}
Литературата за ML-базирани IDS често съобщава точност над 99\%. Такива резултати обаче рядко отчитат три фактора: (1) дали признаците на отделния запис изобщо носят отпечатъка на атаката, (2) реалната честота на атаките в трафика и (3) статистическата несигурност и стабилността на самите метрики. Настоящото изследване поставя тези фактори в центъра на експеримента и търси отговор на следните въпроси:

\begin{enumerate}[label=\textbf{В\arabic*.},leftmargin=2.4em]
  \item Колко точно три принципно различни алгоритъма (ансамбъл с bagging, метод с ядро и ансамбъл с boosting) откриват три атаки с различен характер в трафика на център за данни?
  \item Как редукцията на вектора на признаците (подбор на признаци или PCA) влияе на качеството и на скоростта на моделите?
  \item Статистически значими ли са разликите между алгоритмите, колко стабилни са метриките и има ли преобучение?
  \item Какво е реалното качество на моделите при естественото съотношение между атаки и нормален трафик и как може да се избере прагът на решение?
  \item Кои признаци на потока носят информацията за всяка атака и от какво зависи постижимата точност?
\end{enumerate}

\section*{Принос}
Работата има следния принос:
\begin{itemize}
  \item напълно възпроизводим експериментален конвейер с отворен код, който използва \emph{само реални} записи от CSE-CIC-IDS2018 и CIC-IDS2017 и предпазва от изтичане на информация между обучение и тест;
  \item обосновка на избора на набори от данни: предварително изследване върху CICIoT2023 показва количествено защо Port Scan там не може да се открие с точност над около 85\%, а същата атака в потоковото представяне на CIC-IDS2017 се открива с над 99\%;
  \item декларативна архитектура (регистри на атаки и алгоритми), при която нова атака, нов алгоритъм или атака от друг набор се добавят с един запис;
  \item девет модела, оценени по TP, TN, FP, FN, Accuracy, Precision, Recall, F1 и Specificity с тълкуване „на 100 случая“, 95\% доверителни интервали, тест на McNemar за всяка двойка алгоритми и явна диагностика на преобучението;
  \item пет контролни експеримента (E1--E5) и девет разширени анализа (A1--A9), сред които криви на обучение, сравнение на подбор на признаци и PCA, избор на праг при естествено съотношение и пренос между атаките;
  \item пет проверки за приложимост (A10--A14): без признаците-отпечатъци на инструмента, хронологично разделяне, пренос между подвидовете на атаката, нормален трафик от непознат ден и комбиниран детектор XGBoost или SVM.
\end{itemize}

\section*{Структура на доклада}
Структурата следва зададения план. \Cref{ch:method} описва целта и методологията: избора на атаки и набор от данни, предварителната обработка, математическите основи на алгоритмите и метриките за оценка. \Cref{ch:framework} представя експерименталната технологична рамка, а \cref{ch:plan} описва плана на експериментите с фаза на машинно обучение и оперативна фаза. \Cref{ch:results} съдържа имплементацията и резултатите, а \cref{ch:analysis} ги анализира. Следват заключение и информационни източници (\cref{ch:refs}). Приложенията съдържат речник на признаците, пълните резултати от настройката, таблиците на анализите A3 и A4, фрагменти от кода, допълнителни фигури и инструкции за възпроизвеждане.

% ============================== 1. ЦЕЛ И МЕТОДОЛОГИЯ ==============================
\chapter{Цел на изследването и методология за постигане на целта}\label{ch:method}

\begin{keybox}[Цел на изследването]
Да се синтезират, обучат и експериментално сравнят девет модела за откриване на прониквания в мрежовия трафик на център за данни: три атаки (DDoS, BruteForce и Bot), всяка с три алгоритъма (Random Forest, SVM и XGBoost). Моделите се оценяват по TP, TN, FP, FN, Accuracy, Precision и Sensitivity (Recall), допълнени с F1 и Specificity. Методите за обработка и редукция на данните и рискът от преобучение се анализират количествено.
\end{keybox}

Методологията следва цикъла на ML-базирана IDS: избор на атаки и данни, предварителна обработка с редукция на размерността, синтез и настройка на моделите, тестване върху невиждани данни (оперативна фаза) и статистическа оценка~\cite{buczak2016,arp2022}. Всички случайни процеси използват фиксирано начало на генератора (seed~42), така че експериментът се възпроизвежда изцяло.

% ---------------------------------------------------------------------------
\section{Избор на атаки и набори от данни}\label{sec:attacks}

Трите атаки покриват три различни етапа и цели на нападателя в един център за данни.

\emph{Distributed Denial of Service (DDoS)} е шумна, обемна атака: множество атакуващи машини едновременно изпращат трафик, който изчерпва пропускателната способност, таблицата на връзките или ресурсите на приложението~\cite{mirkovic2004}. В използвания набор DDoS е проведена с инструментите HOIC (наводнение с HTTP заявки) и LOIC-UDP (наводнение с UDP пакети). Двете разновидности образуват един клас „DDoS“, защото практическата IDS трябва да разпознава семейството атаки, а не отделния инструмент.

\emph{Brute Force (BruteForce)}, атака по речник, е опит за неоторизиран достъп: нападателят автоматично пробва хиляди двойки потребител--парола срещу услуга за вход. В набора тя е проведена с инструмента Patator срещу FTP и SSH сървъри~\cite{sharafaldin2018}. Всеки опит е кратка, правилно оформена сесия, но опитите се повтарят с еднаква структура.

\emph{Botnet (Bot)} е етапът след успешното проникване. Заразената машина поддържа връзка с команден сървър, получава команди, изпраща снимки на екрана и откраднати данни. В набора е използван Botnet Ares. Трафикът му е с малък обем и периодичен, затова се различава от DDoS по характер.

Изискването е работа \emph{само с реални и достоверни данни}. Затова са сравнени публичните набори с реално уловен трафик (\cref{tab:datasets}).

\begin{table}[H]
\centering\small
\caption{Сравнение на публичните набори от данни за IDS}\label{tab:datasets}
\begin{tabularx}{\textwidth}{@{}l c L c@{}}
\toprule
\textsf{\textbf{Набор}} & \textsf{\textbf{Година}} & \textsf{\textbf{Среда и бележки}} & \textsf{\textbf{Център за данни}}\\
\midrule
KDD Cup 99 / NSL-KDD & 1999/2009 & военна мрежа от 90-те; остарели атаки, много дубликати~\cite{tavallaee2009} & не\\
CIC-IDS2017 & 2017 & корпоративна мрежа, 5 дни реален трафик, достъп след регистрация & частично\\
\rowcolor{boxbg}CSE-CIC-IDS2018 & 2018 & \textbf{облачна инфраструктура в AWS}: 420 машини, 30 сървъра, 10 дни, 14 вида атаки; публичен достъп~\cite{sharafaldin2018,cse2018web} & \textbf{да}\\
CICIoT2023 & 2023 & 105 реални IoT устройства, 33 атаки~\cite{neto2023}; изследван предварително & не\\
\bottomrule
\end{tabularx}
\end{table}

\paragraph{Предварително изследване върху CICIoT2023.} Първоначално изследването беше проведено върху IoT набора CICIoT2023~\cite{neto2023,cic2023web} с атаките DDoS и Port Scan (Recon-PortScan)~\cite{lyon2009}. DDoS се откриваше с Accuracy над 99,9\%, но Port Scan остана около 84\% дори след регуляризация, подбор на признаци, нови признаци и почистване на шума в обучаващите данни (най-добре около 85\%). Причината е в данните: всеки запис описва прозорец от пакети в рамките на един поток, а отделната проба на скенера изглежда като нормална кратка TCP сесия. Затова изследването е пренесено върху CSE-CIC-IDS2018, който е записан в облачен център за данни и е по-подходящ и по тема.

\paragraph{Избраният набор CSE-CIC-IDS2018.} Наборът е създаден съвместно от Communications Security Establishment (CSE) и Canadian Institute for Cybersecurity (University of New Brunswick)~\cite{sharafaldin2018,cse2018web}. Атакуваната организация е изградена в Amazon Web Services (AWS) с 5 отдела, 420 машини и 30 сървъра, а атакуващата инфраструктура е от 50 машини. Записани са 10 работни дни трафик (14.02--02.03.2018). От записите (\texttt{pcap}) инструментът CICFlowMeter-V3 извлича по 80 признака на всеки двупосочен мрежов поток: продължителност, брой и размер на пакетите в двете посоки, интервали между пакетите, TCP флагове, размер на TCP прозореца, периоди на активност и бездействие. Наборът е публикуван в каталога Registry of Open Data on AWS и се изтегля без регистрация~\cite{aws2018}. Той е сред най-използваните в изследванията на IDS върху големи данни~\cite{leevy2020}.

Използвани са три дневни файла (\cref{tab:class-counts}): 21.02.2018 (DDoS), 14.02.2018 (BruteForce) и 02.03.2018 (Bot), общо \val{rows-read} записа. Нормалният трафик от трите дни се обединява. Етикетът на всеки запис е в колоната \texttt{Label}. Две колони не се използват като признаци. \texttt{Timestamp} е моментът на улавяне, а не свойство на трафика. \texttt{Dst Port} е номерът на порта на жертвата. Той би позволил на модела да „запомни“ порта на атакуваната услуга (21 за FTP, 22 за SSH), вместо да научи поведението на потока. Остават \val{n-features} числови признака. Синтетични, генерирани или аугментирани записи не се използват никъде, включително в автоматичните тестове на кода.

\begin{table}[H]
\centering\small
\caption{Изтеглени записи по етикет преди и след почистване}\label{tab:class-counts}
\fitbox{%
\begin{tabular}{@{}lllrrr@{}}
\toprule
\headrow \hd{Набор} & \hd{Етикет} & \hd{Задача} & \hd{Изтеглени} & \hd{След почистване} & \hd{Премахнати}\\
\midrule
CSE-CIC-IDS2018 & \texttt{Benign} & нормален (DDoS, BruteForce) & 1\,790\,843 & 1\,472\,334 & 17,8\%\\
\rowcolor{rowalt}CSE-CIC-IDS2018 & \texttt{DDOS attack-HOIC} & DDoS & 686\,012 & 198\,861 & 71,0\%\\
CSE-CIC-IDS2018 & \texttt{DDOS attack-LOIC-UDP} & DDoS & 1\,730 & 1\,730 & 0,0\%\\
\rowcolor{rowalt}CSE-CIC-IDS2018 & \texttt{FTP-BruteForce} & BruteForce & 193\,360 & 53 & 100,0\%\\
CSE-CIC-IDS2018 & \texttt{SSH-Bruteforce} & BruteForce & 187\,589 & 94\,048 & 49,9\%\\
\rowcolor{rowalt}CIC-IDS2017 & \texttt{BENIGN} & нормален (PortScan) & 979\,288 & 893\,304 & 8,8\%\\
CIC-IDS2017 & \texttt{Portscan} & PortScan & 159\,066 & 1\,683 & 98,9\%\\
\rowcolor{rowalt}двата набора & други етикети & не се използват & 563\,659 & -- & --\\
\textbf{Общо} &  &  & \textbf{4\,561\,547} & \textbf{2\,662\,013} & \\
\bottomrule
\end{tabular}}%

\end{table}
\interprettab{Всеки ред е един етикет от оригиналните файлове. Най-много записи губят атаките BruteForce: \texttt{FTP-BruteForce} съдържа почти само еднакви потоци (неуспешни опити за вход с еднаква структура), затова след премахването на дубликатите от него остават само няколко десетки уникални записа. Нормалният трафик също съдържа много повтарящи се кратки потоци. Дубликатите се премахват, за да не попадне един и същ запис едновременно в обучаващите и в тестовите данни.}

Признаците образуват седем семейства: обем (брой и размер на пакетите в двете посоки, подпотоци), дължини на пакетите (минимум, максимум, средна, стандартно отклонение), интензивност (байтове и пакети в секунда), време (продължителност на потока, интервали между пакетите IAT, периоди на активност и бездействие), TCP флагове (броячи на FIN, SYN, RST, PSH, ACK, URG, ECE), заглавка (протокол, дължина на заглавките, начален размер на TCP прозореца) и насрещни блокове (bulk)~\cite{sharafaldin2018}. Пълният речник на признаците е в хранилището на проекта (\texttt{results/feature\_dictionary.md})~\cite{repo}.

% ---------------------------------------------------------------------------
\section{Предварителна обработка на данните}\label{sec:preprocessing}

Нека суровите данни са $\mathcal{D}=\{(\vect{x}_i, \ell_i)\}_{i=1}^{N}$, където $\vect{x}_i\in\R^{d}$ е векторът на признаците ($d=\val{n-features}$), а $\ell_i$ е оригиналният етикет. Обработката е последователност от трансформации, които превръщат $\mathcal{D}$ в обучаващи и тестови данни за всяка бинарна задача~\cite{arp2022}.

\paragraph{Почистване (премахване на шума).} Премахват се записите с липсващи или безкрайни стойности и пълните дубликати:
\begin{equation}\label{eq:clean}
\mathcal{D}' = \operatorname{unique}\bigl\{(\vect{x}_i,\ell_i)\in\mathcal{D} \;:\; x_{ij}\in\R \ \ \forall j=1,\dots,d\bigr\}.
\end{equation}
Дубликатите се премахват \emph{преди} разделянето на обучаващи и тестови данни, иначе един и същ запис може да попадне в двете части и да надцени качеството~\cite{engelen2021,liu2022}.

\paragraph{Бинарни задачи и почистване на ниво задача.} За атака $a$ с множество от оригинални етикети $\mathcal{L}_a$ (напр. \texttt{FTP-BruteForce} и \texttt{SSH-Bruteforce} за BruteForce) бинарният етикет е
\begin{equation}\label{eq:binary}
y_i = \mathbf{1}\left[\ell_i\in\mathcal{L}_a\right],\qquad i\in\mathcal{I}_a=\{i:\ \ell_i\in\mathcal{L}_a\ \lor\ \ell_i=\txt{Benign}\}.
\end{equation}
Атаката е положителният клас ($y=1$), а нормалният трафик е отрицателният ($y=0$). Някои вектори на признаците се срещат и с етикет на атака, и с етикет на нормален трафик. Това е неразрешимо противоречие (шум в етикетите), известно за наборите на CIC~\cite{liu2022}. Затова в рамките на всяка задача се запазва по едно копие от всеки вектор, а всички вектори с противоречиви етикети се премахват:
\begin{equation}\label{eq:conflict}
\mathcal{C}_a=\bigl\{\vect{x}:\ \exists\,i,k\in\mathcal{I}_a,\ \vect{x}_i=\vect{x}_k=\vect{x},\ y_i\neq y_k\bigr\},\qquad \mathcal{I}_a'=\operatorname{unique}\bigl\{i\in\mathcal{I}_a:\ \vect{x}_i\notin\mathcal{C}_a\bigr\}.
\end{equation}
Така нито един вектор не може да се появи едновременно в обучаващите и в тестовите данни.

\paragraph{Балансиране и разделяне.} От всеки клас се избират случайно без връщане $n_a=\min(|\mathcal{I}_a^{+}|,|\mathcal{I}_a^{-}|,\,10^5)$ записа, така че задачата е балансирана 1:1. Това предпазва моделите от тривиалното решение „всичко е нормално“ и прави Accuracy пригодна метрика, защото при симетрични данни FP и FN имат сравнима тежест~\cite{borovska2022}. Реалистичното небалансирано съотношение се изследва отделно (експеримент E2). Извадката се разделя стратифицирано:
\begin{equation}\label{eq:split}
\mathcal{S}_a=\mathcal{T}_a\,\dot\cup\,\mathcal{V}_a,\qquad |\mathcal{T}_a| = 0{,}7\,|\mathcal{S}_a|,\qquad \tfrac{1}{|\mathcal{T}_a|}\textstyle\sum_{i\in\mathcal{T}_a}y_i=\tfrac{1}{|\mathcal{V}_a|}\sum_{i\in\mathcal{V}_a}y_i=\tfrac12,
\end{equation}
където $\mathcal{T}_a$ е обучаващото множество (фаза на машинно обучение), а $\mathcal{V}_a$ е тестовото (оперативна фаза).

\paragraph{Редукция, етап 1: филтър по дисперсия.} Признаците с дисперсия $s_j^2\le 10^{-8}$ върху обучаващите данни се премахват като константни, където
\begin{equation}\label{eq:var}
s_j^2=\frac{1}{n-1}\sum_{i=1}^{n}\left(x_{ij}-\bar{x}_j\right)^2,\qquad \bar{x}_j=\frac{1}{n}\sum_{i=1}^{n}x_{ij}.
\end{equation}
В CSE-CIC-IDS2018 това са \val{lowvar-DDoS} признака (броячите на насрещни блокове и някои флагове), които CICFlowMeter-V3 записва винаги с нула. Остават $d_f=\val{n-filtered}$ признака.

\paragraph{Редукция, етап 2: подбор на признаци по важност.} Вторият етап избира $k=\val{n-selected}$ оригинални признака, които носят най-много информация за задачата. Върху обучаващите данни се обучава помощен ансамбъл от градиентно усилени дървета (\cref{eq:xgb-obj}), а важността на признак $j$ е средното подобрение на целевата функция (gain) при разклоненията по него:
\begin{equation}\label{eq:gain-imp}
I_j=\frac{1}{|\mathcal{R}_j|}\sum_{r\in\mathcal{R}_j}\operatorname{Gain}(r),\qquad \tilde I_j=\frac{I_j}{\sum_{m}I_m},\qquad \mathcal{J}^{*}=\operatorname{top}_k\bigl\{\tilde I_j\bigr\},
\end{equation}
където $\mathcal{R}_j$ е множеството от разклонения по признак $j$, а $\operatorname{Gain}$ е дефиниран в \cref{eq:xgb-gain}. Подборът запазва оригиналните, интерпретируеми признаци. Той отчита връзката с етикета, докато PCA търси само направленията с голяма дисперсия, които не са непременно полезни за класификацията~\cite{guyon2003}. В предварителното изследване подборът подобри точността при Port Scan с около 3 пункта спрямо PCA. Затова той е основният метод, а PCA е контролен вариант (експеримент E5).

\paragraph{Стандартизация.} Признаците имат много различни мащаби (микросекунди, байтове, броячи), а SVM и PCA са чувствителни към мащаба. Затова всеки признак се стандартизира с параметрите от обучаващите данни:
\begin{equation}\label{eq:zscore}
z_{ij}=\frac{x_{ij}-\bar{x}_j}{s_j}.
\end{equation}

\paragraph{Контролен метод: анализ на главните компоненти (PCA).} PCA търси ортогонални направления с максимална дисперсия~\cite{jolliffe2016}. За стандартизираната матрица $\vect{Z}\in\R^{n\times d_f}$ ковариационната матрица и нейното спектрално разлагане са
\begin{equation}\label{eq:cov}
\bm{\Sigma}=\frac{1}{n-1}\vect{Z}^{\top}\vect{Z},\qquad \bm{\Sigma}\,\vect{w}_k=\lambda_k\vect{w}_k,\qquad \lambda_1\ge\lambda_2\ge\dots\ge\lambda_{d_f}\ge 0.
\end{equation}
Запазва се най-малкият брой компоненти, който обяснява поне 95\% от дисперсията, и всеки запис се проектира върху тях:
\begin{equation}\label{eq:kstar}
\eta(k)=\frac{\sum_{m=1}^{k}\lambda_m}{\sum_{m=1}^{d_f}\lambda_m},\qquad k^{*}=\min\{k:\ \eta(k)\ge 0{,}95\},\qquad \tilde{\vect{x}}=\left[\vect{w}_1,\dots,\vect{w}_{k^*}\right]^{\top}\vect{z}.
\end{equation}

\begin{notebox}[Предпазване от изтичане на информация]
Филтърът, подборът на признаци и стандартизацията са стъпки на един \texttt{sklearn.Pipeline} заедно с класификатора~\cite{pedregosa2011}. Параметрите им, включително важността на признаците, се оценяват \emph{само} върху обучаващите данни, и то вътре във всеки фолд на кръстосаната валидация. Тестовите данни се трансформират с вече научените параметри.
\end{notebox}

% ---------------------------------------------------------------------------
\section{Избор и дизайн на модели за машинно обучение}\label{sec:models}

Избрани са три алгоритъма от различни семейства. Random Forest усреднява много независими дървета (bagging). XGBoost строи дърветата последователно, като всяко поправя грешките на предишните (boosting). SVM с RBF ядро е метод на максималния марж, който строи гладка нелинейна граница. Различните индуктивни пристрастия правят сравнението информативно~\cite{hastie2009}.

\paragraph{Решаващо дърво и Random Forest.} Решаващото дърво (CART) разделя пространството рекурсивно с прагови правила $x_j\le t$. Във всеки възел се избира разделянето, което най-много намалява индекса на Джини~\cite{breiman1984}:
\begin{equation}\label{eq:gini}
G(S_v)=1-\sum_{c\in\{0,1\}}p_{v,c}^2,\qquad p_{v,c}=\frac{|\{i\in S_v:\ y_i=c\}|}{|S_v|}.
\end{equation}
Единичното дърво има висока дисперсия. Random Forest~\cite{breiman2001} обучава $B$ дървета, всяко върху bootstrap извадка от обучаващите данни, като при всяко разделяне разглежда само случайно подмножество от $m$ признака. Вероятността за атака и решението на ансамбъла са
\begin{equation}\label{eq:rf}
\hat{p}(\vect{x})=\frac{1}{B}\sum_{b=1}^{B}\hat{p}_b(\vect{x}),\qquad \hat{y}(\vect{x})=\mathbf{1}\left[\hat{p}(\vect{x})>0{,}5\right],
\end{equation}
където $\hat{p}_b(\vect{x})$ е делът на атаките в листото на дърво $b$. Ако дърветата имат дисперсия $\sigma^2$ и средна двойкова корелация $\rho$, дисперсията на средното им е~\cite{hastie2009}
\begin{equation}\label{eq:rf-var}
\operatorname{Var}\bigl[\hat{p}\bigr]=\rho\,\sigma^2+\frac{1-\rho}{B}\,\sigma^2,
\end{equation}
т.е. ансамбълът е толкова по-стабилен, колкото по-слабо корелирани са дърветата. Минималният брой записи в листо и максималната дълбочина ограничават сложността на всяко дърво и така намаляват преобучението.

\paragraph{Градиентно усилване (XGBoost).} XGBoost~\cite{chen2016} е адитивен модел от $K$ регресионни дървета $f_k$, построени последователно~\cite{friedman2001}. Оценката (логит) и вероятността за атака са
\begin{equation}\label{eq:xgb-model}
F(\vect{x})=\sum_{k=1}^{K}\nu\,f_k(\vect{x}),\qquad \hat{p}(\vect{x})=\sigma\bigl(F(\vect{x})\bigr)=\frac{1}{1+e^{-F(\vect{x})}},
\end{equation}
където $\nu$ е скоростта на обучение. Всяко ново дърво минимизира регуляризираната целева функция
\begin{equation}\label{eq:xgb-obj}
\mathcal{L}^{(k)}=\sum_{i=1}^{n}l\bigl(y_i,\ F^{(k-1)}(\vect{x}_i)+f_k(\vect{x}_i)\bigr)+\Omega(f_k),\qquad \Omega(f)=\gamma T+\tfrac12\lambda\sum_{t=1}^{T}w_t^2,
\end{equation}
където $l$ е логистичната загуба, $T$ е броят листа, а $w_t$ са стойностите в листата. Членът $\Omega$ наказва сложните дървета и е вградената защита срещу преобучение. С приближение от втори ред с градиентите $g_i=\partial l/\partial F$ и вторите производни $h_i=\partial^2 l/\partial F^2$ оптималната стойност на листо и подобрението от разклонение на множеството $S$ на $S_L$ и $S_R$ са
\begin{equation}\label{eq:xgb-gain}
w_t^{*}=-\frac{\sum_{i\in S_t}g_i}{\sum_{i\in S_t}h_i+\lambda},\qquad
\operatorname{Gain}=\frac12\left[\frac{G_L^2}{H_L+\lambda}+\frac{G_R^2}{H_R+\lambda}-\frac{(G_L+G_R)^2}{H_L+H_R+\lambda}\right]-\gamma,
\end{equation}
където $G$ и $H$ са сумите на $g_i$ и $h_i$ в съответния възел. Допълнително всяко дърво се обучава върху случайни 80\% от записите и 80\% от признаците (subsample, colsample), а параметърът \texttt{min\_child\_weight} забранява листа с малка сума на $h_i$.

\paragraph{Метод на опорните вектори (SVM).} SVM търси хиперравнината $\vect{w}^{\top}\phi(\vect{x})+b=0$, която разделя класовете с максимален марж~\cite{cortes1995}. С етикети $t_i=2y_i-1\in\{-1,+1\}$ задачата с мек марж е
\begin{equation}\label{eq:svm-primal}
\min_{\vect{w},b,\bm{\xi}}\ \frac12\|\vect{w}\|^2+C\sum_{i=1}^{n}\xi_i\qquad\txt{при}\qquad t_i\left(\vect{w}^{\top}\phi(\vect{x}_i)+b\right)\ge 1-\xi_i,\quad \xi_i\ge 0.
\end{equation}
Параметърът $C$ регулира компромиса между широк марж и малко грешки. Дуалната задача съдържа данните само чрез скаларни произведения, които се заменят с ядро~\cite{scholkopf2002}:
\begin{equation}\label{eq:svm-dual}
\max_{\bm{\alpha}}\ \sum_{i=1}^{n}\alpha_i-\frac12\sum_{i=1}^{n}\sum_{k=1}^{n}\alpha_i\alpha_k t_i t_k K(\vect{x}_i,\vect{x}_k)\qquad\txt{при}\qquad 0\le\alpha_i\le C,\ \ \sum_{i}\alpha_i t_i=0.
\end{equation}
Използва се радиалното базисно ядро, а решаващата функция зависи само от опорните вектори ($\alpha_i>0$):
\begin{equation}\label{eq:rbf}
K(\vect{x},\vect{x}')=\exp\left(-\gamma\,\|\vect{x}-\vect{x}'\|^2\right),\qquad f(\vect{x})=\sum_{i:\ \alpha_i>0}\alpha_i t_i K(\vect{x}_i,\vect{x})+b.
\end{equation}
Малко $\gamma$ дава гладка граница, а голямо $\gamma$ дава локална граница с риск от преобучение. Обучението (SMO в LIBSVM~\cite{chang2011}) има сложност между $O(n^2)$ и $O(n^3)$, а класификацията на един запис е $O(n_{SV}\cdot d)$. Затова SVM се обучава върху ограничена извадка до \num{50000} записа.

\paragraph{Избор на хиперпараметри и защита от преобучение.} Хиперпараметрите $\bm{\theta}$ се избират с обхождане на решетка и $K=5$-кратна стратифицирана кръстосана валидация~\cite{kohavi1995}. За всяка конфигурация се изчисляват средната оценка и стандартната грешка по фолдовете:
\begin{equation}\label{eq:cv}
\overline{\operatorname{F1}}(\bm{\theta})=\frac{1}{K}\sum_{k=1}^{K}\operatorname{F1}\Bigl(h_{\bm{\theta}}^{(\mathcal{T}\setminus\mathcal{F}_k)},\ \mathcal{F}_k\Bigr),\qquad \operatorname{SE}(\bm{\theta})=\frac{s_{\operatorname{F1}}(\bm{\theta})}{\sqrt{K}},
\end{equation}
където $h_{\bm{\theta}}^{(\mathcal{A})}$ е целият конвейер, обучен върху $\mathcal{A}$, а $\mathcal{F}_k$ е $k$-тият фолд. Вместо конфигурацията с най-висока оценка се избира \emph{най-простата} конфигурация, чиято оценка е в рамките на една стандартна грешка от най-добрата (правилото „една стандартна грешка“~\cite{breiman1984,hastie2009}):
\begin{equation}\label{eq:one-se}
\bm{\theta}^{*}=\argmin_{\bm{\theta}\in\Theta}\Bigl\{c(\bm{\theta}):\ \overline{\operatorname{F1}}(\bm{\theta})\ \ge\ \overline{\operatorname{F1}}(\hat{\bm{\theta}})-\operatorname{SE}(\hat{\bm{\theta}})\Bigr\},\qquad \hat{\bm{\theta}}=\argmax_{\bm{\theta}\in\Theta}\overline{\operatorname{F1}}(\bm{\theta}),
\end{equation}
където $c(\bm{\theta})$ е мярка за сложност: по-голям минимален брой записи в листо и по-малка дълбочина за дърветата, по-малки $C$ и $\gamma$ за SVM. Конфигурациите в рамките на една стандартна грешка са статистически неразличими, а по-простата обобщава по-добре. Критерият F1 балансира пропуснатите атаки и фалшивите аларми (\cref{eq:f1}). Преобучението се проверява явно: за всеки модел се сравняват точността върху обучаващите данни, оценката от кръстосаната валидация и точността върху тестовите данни.

% ---------------------------------------------------------------------------
\section{Експериментална оценка на моделите}\label{sec:metrics}

Ефективността на моделите се оценява чрез стандартните метрики за качество на класификацията в системите за откриване на проникване: TP, TN, FP, FN, Precision, Accuracy и Sensitivity (Recall). Те показват както общата точност на модела, така и доколко той открива реалните атаки и избягва фалшивите аларми. Резултатът от класификацията на тестовите записи се обобщава в матрица на объркване (\cref{tab:cm}). Положителният клас е атаката. По аналогия с медицинската диагноза~\cite{borovska2022} TP и TN са желаните решения, FP е фалшива тревога, а FN (пропусната атака) е най-нежеланата грешка, защото атаката продължава незабелязано.

\begin{table}[H]
\centering\small
\caption{Матрица на объркване: определение на TP, TN, FP и FN}\label{tab:cm}
\begin{tabular}{@{}l l c c@{}}
\toprule
 & & \multicolumn{2}{c}{\textsf{\textbf{Предсказан клас}}}\\
\cmidrule(l){3-4}
 & & Атака ($\hat y=1$) & Нормален ($\hat y=0$)\\
\midrule
\multirow{2}{*}{\textsf{\textbf{Реален клас}}} & Атака ($y=1$) & \cellcolor{boxbg}TP: открита атака & FN: пропусната атака\\
 & Нормален ($y=0$) & FP: фалшива аларма & \cellcolor{boxbg}TN: правилно пропуснат\\
\bottomrule
\end{tabular}
\end{table}

Метриките са изброени в реда на учебния материал~\cite{borovska2022}. Всяка отговаря на различен въпрос и се тълкува „на 100 случая“.

\paragraph{Акуратност (Accuracy)} е делът на всички верни решения:
\begin{equation}\label{eq:accuracy}
\operatorname{Accuracy}=\frac{TP+TN}{TP+TN+FP+FN}.
\end{equation}
Accuracy~=~0,99 означава, че от 100 тестови записа 99 са класифицирани правилно. Метриката е най-интуитивна, но е подходяща главно за симетрични данни, при които FP и FN са близки.

\paragraph{Прецизност (Precision)} е делът на истинските атаки сред всички аларми:
\begin{equation}\label{eq:precision}
\operatorname{Precision}=\frac{TP}{TP+FP}.
\end{equation}
Precision~=~0,95 означава, че от 100 аларми 95 са истински атаки, а 5 са фалшиви. В примера на материала за Precision е описано тълкуване на Recall. Тук е използвано коректното тълкуване: делът на верните сред положителните решения~\cite{fawcett2006}.

\paragraph{Чувствителност (Sensitivity, Recall)} е делът на откритите атаки:
\begin{equation}\label{eq:sensitivity}
\operatorname{Sensitivity}=\operatorname{Recall}=\operatorname{TPR}=\frac{TP}{TP+FN}.
\end{equation}
Recall~=~0,97 означава, че от 100 атаки 97 се откриват, а 3 се пропускат. Тъй като FN е най-нежеланата грешка, Recall е водещата метрика за сигурността.

\paragraph{F1 оценка} е хармоничното средно на Precision и Recall:
\begin{equation}\label{eq:f1}
\operatorname{F1}=2\cdot\frac{\operatorname{Precision}\cdot\operatorname{Recall}}{\operatorname{Precision}+\operatorname{Recall}}=\frac{2\,TP}{2\,TP+FP+FN}.
\end{equation}
F1 е висока само когато и двете съставки са високи. Затова тя е предпочитаната обобщена метрика при несиметрични данни, при които Accuracy подвежда~\cite{borovska2022,saito2015}.

\paragraph{Специфичност (Specificity)} е делът на нормалния трафик, който не предизвиква аларма:
\begin{equation}\label{eq:spec}
\operatorname{Specificity}=\operatorname{TNR}=\frac{TN}{TN+FP},\qquad \operatorname{FPR}=1-\operatorname{Specificity}.
\end{equation}
Specificity~=~0,98 означава, че от 100 нормални записа 98 минават без аларма, а 2 предизвикват фалшива тревога.

\paragraph{ROC крива и реална честота на атаките.} Трите модела дават непрекъсната оценка $s(\vect{x})$ ($\hat p$ при RF и XGBoost и $f$ при SVM). При праг $\theta$ решението е $\hat y=\mathbf{1}[s(\vect{x})\ge\theta]$. ROC кривата съдържа точките $(\operatorname{FPR}(\theta),\operatorname{TPR}(\theta))$ за всички прагове, а площта под нея е вероятността случайна атака да получи по-висока оценка от случаен нормален запис~\cite{fawcett2006}:
\begin{equation}\label{eq:auc}
\operatorname{AUC}=\int_0^1\operatorname{TPR}\,d\operatorname{FPR}=P\bigl(s(\vect{x}^{+})>s(\vect{x}^{-})\bigr).
\end{equation}
Precision зависи и от дела $\pi$ на атаките в трафика. От теоремата на Бейс следва~\cite{axelsson2000}
\begin{equation}\label{eq:baserate}
\operatorname{Precision}(\pi)=\frac{\operatorname{TPR}\cdot\pi}{\operatorname{TPR}\cdot\pi+\operatorname{FPR}\cdot(1-\pi)}.
\end{equation}
При малко $\pi$ дори нисък FPR дава повече фалшиви, отколкото истински аларми. Това е \emph{проблемът на базовата честота}, изследван в експеримент E2 с Precision--Recall криви~\cite{saito2015}. Там се избират и работни прагове: най-високият праг, при който Recall достига целева стойност (0,90 или 0,95), за да се ограничат пропуснатите атаки.

% ---------------------------------------------------------------------------
\section{Графично представяне и анализ на резултатите}\label{sec:stats}

Резултатите се анализират в два слоя: статистическа проверка на надеждността на метриките и графично представяне на всички етапи на експеримента~\cite{efron1993}.

\paragraph{Bootstrap доверителни интервали.} Несигурността на метриките се оценява с непараметричен bootstrap~\cite{efron1993}: от $n$-те тестови записа се теглят с връщане $R=\num{1000}$ извадки с размер $n$, върху всяка се изчислява метриката $\hat M^{*}_r$, а 95\% перцентилният интервал е
\begin{equation}\label{eq:bootstrap}
\operatorname{CI}_{95\%}(M)=\left[\hat{M}^{*}_{(0{,}025)},\ \hat{M}^{*}_{(0{,}975)}\right].
\end{equation}

\paragraph{Тест на McNemar.} Всяка двойка класификатори A и B се сравнява върху едни и същи тестови записи по записите, по които се различават~\cite{mcnemar1947,dietterich1998}: $b$ е броят записи, верни при A и грешни при B, а $c$ е обратното. При нулевата хипотеза за еднакъв процент грешки
\begin{equation}\label{eq:mcnemar}
\chi^2=\frac{\left(|b-c|-1\right)^2}{b+c}\ \sim\ \chi^2_1,
\end{equation}
а при $b+c<25$ се използва точният биномен тест.

\paragraph{Едномерна разделимост.} За всеки признак се сравняват емпиричните функции на разпределение при атака ($F_1$) и нормален трафик ($F_0$) със статистиката на Колмогоров--Смирнов
\begin{equation}\label{eq:ks}
D_j=\sup_{t\in\R}\ \bigl|F_{1,j}(t)-F_{0,j}(t)\bigr|\in[0,1].
\end{equation}
$D_j=1$ означава, че един праг по признак $j$ разделя класовете идеално~\cite{hastie2009}.

\paragraph{Графично представяне.} Всички графики се генерират автоматично след всяко изпълнение с еднакво цветово кодиране: Random Forest в синьо, SVM в оранжево, XGBoost в зелено, атаката в червено, нормалният трафик в сиво. Те покриват данните и обработката, качеството на моделите, поведението им при обучение и работата при реално съотношение. След всяка фигура в \cref{ch:results} следва обяснение какво е показано~\cite{repo}.

% ============================== 2. ТЕХНОЛОГИЧНА РАМКА ==============================
\chapter{Експериментална технологична рамка}\label{ch:framework}

\section{Хардуер и софтуер}
Експериментът е реализиран на Python със стандартни библиотеки за анализ на данни и машинно обучение (\cref{tab:env}). Точните версии са фиксирани в \texttt{requirements.txt}, а параметрите на средата се записват автоматично при всяко изпълнение във файла \path{results/environment.json}.

\begin{table}[!htbp]
\centering\small
\caption{Експериментална среда}\label{tab:env}
\fitbox{%
\begin{tabularx}{\textwidth}{@{}lL@{}}
\toprule
\headrow \hd{Компонент} & \hd{Версия / описание}\\
\midrule
Процесор & AMD64 Family 25 Model 97 Stepping 2, AuthenticAMD (16 ядра)\\
\rowcolor{rowalt}Оперативна памет & 0,0 GB\\
Операционна система & Windows-11-10.0.26300-SP0\\
\rowcolor{rowalt}Език & Python 3.13.16\\
Анализ на данни & pandas 3.0.5, NumPy 2.5.3, SciPy 1.18.1\\
\rowcolor{rowalt}Машинно обучение & scikit-learn 1.9.1 (RandomForestClassifier, SVC/LIBSVM, PCA), XGBoost 3.4.1\\
Визуализация & matplotlib, seaborn\\
\rowcolor{rowalt}Документация & XeLaTeX (този доклад), pptxgenjs (презентация)\\
Тестове & pytest (8 теста върху реални записи)\\
\bottomrule
\end{tabularx}}%

\end{table}

\section{Архитектура на програмната реализация}
Кодът е организиран като малък пакет с ясно разделени отговорности (\cref{fig:architecture}). Централната идея е \emph{декларативното описание} на експеримента. Атаките и алгоритмите не са вписани в логиката на конвейера, а са записи в два регистъра (\texttt{src/registry.py}). Класът \texttt{Experiment} обхожда регистрите и за всяка двойка (атака, алгоритъм) изпълнява едни и същи стъпки. Така при разширяването на екипа третата атака (BruteForce, по-късно и Bot) и третият алгоритъм (XGBoost) бяха добавени с по един нов запис в регистъра, без промяна на останалия код. Смяната на набора от данни изисква само нов зареждащ модул и нови записи за атаките.

\begin{figure}[!htbp]
\centering
\fitbox{\begin{tikzpicture}[font=\sffamily\small,
  bxmod/.style={draw=accent,fill=boxbg,rounded corners=2pt,minimum width=3.2cm,minimum height=1.05cm,align=center,line width=0.7pt},
  bxcore/.style={bxmod,draw=navy,fill=navy,text=white,font=\sffamily\small\bfseries},
  bxout/.style={bxmod,draw=orange,fill=white},
  bxarr/.style={-{Latex[length=2.2mm]},color=muted,line width=0.7pt}]
\node[bxmod] (cfg) at (0,0) {config.py\\[-1pt]\scriptsize параметри, seed};
\node[bxmod] (reg) at (0,-1.6) {registry.py\\[-1pt]\scriptsize AttackSpec, AlgorithmSpec};
\node[bxcore] (exp) at (4.4,-0.8) {experiment.py\\[-1pt]\scriptsize клас Experiment};
\node[bxmod] (pre) at (9,1.2) {preprocess.py\\[-1pt]\scriptsize зареждане, почистване, задачи,\\[-2pt]\scriptsize FeatureFilter, ImportanceSelector};
\node[bxmod] (trn) at (9,-0.25) {train.py\\[-1pt]\scriptsize Pipeline, GridSearchCV, 1-SE};
\node[bxmod] (evl) at (9,-1.7) {evaluate.py\\[-1pt]\scriptsize оперативна фаза, метрики};
\node[bxmod] (ana) at (9,-3.15) {analysis.py, stats.py\\[-1pt]\scriptsize A1--A9, bootstrap, McNemar};
\node[bxout] (plt) at (4.4,-3.15) {plots.py, plots\_ext.py\\[-1pt]\scriptsize 26 фигури};
\node[bxout] (res) at (0,-3.15) {results/, models/\\[-1pt]\scriptsize CSV, JSON, joblib};
\node[bxmod,draw=muted] (cli) at (4.4,1.2) {run\_all.py\\[-1pt]\scriptsize команден ред (CLI)};
\draw[bxarr] (cli) -- (exp);
\draw[bxarr] (cfg) -- (exp);
\draw[bxarr] (reg) -- (exp);
\draw[bxarr] (exp.east) -- ++(0.8,0) |- (pre.west);
\draw[bxarr] (exp.east) -- ++(0.8,0) |- (trn.west);
\draw[bxarr] (exp.east) -- ++(0.8,0) |- (evl.west);
\draw[bxarr] (exp.east) -- ++(0.8,0) |- (ana.west);
\draw[bxarr] (exp) -- (plt);
\draw[bxarr] (plt) -- (res);
\draw[bxarr] (exp.south west) to[bend right=15] (res.north east);
\node[bxmod,draw=muted,fill=white] (tst) at (13,-0.8) {tests/\\[-1pt]\scriptsize pytest върху\\[-2pt]\scriptsize реални записи};
\draw[bxarr,dashed] (tst.west) -- ++(-0.6,0);
\end{tikzpicture}}
\caption{Архитектура на програмната реализация}\label{fig:architecture}
\vspace{2pt}
\interpret{Модулите в синьо реализират обработката и моделите, тъмният модул \texttt{experiment.py} ги оркестрира, а модулите в оранжево произвеждат резултатите. Конфигурацията и регистрите са единственото място, където се описва \emph{какво} се изследва. Останалите модули описват \emph{как}. Автоматичните тестове проверяват основните свойства на конвейера върху реални записи.}
\end{figure}

Обработката и класификаторът образуват един обект \texttt{Pipeline} (филтър, подбор на признаци, стандартизация, класификатор). Той се обучава, записва и зарежда като цяло, затова в оперативната фаза новите записи минават през точно същите трансформации~\cite{pedregosa2011}. Класът \texttt{Experiment} има четири етапа: подготовка на данните, деветте модела с контролните експерименти E1--E5 и диагностиката на преобучението, разширените анализи A1--A9 и фигурите. Проверките за приложимост A10--A14 са в модула \texttt{src/robustness.py} и се изпълняват след основния експеримент (около 8 минути). Изпълнението се стартира от командния ред с \texttt{python run\_all.py} (около 25 минути), като атаките и алгоритмите могат да се избират с параметри. Инструкциите за възпроизвеждане са в \cref{app:repro}~\cite{repo}.

\section{Възпроизводимост и проверки}
Възпроизводимостта се осигурява с фиксиран seed (42) за всяка случайна операция, фиксирани версии на библиотеките в \texttt{requirements.txt} и пълен журнал на междинните резултати в CSV/JSON файлове~\cite{arp2022}. Осем автоматични теста (\texttt{tests/}) проверяват върху реални записи от CSE-CIC-IDS2018, че обучаващите и тестовите вектори не се пресичат дори по съдържание, че класовете са балансирани, че подборът на признаци се обучава само върху обучаващите данни, че филтърът оставя само признаци с $|r|\le0{,}95$, че метриките съвпадат с дефинициите (\ref{eq:accuracy})--(\ref{eq:spec}), че праговете по Recall постигат целта и че bootstrap интервалите и тестът на McNemar са коректни. Независима проверка зарежда записаните модели, построява тестовите данни наново и получава същите TP, TN, FP и FN като в резултатите.

% ============================== 3. ПЛАН НА ЕКСПЕРИМЕНТИТЕ ==============================
\chapter{План на експериментите (фаза на машинно обучение и оперативна фаза)}\label{ch:plan}

Експериментите следват работния цикъл на реална IDS и се разделят на две фази: \emph{фаза на машинно обучение}, в която моделите се синтезират, и \emph{оперативна фаза}, в която записаните модели работят върху нови данни (\cref{fig:pipeline}). Около основния експеримент са планирани пет контролни експеримента (E1--E5) девет разширени анализа (A1--A9) и пет проверки за приложимост (A10--A14).

\begin{figure}[!htbp]
\centering
\includegraphics[width=\textwidth]{00_pipeline.png}
\caption{Експериментален конвейер: от суровите CSV файлове до метриките}\label{fig:pipeline}
\vspace{2pt}
\interpret{Горният ред показва подготовката на данните, обща за двете фази. Синята област е фазата на машинно обучение: върху 70\% от данните се обучава един \texttt{Pipeline} (филтър, подбор на признаци, стандартизация, класификатор), който се записва на диска. Оранжевата област е оперативната фаза: записаният модел се зарежда и класифицира останалите 30\%, които не е виждал, на пакети по \num{10000} записа, както би работила IDS върху поток от записи.}
\end{figure}

\section{Фаза на машинно обучение}
За всяка от трите атаки и всеки от трите алгоритъма се изпълняват следните стъпки:
\begin{enumerate}
  \item формиране на балансираната бинарна задача (\cref{eq:binary}) и стратифицирано разделяне 70/30 (\cref{eq:split});
  \item избор на хиперпараметрите с 5-кратна кръстосана валидация и правилото „една стандартна грешка“ (\cref{eq:cv,eq:one-se}). Решетките са: за RF брой дървета $\{200, 300\}$, максимална дълбочина $\{\infty, 25\}$, минимален брой записи в листо $\{2, 5, 10\}$ и дял на признаците при разделяне $\{\sqrt{d}, 0{,}5\}$ (24 комбинации); за XGBoost брой дървета $\{300, 800\}$, дълбочина $\{6, 10\}$, скорост на обучение $\{0{,}03;\ 0{,}1\}$ и \texttt{min\_child\_weight} $\{1, 5\}$ (16 комбинации); за SVM $C\in\{1, 10, 100\}$ и $\gamma\in\{\txt{scale}, 0{,}01;\ 0{,}1\}$ (9 комбинации). RF и XGBoost се настройват върху \num{50000} обучаващи записа, а SVM върху \num{20000}. Всяка комбинация се оценява с целия конвейер;
  \item обучение на финалния конвейер с избраните параметри: RF и XGBoost върху всички обучаващи записи, SVM върху \num{50000} от тях заради сложността $O(n^2)$--$O(n^3)$;
  \item диагностика на преобучението: точност върху обучаващите данни срещу оценката от кръстосаната валидация и точността върху тестовите данни;
  \item запис на конвейера във файл (\texttt{models/<алгоритъм>\_<атака>.joblib}).
\end{enumerate}

\section{Оперативна фаза (фаза на тестване/експлоатация)}
\begin{enumerate}
  \item зареждане на записания модел от диска, без достъп до обучаващите данни;
  \item подаване на тестовите записи (30\%) на последователни пакети по \num{10000} и измерване на времето за класификация;
  \item сравняване на предсказанията с реалните етикети и изчисляване на матрицата на объркване и метриките (\ref{eq:accuracy})--(\ref{eq:auc});
  \item статистическа оценка: bootstrap интервали (\cref{eq:bootstrap}) и тест на McNemar за всяка двойка алгоритми (\cref{eq:mcnemar}).
\end{enumerate}

\section{Контролни експерименти и разширени анализи}
Контролните експерименти проверяват конкретни методологични решения, а разширените анализи обясняват поведението на моделите (\cref{tab:experiments}).

\begin{table}[!htbp]
\centering\small
\caption{План на допълнителните експерименти и анализи}\label{tab:experiments}
\begin{tabularx}{\textwidth}{@{}l p{4.4cm} L@{}}
\toprule
\textsf{\textbf{№}} & \textsf{\textbf{Експеримент}} & \textsf{\textbf{Въпрос, на който отговаря}}\\
\midrule
E1 & Пълен вектор (без редукция) & Каква е цената или ползата от редукцията на размерността?\\
E2 & Тест при естествено съотношение & Колко полезни са алармите при реалния дял на атаките (\cref{eq:baserate})?\\
E3 & Стабилност при кръстосана валидация & Колко се колебае качеството между фолдовете?\\
E4 & Без времевите признаци на потока & Зависят ли моделите от продължителността и интервалите между пакетите, които са най-чувствителни към натоварването на мрежата?\\
E5 & PCA вместо подбор на признаци & Кой метод за редукция е по-подходящ?\\
\midrule
A1 & Профил и KS разделимост & Кои признаци различават атаката от нормалния трафик (\cref{eq:ks})?\\
A2 & 2D проекции (PCA, t-SNE) & Разделими ли са класовете геометрично?\\
A3 & Криви на обучение & Колко данни са нужни и би ли помогнало повече обучение?\\
A4 & Брой запазени измерения & Колко признака или компоненти са необходими и какво печелим по скорост?\\
A5 & Праг при естествено съотношение & Какъв праг осигурява дадено ниво на Precision или Recall?\\
A6 & Подвидове на атаките & Кои разновидности (HOIC, LOIC-UDP, FTP, SSH) се пропускат?\\
A7 & Пренос между атаките & Учат ли моделите конкретната атака или общо „аномалия“?\\
A8 & Пермутационна важност & На кои признаци реално разчита всеки модел?\\
A9 & Bootstrap и McNemar & Значими ли са разликите между алгоритмите?\\
\midrule
A10 & Без признаци-отпечатъци на TCP стека & Разпознават ли моделите атаката или само инструмента и операционната система на атакуващия?\\
A11 & Хронологично разделяне & Открива ли моделът атаката и по-късно във времето, а не само в случайно избрани записи?\\
A12 & Пренос между подвидове & Открива ли модел, обучен върху един инструмент (HOIC, SSH), атака със същата цел, но с друг инструмент?\\
A13 & Нормален трафик от непознат ден & Колко фалшиви аларми дава моделът при нормален трафик, различен от този при обучението?\\
A14 & Комбиниран детектор (XGBoost или SVM) & Може ли съчетанието на двата алгоритъма да открива и нови разновидности на атаката?\\
\bottomrule
\end{tabularx}
\end{table}

Проверките A10--A14 се изпълняват отделно (\texttt{python -m src.robustness}) с хиперпараметрите, избрани в основния експеримент. За тях записите се подреждат по време. В CSV файловете на CSE-CIC-IDS2018 часът е записан в 12-часов формат без означение преди и след обяд. Затова часовете от 1 до 7 се приемат за следобедни: така атаките попадат точно в публикувания график на набора (например LOIC-UDP от 10:09 до 10:43 и HOIC от 14:11 до 14:33 на 21.02.2018)~\cite{sharafaldin2018}. При повтарящ се вектор на признаците се запазва най-ранното му срещане, за да не попадне един и същ запис и в ранните, и в късните данни. Нормалният трафик в A10, A11 и A14 е от деня на съответната атака, за да има смисъл подреждането по време; в A13 се сравнява нормален трафик от един и от два дни.


\paragraph{Детайли на анализите.} При A3 се обучават модели върху $\{10^3, 2\cdot10^3, 5\cdot10^3, 10^4, 2\cdot10^4, 5\cdot10^4, 10^5\}$ записа (SVM до $5\cdot10^4$). При A4 броят запазени измерения е $k\in\{2,5,10,15,20,30,40\}$ за двата метода (подбор на признаци и PCA) плюс контрол без редукция. Двата анализа използват обучаваща извадка до \num{20000} записа (за A4) и обща тестова извадка от \num{10000} невиждани записа. При E2 се използват всички записи, които не са участвали в обучението, смесени в естественото съотношение на класовете в изтеглените файлове (\val{nat-share-DDoS} DDoS, \val{nat-share-BruteForce} BruteForce и \val{nat-share-Bot} Bot). При A5 за всяко целево ниво на Precision се търси прагът $\theta$ с максимален Recall, а за целеви Recall от 0,90 и 0,95 се търси най-високият праг, който го постига. Пермутационната важност (A8) се изчислява върху \num{2000} тестови записа с 3 повторения~\cite{fisher2019}.

% ============================== 4. ИМПЛЕМЕНТАЦИЯ И РЕЗУЛТАТИ ==============================
\chapter{Имплементация и експериментални резултати}\label{ch:results}

Главата представя резултатите в реда на експерименталния конвейер: данни и обработка, фаза на машинно обучение и проверка за преобучение, оперативна фаза, статистическа надеждност, ефект на редукцията и работа при естествено съотношение. Всички числа, таблици и фигури са генерирани автоматично от последното изпълнение на \texttt{run\_all.py}~\cite{repo}.

% ---------------------------------------------------------------------------
\section{Данни и предварителна обработка}\label{sec:r-data}

От изтеглените \val{rows-read} записа \val{nan-rows} съдържат липсващи или безкрайни стойности, а \val{dup-rows} (\val{dup-share}) са пълни дубликати. След почистването остават \val{rows-clean} уникални записа (\cref{tab:class-counts}). Почистването на ниво задача (\cref{eq:conflict}) премахва още \val{task-conf-Bot} записа (\val{task-confvec-Bot} вектора) с противоречиви етикети в задачата Bot, \val{task-conf-DDoS} в задачата DDoS и \val{task-dup-BruteForce} дубликата в задачата BruteForce. От всеки клас се изтеглят по \val{per-class-DDoS} записа за DDoS и Bot и \val{per-class-BruteForce} за BruteForce, т.е. всички уникални записи на BruteForce. Така обучаващите множества съдържат \val{train-rows-DDoS}, \val{train-rows-BruteForce} и \val{train-rows-Bot} записа, а тестовите \val{test-rows-DDoS}, \val{test-rows-BruteForce} и \val{test-rows-Bot} (\cref{fig:classes}).

\begin{figure}[H]
\centering
\includegraphics[width=\textwidth]{01_class_distribution.png}
\caption{Разпределение на класовете в изтеглените данни и след балансирането (логаритмична скала)}\label{fig:classes}
\vspace{2pt}
\interpret{Сивите стълбове са нормалният трафик, а червените атаката. В оригиналните файлове нормалният трафик е 2,6 (DDoS) до 6,3 пъти (Bot) повече от атаката. Вдясно е балансираната извадка 1:1, върху която се обучават и тестват моделите. Оригиналното съотношение се използва отново в експеримент E2.}
\end{figure}

\Cref{fig:ks} подрежда признаците по KS статистиката от \cref{eq:ks} (анализ A1). При BruteForce един признак, \texttt{\val{ks-top-BruteForce}}, сам разделя класовете ($D=\val{ks-topD-BruteForce}$), а средното $D$ на петте най-добри признака е \val{ks-top5-BruteForce}. При DDoS най-разделим е \texttt{\val{ks-top-DDoS}} ($D=\val{ks-topD-DDoS}$), а при Bot \texttt{\val{ks-top-Bot}} ($D=\val{ks-topD-Bot}$). Признаци с $D>0{,}5$ има \val{ks-n05-DDoS} за DDoS, \val{ks-n05-BruteForce} за BruteForce и \val{ks-n05-Bot} за Bot.

\begin{figure}[H]
\centering
\includegraphics[width=\textwidth]{12_ks_separability.png}
\caption{Едномерна разделимост на признаците: KS статистика $D$ между атаката и нормалния трафик}\label{fig:ks}
\vspace{2pt}
\interpret{Дължината на стълба е максималната разлика между кумулативните разпределения на признака при атака и при нормален трафик. Стойност близо до 1 означава, че само този признак почти разделя класовете. При BruteForce много признаци имат $D>0{,}9$: опитите за вход са кратки потоци с еднаква структура. При DDoS и Bot най-добрите признаци са около 0,75, т.е. всеки сам по себе си не стига, но заедно разделят класовете.}
\end{figure}

Редукцията свива вектора на признаците в два етапа, както е обобщено в \cref{tab:reduction}. Филтърът премахва \val{lowvar-DDoS} константни признака (\texttt{\val{lowvar-names}}), а подборът запазва \val{n-selected} от останалите \val{n-filtered}. Избраните признаци носят \val{sel-share-DDoS} (DDoS), \val{sel-share-BruteForce} (BruteForce) и \val{sel-share-Bot} (Bot) от общата важност в помощния модел. Най-важни са \texttt{\val{sel-top3-DDoS}} за DDoS, \texttt{\val{sel-top3-BruteForce}} за BruteForce и \texttt{\val{sel-top3-Bot}} за Bot.

\begin{table}[H]
\centering\small
\caption{Редукция на размерността по задача: филтър и подбор на признаци}\label{tab:reduction}
\fitbox{%
\begin{tabular}{@{}lccccc>{\raggedright\arraybackslash}p{6.2cm}@{}}
\toprule
\headrow \hd{Задача} & \hd{Вход} & \hd{Нулева\\дисперсия} & \hd{След\\филтър} & \hd{Избрани} & \hd{Дял на\\важността} & \hd{Избрани признаци}\\
\midrule
DDoS & 77 & 10 & 67 & 20 & 99,9\% & \scriptsize\texttt{Protocol, Flow Duration, Tot Fwd Pkts, Tot Bwd Pkts, TotLen Bwd Pkts, Bwd Pkt Len Mean, Bwd Pkt Len Std, Flow Pkts/s, Flow IAT Max, Flow IAT Min, Fwd IAT Tot, Fwd IAT Max, Fwd IAT Min, Fwd Header Len, Fwd Pkts/s, Pkt Len Mean, Pkt Len Std, Init Fwd Win Byts, Init Bwd Win Byts, Fwd Act Data Pkts}\\
\rowcolor{rowalt}BruteForce & 77 & 10 & 67 & 20 & 100,0\% & \scriptsize\texttt{Flow Duration, Tot Fwd Pkts, Tot Bwd Pkts, TotLen Bwd Pkts, Bwd Pkt Len Max, Bwd Pkt Len Mean, Flow Pkts/s, Flow IAT Mean, Flow IAT Max, Flow IAT Min, Fwd IAT Min, Bwd IAT Mean, Bwd IAT Min, Fwd Header Len, Fwd Pkts/s, Pkt Len Std, Init Fwd Win Byts, Init Bwd Win Byts, Fwd Act Data Pkts, Fwd Seg Size Min}\\
Bot & 77 & 10 & 67 & 20 & 99,5\% & \scriptsize\texttt{Protocol, Tot Fwd Pkts, Fwd Pkt Len Max, Fwd Pkt Len Mean, Bwd Pkt Len Max, Bwd Pkt Len Mean, Flow Byts/s, Flow IAT Std, Flow IAT Max, Flow IAT Min, Fwd IAT Std, Fwd IAT Min, Fwd Pkts/s, Pkt Len Max, Pkt Len Mean, RST Flag Cnt, Pkt Size Avg, Init Fwd Win Byts, Active Mean, Idle Min}\\
\bottomrule
\end{tabular}}%

\vspace{2pt}
\interprettab{За всяка задача е показан броят признаци на входа, броят премахнати като константни, броят след филтъра, броят избрани признаци и делът на общата важност, който те носят. Векторът се свива от \val{n-features} до \val{n-selected} измерения (с 74\%), като почти цялата информация се запазва. Подборът се извършва отделно за всяка задача, защото различните атаки се разпознават по различни признаци.}
\end{table}

Двумерните проекции PCA и t-SNE~\cite{vandermaaten2008} (анализ A2) показват, че и при трите атаки атакуващите записи образуват компактни групи, отделени от нормалния трафик. За контролния метод PCA 95\% от дисперсията се запазват с \val{pca-DDoS} компонента за DDoS, \val{pca-BruteForce} за BruteForce и \val{pca-Bot} за Bot (\cref{eq:kstar}).

% ---------------------------------------------------------------------------
\section{Фаза на машинно обучение и проверка за преобучение}\label{sec:r-ml}

\Cref{tab:cv} показва избраните хиперпараметри. Правилото „една стандартна грешка“ (\cref{eq:one-se}) избира по-проста конфигурация от тази с най-висока оценка в няколко случая (например SVM за Bot с $C=1$ вместо $C=10$), без статистически значима загуба. Разсейването на F1 между петте фолда (експеримент E3) е в рамките на стотни от процента при DDoS и BruteForce и около една десета от процента при Bot.

\begin{table}[H]
\centering\small
\caption{Резултат от GridSearchCV: избрани хиперпараметри и CV F1}\label{tab:cv}
\fitbox{%
\begin{tabular}{@{}llp{4.4cm}rrrrr@{}}
\toprule
\headrow \hd{Модел} & \hd{Атака} & \hd{Избрани хиперпараметри} & \hd{CV F1\\избрана} & \hd{CV F1\\най-висока} & \hd{Ст. откл.} & \hd{Мин.--макс. по фолд} & \hd{Време, s}\\
\midrule
Random Forest & DDoS & \scriptsize\texttt{max\_depth=25, max\_features=0.5, min\_samples\_leaf=2, n\_estimators=200} & 99,99\% & 99,99\% & 0,01\% & 99,98\% -- 100,00\% & 241\\
\rowcolor{rowalt}SVM (RBF) & DDoS & \scriptsize\texttt{C=10, gamma=0.1} & 99,96\% & 99,97\% & 0,03\% & 99,90\% -- 100,00\% & 12\\
XGBoost & DDoS & \scriptsize\texttt{learning\_rate=0.03, max\_depth=6, min\_child\_weight=1, n\_estimators=300} & 100,00\% & 100,00\% & 0,00\% & 100,00\% -- 100,00\% & 61\\
\rowcolor{rowalt}Random Forest & BruteForce & \scriptsize\texttt{max\_depth=25, max\_features=sqrt, min\_samples\_leaf=2, n\_estimators=200} & 99,99\% & 99,99\% & 0,01\% & 99,97\% -- 100,00\% & 189\\
SVM (RBF) & BruteForce & \scriptsize\texttt{C=10, gamma=scale} & 99,96\% & 99,97\% & 0,03\% & 99,92\% -- 100,00\% & 10\\
\rowcolor{rowalt}XGBoost & BruteForce & \scriptsize\texttt{learning\_rate=0.03, max\_depth=6, min\_child\_weight=5, n\_estimators=300} & 99,99\% & 99,99\% & 0,01\% & 99,98\% -- 100,00\% & 53\\
Random Forest & Bot & \scriptsize\texttt{max\_depth=25, max\_features=sqrt, min\_samples\_leaf=2, n\_estimators=200} & 99,93\% & 99,94\% & 0,03\% & 99,90\% -- 99,98\% & 226\\
\rowcolor{rowalt}SVM (RBF) & Bot & \scriptsize\texttt{C=1, gamma=scale} & 99,47\% & 99,49\% & 0,10\% & 99,35\% -- 99,60\% & 15\\
XGBoost & Bot & \scriptsize\texttt{learning\_rate=0.1, max\_depth=6, min\_child\_weight=5, n\_estimators=800} & 99,96\% & 99,97\% & 0,02\% & 99,94\% -- 99,99\% & 70\\
\bottomrule
\end{tabular}}%
\par\smallskip{\footnotesize\color{muted}Избраната конфигурация е най-простата в рамките на една стандартна грешка от най-високата CV оценка (правило 1-SE).}

\vspace{2pt}
\interprettab{За всеки модел е посочена избраната конфигурация, нейното средно F1 по петте фолда, най-високото F1 в решетката, стандартното отклонение, най-ниската и най-високата стойност по фолд и времето за търсене. Всички модели постигат CV F1 над 99,4\%. Най-ниско е F1 на SVM при Bot (\val{SVM-Bot-cv}).}
\end{table}

\Cref{tab:overfit} проверява явно за преобучение. Преобученият модел има много по-висока точност върху обучаващите данни, отколкото върху новите. Тук разликата е най-много \val{gap-max} процентни пункта, а F1 от кръстосаната валидация съвпада с F1 върху тестовите данни. Следователно моделите обобщават, а не запомнят обучаващите записи.

\begin{table}[H]
\centering\small
\caption{Проверка за преобучение: обучение, кръстосана валидация и тест}\label{tab:overfit}
\fitbox{%
\begin{tabular}{@{}llrrrrr@{}}
\toprule
\headrow \hd{Модел} & \hd{Атака} & \hd{Accuracy\\обучение} & \hd{F1\\5-CV} & \hd{Accuracy\\тест} & \hd{F1\\тест} & \hd{Разлика обуч.$-$тест,\\пр. п.}\\
\midrule
Random Forest & DDoS & 100,000\% & 99,994\% & 100,000\% & 100,000\% & $+$0,000\\
\rowcolor{rowalt}SVM (RBF) & DDoS & 99,954\% & 99,955\% & 99,950\% & 99,950\% & $+$0,004\\
XGBoost & DDoS & 100,000\% & 100,000\% & 100,000\% & 100,000\% & $+$0,000\\
\rowcolor{rowalt}Random Forest & BruteForce & 99,998\% & 99,992\% & 99,998\% & 99,998\% & $+$0,000\\
SVM (RBF) & BruteForce & 99,972\% & 99,965\% & 99,986\% & 99,986\% & $-$0,014\\
\rowcolor{rowalt}XGBoost & BruteForce & 99,997\% & 99,990\% & 99,998\% & 99,998\% & $-$0,001\\
Random Forest & Bot & 99,983\% & 99,932\% & 99,967\% & 99,967\% & $+$0,016\\
\rowcolor{rowalt}SVM (RBF) & Bot & 99,586\% & 99,473\% & 99,598\% & 99,597\% & $-$0,012\\
XGBoost & Bot & 99,991\% & 99,960\% & 99,982\% & 99,982\% & $+$0,010\\
\bottomrule
\end{tabular}}%

\vspace{2pt}
\interprettab{Колоните сравняват Accuracy върху обучаващите данни, F1 от 5-кратната кръстосана валидация (данни, невиждани при обучението на съответния фолд) и Accuracy и F1 върху тестовите данни от оперативната фаза. Отрицателна разлика означава, че моделът е малко по-точен върху тестовите данни, отколкото върху обучаващите, което е в рамките на случайността.}
\end{table}

Кривите на обучение (анализ A3, \texttt{results/learning\_curves.csv}) показват, че още с \val{lc-RF-DDoS-nfirst} записа всички модели надхвърлят 99,3\%. С повече данни F1 при Bot расте до \val{lc-RF-Bot-last} (RF) и \val{lc-XGB-Bot-last} (XGBoost), а SVM расте по-бавно (\val{lc-SVM-Bot-first} до \val{lc-SVM-Bot-last}) и времето му за обучение расте по-бързо от линейно.

% ---------------------------------------------------------------------------
\section{Оперативна фаза: основни резултати}\label{sec:r-op}

\Cref{tab:main} съдържа всички величини от матрицата на объркване и метриките (\ref{eq:accuracy})--(\ref{eq:spec}) на деветте модела. Те са изчислени от записаните модели върху тестови записи, които не са участвали в обучението.

\begin{table}[H]
\centering\small
\caption{Основни резултати в оперативната фаза}\label{tab:main}
\fitbox{%
\begin{tabular}{@{}llrrrrrrrrr@{}}
\toprule
\headrow \hd{Модел} & \hd{Атака} & \hd{TP} & \hd{TN} & \hd{FP} & \hd{FN} & \hd{Accuracy} & \hd{Precision} & \hd{Recall\\(Sensitivity)} & \hd{F1} & \hd{Specificity}\\
\midrule
Random Forest & PortScan & 504 & 504 & 1 & 1 & 99,80\% & 99,80\% & 99,80\% & 99,80\% & 99,80\%\\
\rowcolor{rowalt}SVM (RBF) & PortScan & 501 & 503 & 2 & 4 & 99,41\% & 99,60\% & 99,21\% & 99,40\% & 99,60\%\\
XGBoost & PortScan & 504 & 503 & 2 & 1 & 99,70\% & 99,60\% & 99,80\% & 99,70\% & 99,60\%\\
\rowcolor{rowalt}Random Forest & DDoS & 30\,000 & 30\,000 & 0 & 0 & 100,00\% & 100,00\% & 100,00\% & 100,00\% & 100,00\%\\
SVM (RBF) & DDoS & 29\,988 & 29\,982 & 18 & 12 & 99,95\% & 99,94\% & 99,96\% & 99,95\% & 99,94\%\\
\rowcolor{rowalt}XGBoost & DDoS & 30\,000 & 30\,000 & 0 & 0 & 100,00\% & 100,00\% & 100,00\% & 100,00\% & 100,00\%\\
Random Forest & BruteForce & 28\,227 & 28\,229 & 0 & 1 & 100,00\% & 100,00\% & 100,00\% & 100,00\% & 100,00\%\\
\rowcolor{rowalt}SVM (RBF) & BruteForce & 28\,221 & 28\,228 & 1 & 7 & 99,99\% & 100,00\% & 99,98\% & 99,99\% & 100,00\%\\
XGBoost & BruteForce & 28\,227 & 28\,229 & 0 & 1 & 100,00\% & 100,00\% & 100,00\% & 100,00\% & 100,00\%\\
\bottomrule
\end{tabular}}%

\vspace{2pt}
\interprettab{TP и TN са верните решения, FP са фалшивите аларми, а FN са пропуснатите атаки. Всички модели имат Accuracy над 99,5\%. Random Forest и XGBoost не правят нито една грешка при DDoS и по една грешка при BruteForce. Най-много грешки има SVM при Bot: \val{SVM-Bot-FN} пропуснати атаки и \val{SVM-Bot-FP} фалшиви аларми от \val{test-rows-Bot} записа.}
\end{table}

\Cref{tab:per100} превежда метриките в броя случаи от 100, както препоръчва учебният материал~\cite{borovska2022}. Това тълкуване е по-разбираемо за оператора на IDS от процентите.

\begin{table}[H]
\centering\scriptsize
\caption{Тълкуване на метриките „на 100 случая“}\label{tab:per100}
\fitbox{%
\begin{tabular}{@{}llrrrrrrrrr@{}}
\toprule
\headrow \hd{Метрика} & \hd{Какво означава} & \hd{Random Forest PortScan} & \hd{SVM (RBF) PortScan} & \hd{XGBoost PortScan} & \hd{Random Forest DDoS} & \hd{SVM (RBF) DDoS} & \hd{XGBoost DDoS} & \hd{Random Forest BruteForce} & \hd{SVM (RBF) BruteForce} & \hd{XGBoost BruteForce}\\
\midrule
Accuracy & правилни решения от 100 тестови записа & 100 от 100 & 99 от 100 & 100 от 100 & 100 от 100 & 100 от 100 & 100 от 100 & 100 от 100 & 100 от 100 & 100 от 100\\
\rowcolor{rowalt}Precision & истински атаки от 100 вдигнати аларми & 100 от 100 & 100 от 100 & 100 от 100 & 100 от 100 & 100 от 100 & 100 от 100 & 100 от 100 & 100 от 100 & 100 от 100\\
Recall & открити атаки от 100 реални атаки & 100 от 100 & 99 от 100 & 100 от 100 & 100 от 100 & 100 от 100 & 100 от 100 & 100 от 100 & 100 от 100 & 100 от 100\\
\rowcolor{rowalt}Specificity & нормални записи без аларма от 100 нормални & 100 от 100 & 100 от 100 & 100 от 100 & 100 от 100 & 100 от 100 & 100 от 100 & 100 от 100 & 100 от 100 & 100 от 100\\
\bottomrule
\end{tabular}}%

\vspace{2pt}
\interprettab{При закръгляне до цял случай от 100 всички модели отговарят правилно в практически всичките 100 случая. Единственото изключение е SVM при Bot: от 100 атаки на Bot \val{SVM-Bot-Recall-100} се откриват.}
\end{table}

\begin{figure}[H]
\centering
\includegraphics[width=0.86\textwidth]{04_confusion_matrices.png}
\caption{Матрици на объркване на деветте модела в оперативната фаза}\label{fig:cm}
\vspace{2pt}
\interpret{Редовете са реалният клас, а колоните предсказаният. Всяка клетка съдържа броя записи и процента от реалния клас. Диагоналът е практически 100\% навсякъде. Грешките на RF и XGBoost се броят на пръсти. SVM при Bot пропуска \val{SVM-Bot-FN} записа (\val{SVM-Bot-FNR} от атаките).}
\end{figure}

ROC кривите на RF и XGBoost са на горния ляв ръб дори при Bot (AUC \val{RF-Bot-AUC} и \val{XGB-Bot-AUC}), а SVM при Bot достига пълен Recall малко по-късно (AUC \val{SVM-Bot-AUC}).

Скоростта в експлоатация се различава съществено. При Bot XGBoost класифицира \val{XGB-Bot-throughput} записа в секунда, Random Forest \val{RF-Bot-throughput}, а SVM \val{SVM-Bot-throughput}. Класификацията със SVM изисква изчисляване на ядрото (\cref{eq:rbf}) спрямо всеки опорен вектор~\cite{chang2011}, а XGBoost използва малко на брой плитки дървета.

% ---------------------------------------------------------------------------
\section{Статистическа надеждност (A9)}\label{sec:r-stat}

\begin{table}[H]
\centering\small
\caption{Метрики с 95\% bootstrap доверителни интервали}\label{tab:ci}
\fitbox{%
\begin{tabular}{@{}llcccc@{}}
\toprule
\headrow \hd{Модел} & \hd{Атака} & \hd{Precision} & \hd{Accuracy} & \hd{Sensitivity} & \hd{F1}\\
\midrule
Random Forest & PortScan & \makecell{99,80\%\\[-1pt]{\scriptsize\color{muted}[99,37\%; 100,00\%]}} & \makecell{99,80\%\\[-1pt]{\scriptsize\color{muted}[99,50\%; 100,00\%]}} & \makecell{99,80\%\\[-1pt]{\scriptsize\color{muted}[99,39\%; 100,00\%]}} & \makecell{99,80\%\\[-1pt]{\scriptsize\color{muted}[99,50\%; 100,00\%]}}\\
\rowcolor{rowalt}SVM (RBF) & PortScan & \makecell{99,60\%\\[-1pt]{\scriptsize\color{muted}[98,99\%; 100,00\%]}} & \makecell{99,41\%\\[-1pt]{\scriptsize\color{muted}[98,91\%; 99,80\%]}} & \makecell{99,21\%\\[-1pt]{\scriptsize\color{muted}[98,37\%; 99,81\%]}} & \makecell{99,40\%\\[-1pt]{\scriptsize\color{muted}[98,90\%; 99,81\%]}}\\
XGBoost & PortScan & \makecell{99,60\%\\[-1pt]{\scriptsize\color{muted}[98,99\%; 100,00\%]}} & \makecell{99,70\%\\[-1pt]{\scriptsize\color{muted}[99,41\%; 100,00\%]}} & \makecell{99,80\%\\[-1pt]{\scriptsize\color{muted}[99,39\%; 100,00\%]}} & \makecell{99,70\%\\[-1pt]{\scriptsize\color{muted}[99,37\%; 100,00\%]}}\\
\rowcolor{rowalt}Random Forest & DDoS & \makecell{100,00\%\\[-1pt]{\scriptsize\color{muted}[100,00\%; 100,00\%]}} & \makecell{100,00\%\\[-1pt]{\scriptsize\color{muted}[100,00\%; 100,00\%]}} & \makecell{100,00\%\\[-1pt]{\scriptsize\color{muted}[100,00\%; 100,00\%]}} & \makecell{100,00\%\\[-1pt]{\scriptsize\color{muted}[100,00\%; 100,00\%]}}\\
SVM (RBF) & DDoS & \makecell{99,94\%\\[-1pt]{\scriptsize\color{muted}[99,91\%; 99,96\%]}} & \makecell{99,95\%\\[-1pt]{\scriptsize\color{muted}[99,93\%; 99,97\%]}} & \makecell{99,96\%\\[-1pt]{\scriptsize\color{muted}[99,94\%; 99,98\%]}} & \makecell{99,95\%\\[-1pt]{\scriptsize\color{muted}[99,93\%; 99,97\%]}}\\
\rowcolor{rowalt}XGBoost & DDoS & \makecell{100,00\%\\[-1pt]{\scriptsize\color{muted}[100,00\%; 100,00\%]}} & \makecell{100,00\%\\[-1pt]{\scriptsize\color{muted}[100,00\%; 100,00\%]}} & \makecell{100,00\%\\[-1pt]{\scriptsize\color{muted}[100,00\%; 100,00\%]}} & \makecell{100,00\%\\[-1pt]{\scriptsize\color{muted}[100,00\%; 100,00\%]}}\\
Random Forest & BruteForce & \makecell{100,00\%\\[-1pt]{\scriptsize\color{muted}[100,00\%; 100,00\%]}} & \makecell{100,00\%\\[-1pt]{\scriptsize\color{muted}[99,99\%; 100,00\%]}} & \makecell{100,00\%\\[-1pt]{\scriptsize\color{muted}[99,99\%; 100,00\%]}} & \makecell{100,00\%\\[-1pt]{\scriptsize\color{muted}[99,99\%; 100,00\%]}}\\
\rowcolor{rowalt}SVM (RBF) & BruteForce & \makecell{100,00\%\\[-1pt]{\scriptsize\color{muted}[99,99\%; 100,00\%]}} & \makecell{99,99\%\\[-1pt]{\scriptsize\color{muted}[99,98\%; 99,99\%]}} & \makecell{99,98\%\\[-1pt]{\scriptsize\color{muted}[99,95\%; 99,99\%]}} & \makecell{99,99\%\\[-1pt]{\scriptsize\color{muted}[99,98\%; 99,99\%]}}\\
XGBoost & BruteForce & \makecell{100,00\%\\[-1pt]{\scriptsize\color{muted}[100,00\%; 100,00\%]}} & \makecell{100,00\%\\[-1pt]{\scriptsize\color{muted}[99,99\%; 100,00\%]}} & \makecell{100,00\%\\[-1pt]{\scriptsize\color{muted}[99,99\%; 100,00\%]}} & \makecell{100,00\%\\[-1pt]{\scriptsize\color{muted}[99,99\%; 100,00\%]}}\\
\bottomrule
\end{tabular}}%
\par\smallskip{\footnotesize\color{muted}В квадратни скоби: 95\% перцентилен bootstrap интервал по формула (\ref{eq:bootstrap}), 1000 повторни извадки.}

\vspace{2pt}
\interprettab{Всяка клетка съдържа стойността на метриката и интервала, в който тя би попаднала в 95\% от повторенията на теста с нова извадка със същия размер (\cref{eq:bootstrap}). Интервалите са много тесни: полуширината за F1 е най-много $\pm\val{SVM-Bot-F1-hw}$ пункта (SVM при Bot).}
\end{table}

\begin{table}[H]
\centering\small
\caption{Сравнение на алгоритмите по двойки: тест на McNemar и сдвоени bootstrap разлики}\label{tab:mcnemar}
\fitbox{%
\begin{tabular}{@{}llrrrrrrrr@{}}
\toprule
\headrow \hd{Атака} & \hd{Двойка\\A / B} & \hd{И двата\\верни} & \hd{Само A\\верен ($b$)} & \hd{Само B\\верен ($c$)} & \hd{И двата\\грешни} & \hd{$\chi^2$} & \hd{$p$} & \hd{$\Delta$F1} & \hd{$\Delta$Sensitivity}\\
\midrule
DDoS & RF / SVM & 59\,970 & 30 & 0 & 0 & 28,03 & $1{,}2\cdot10^{-7}$ & \makecell{0,050\\[-1pt]{\scriptsize\color{muted}[0,033; 0,068]}} & \makecell{0,040\\[-1pt]{\scriptsize\color{muted}[0,020; 0,063]}}\\
\rowcolor{rowalt}DDoS & RF / XGB & 60\,000 & 0 & 0 & 0 & 0,00 & $1{,}0000$ & \makecell{0,000\\[-1pt]{\scriptsize\color{muted}[0,000; 0,000]}} & \makecell{0,000\\[-1pt]{\scriptsize\color{muted}[0,000; 0,000]}}\\
DDoS & SVM / XGB & 59\,970 & 0 & 30 & 0 & 28,03 & $1{,}2\cdot10^{-7}$ & \makecell{-0,050\\[-1pt]{\scriptsize\color{muted}[-0,068; -0,033]}} & \makecell{-0,040\\[-1pt]{\scriptsize\color{muted}[-0,063; -0,020]}}\\
\rowcolor{rowalt}BruteForce & RF / SVM & 56\,449 & 7 & 0 & 1 & 0,00 & $0{,}0156$ & \makecell{0,012\\[-1pt]{\scriptsize\color{muted}[0,005; 0,023]}} & \makecell{0,021\\[-1pt]{\scriptsize\color{muted}[0,007; 0,039]}}\\
BruteForce & RF / XGB & 56\,456 & 0 & 0 & 1 & 0,00 & $1{,}0000$ & \makecell{0,000\\[-1pt]{\scriptsize\color{muted}[0,000; 0,000]}} & \makecell{0,000\\[-1pt]{\scriptsize\color{muted}[0,000; 0,000]}}\\
\rowcolor{rowalt}BruteForce & SVM / XGB & 56\,449 & 0 & 7 & 1 & 0,00 & $0{,}0156$ & \makecell{-0,012\\[-1pt]{\scriptsize\color{muted}[-0,023; -0,005]}} & \makecell{-0,021\\[-1pt]{\scriptsize\color{muted}[-0,039; -0,007]}}\\
PortScan & RF / SVM & 1\,003 & 5 & 1 & 1 & 1,00 & $0{,}2188$ & \makecell{0,397\\[-1pt]{\scriptsize\color{muted}[0,000; 0,930]}} & \makecell{0,594\\[-1pt]{\scriptsize\color{muted}[0,000; 1,378]}}\\
\rowcolor{rowalt}PortScan & RF / XGB & 1\,007 & 1 & 0 & 2 & 0,00 & $1{,}0000$ & \makecell{0,099\\[-1pt]{\scriptsize\color{muted}[0,000; 0,304]}} & \makecell{0,000\\[-1pt]{\scriptsize\color{muted}[0,000; 0,000]}}\\
PortScan & SVM / XGB & 1\,002 & 2 & 5 & 1 & 2,00 & $0{,}4531$ & \makecell{-0,299\\[-1pt]{\scriptsize\color{muted}[-0,894; 0,198]}} & \makecell{-0,594\\[-1pt]{\scriptsize\color{muted}[-1,378; 0,000]}}\\
\bottomrule
\end{tabular}}%
\par\smallskip{\footnotesize\color{muted}$\Delta$ = A $-$ B в процентни пунктове, в скоби 95\% сдвоен bootstrap интервал.}

\vspace{2pt}
\interprettab{Колоните $b$ и $c$ съдържат записите, по които двата модела не са съгласни (\cref{eq:mcnemar}). При Bot SVM греши значимо по-често от RF ($b=\val{mcn-Bot-RF-SVM-b}$, $c=\val{mcn-Bot-RF-SVM-c}$, $p={}$\val{mcn-Bot-RF-SVM-p}) и от XGBoost ($p={}$\val{mcn-Bot-SVM-XGB-p}), а XGBoost е малко по-точен от RF ($b=\val{mcn-Bot-RF-XGB-b}$, $c=\val{mcn-Bot-RF-XGB-c}$, $p={}$\val{mcn-Bot-RF-XGB-p}). При DDoS и BruteForce RF и XGBoost дават напълно еднакви решения, а SVM прави малко повече грешки (\val{mcn-DDoS-RF-SVM-b} при DDoS и \val{mcn-BruteForce-RF-SVM-b} при BruteForce). Тези разлики са статистически значими, но са под 0,05 процентни пункта и нямат практическо значение.}
\end{table}

% ---------------------------------------------------------------------------
\section{Ефект на редукцията на размерността (E1, E5, A4, E4)}\label{sec:r-red}

\begin{table}[H]
\centering\small
\caption{Пълен срещу редуциран вектор на признаците (E1)}\label{tab:e1}
\fitbox{%
\begin{tabular}{@{}llrrrrr@{}}
\toprule
\headrow \hd{Модел} & \hd{Атака} & \hd{F1 пълен} & \hd{F1 редуциран} & \hd{Разлика, пр. п.} & \hd{Обучение, s} & \hd{µs/запис}\\
\midrule
Random Forest & DDoS & 100,000\% & 100,000\% & $+$0,00 & 5,6 / 4,9 & 3,8 / 3,8\\
\rowcolor{rowalt}SVM (RBF) & DDoS & 99,972\% & 99,950\% & $-$0,02 & 2,0 / 1,6 & 56,4 / 11,2\\
XGBoost & DDoS & 100,000\% & 100,000\% & $+$0,00 & 1,3 / 4,1 & 0,5 / 0,6\\
\rowcolor{rowalt}Random Forest & BruteForce & 99,996\% & 99,998\% & $+$0,00 & 1,4 / 3,2 & 4,0 / 4,1\\
SVM (RBF) & BruteForce & 99,977\% & 99,986\% & $+$0,01 & 0,2 / 1,1 & 3,5 / 2,5\\
\rowcolor{rowalt}XGBoost & BruteForce & 100,000\% & 99,998\% & $-$0,00 & 1,3 / 3,0 & 0,6 / 0,6\\
Random Forest & PortScan & 99,703\% & 99,802\% & $+$0,10 & 0,2 / 0,6 & 36,6 / 37,0\\
\rowcolor{rowalt}SVM (RBF) & PortScan & 99,504\% & 99,405\% & $-$0,10 & 0,0 / 0,4 & 5,4 / 4,1\\
XGBoost & PortScan & 99,703\% & 99,703\% & $+$0,00 & 0,4 / 0,7 & 3,0 / 2,4\\
\bottomrule
\end{tabular}}%
\par\smallskip{\footnotesize\color{muted}Времената са във формат „пълен / редуциран“ вектор.}

\vspace{2pt}
\interprettab{Всеки модел е обучен два пъти с едни и същи хиперпараметри: с всички \val{n-features} признака и с филтър и подбор на \val{n-selected} признака. Разликите в F1 са в рамките на $\pm0{,}03$ пункта, т.е. редукцията не губи информация. Тя ускорява класификацията със SVM \val{SVM-DDoS-speedup} пъти при DDoS и \val{SVM-Bot-speedup} пъти при Bot, защото изчисляването на ядрото е пропорционално на размерността.}
\end{table}

\begin{table}[H]
\centering\small
\caption{Подбор на признаци срещу PCA при едни и същи хиперпараметри (E5)}\label{tab:e5}
\fitbox{%
\begin{tabular}{@{}llrrrr@{}}
\toprule
\headrow \hd{Модел} & \hd{Атака} & \hd{PCA компоненти} & \hd{F1 с PCA} & \hd{F1 с подбор} & \hd{Подбор $-$ PCA, пр. п.}\\
\midrule
Random Forest & DDoS & 22 & 100,000\% & 100,000\% & $+$0,00\\
\rowcolor{rowalt}SVM (RBF) & DDoS & 21 & 99,955\% & 99,950\% & $-$0,01\\
XGBoost & DDoS & 22 & 99,998\% & 100,000\% & $+$0,00\\
\rowcolor{rowalt}Random Forest & BruteForce & 22 & 99,993\% & 99,998\% & $+$0,01\\
SVM (RBF) & BruteForce & 20 & 99,972\% & 99,986\% & $+$0,01\\
\rowcolor{rowalt}XGBoost & BruteForce & 22 & 99,991\% & 99,998\% & $+$0,01\\
Random Forest & Bot & 23 & 99,948\% & 99,967\% & $+$0,02\\
\rowcolor{rowalt}SVM (RBF) & Bot & 22 & 99,600\% & 99,597\% & $-$0,00\\
XGBoost & Bot & 23 & 99,942\% & 99,982\% & $+$0,04\\
\bottomrule
\end{tabular}}%

\vspace{2pt}
\interprettab{PCA запазва 95\% от дисперсията с \val{pca-DDoS}--\val{pca-Bot} компонента. Подборът на \val{n-selected} признака постига същото или малко по-високо F1 с по-малко измерения. Освен това запазва оригиналните признаци, които операторът може да тълкува.}
\end{table}

Анализът A4 (\texttt{results/pca\_sweep.csv}) показва, че при BruteForce два признака стигат за 99,9\%, при DDoS пет признака дават \val{ps-select-RF-DDoS-k5}, а при Bot качеството расте до около 20 признака и след това се изравнява. Затова $k=\val{n-selected}$ е разумният компромис.

Времевите признаци на потока (продължителност, интервали между пакетите, периоди на активност и бездействие) зависят от натоварването на мрежата и от машината, на която е уловен трафикът~\cite{arp2022,engelen2021}. Затова в експеримент E4 моделите са обучени повторно без тях. Промяната на F1 е между \val{SVM-Bot-noiat-dF1} (SVM при Bot) и 0,00 пункта, т.е. в рамките на доверителните интервали (\texttt{results/metrics\_without\_IAT.csv}). Моделите разчитат и на други свойства на трафика, като размер на пакетите, TCP флагове и размер на TCP прозореца, и не зависят от времевите признаци.

% ---------------------------------------------------------------------------
\section{Работа при естествено съотношение на класовете (E2, A5)}\label{sec:r-nat}

В реален трафик атаките не са 50\% от записите. В експеримент E2 записаните модели се прилагат върху записи, които не са участвали в обучението, смесени в естественото съотношение на изтеглените файлове: \val{nat-share-DDoS} DDoS, \val{nat-share-BruteForce} BruteForce и \val{nat-share-Bot} Bot (\cref{tab:e2}).

\begin{table}[H]
\centering\small
\caption{Резултати при естествено съотношение на класовете (E2)}\label{tab:e2}
\fitbox{%
\begin{tabular}{@{}llrrrrrrrr@{}}
\toprule
\headrow \hd{Модел} & \hd{Атака} & \hd{Дял атаки} & \hd{Записи} & \hd{TP} & \hd{FP} & \hd{FN} & \hd{Precision} & \hd{Sensitivity} & \hd{Specificity}\\
\midrule
Random Forest & DDoS & 27,7\% & 200\,000 & 55\,495 & 2 & 0 & 100,0\% & 100,0\% & 100,0\%\\
\rowcolor{rowalt}SVM (RBF) & DDoS & 27,7\% & 200\,000 & 55\,478 & 116 & 17 & 99,8\% & 100,0\% & 99,9\%\\
XGBoost & DDoS & 27,7\% & 200\,000 & 55\,495 & 0 & 0 & 100,0\% & 100,0\% & 100,0\%\\
\rowcolor{rowalt}Random Forest & BruteForce & 17,5\% & 160\,925 & 28\,227 & 4 & 1 & 100,0\% & 100,0\% & 100,0\%\\
SVM (RBF) & BruteForce & 17,5\% & 160\,925 & 28\,221 & 12 & 7 & 100,0\% & 100,0\% & 100,0\%\\
\rowcolor{rowalt}XGBoost & BruteForce & 17,5\% & 160\,925 & 28\,227 & 5 & 1 & 100,0\% & 100,0\% & 100,0\%\\
Random Forest & PortScan & 14,0\% & 3\,613 & 504 & 9 & 1 & 98,2\% & 99,8\% & 99,7\%\\
\rowcolor{rowalt}SVM (RBF) & PortScan & 14,0\% & 3\,613 & 501 & 37 & 4 & 93,1\% & 99,2\% & 98,8\%\\
XGBoost & PortScan & 14,0\% & 3\,613 & 504 & 13 & 1 & 97,5\% & 99,8\% & 99,6\%\\
\bottomrule
\end{tabular}}%

\vspace{2pt}
\interprettab{Recall и Specificity остават практически същите като при балансирания тест, защото не зависят от съотношението на класовете. Precision се променя по \cref{eq:baserate}, но остава висока, защото фалшивите аларми са много редки. Най-ниска е при SVM за Bot: \val{SVM-Bot-nat-Precision} (\val{SVM-Bot-nat-FP} фалшиви срещу \val{SVM-Bot-nat-TP} истински аларми). XGBoost при DDoS не вдига нито една фалшива аларма.}
\end{table}

\Cref{tab:ops} показва работните точки за най-трудната атака, \val{hard-attack}, по двата критерия. Праговете по Precision ограничават шума от аларми, а праговете по Recall ограничават пропуснатите атаки, които са най-нежеланата грешка~\cite{borovska2022,axelsson2000}.

\begin{table}[H]
\centering\small
\caption{Работни точки за Bot при естествено съотношение: праг по целева Precision или по целеви Recall}\label{tab:ops}
\fitbox{%
\begin{tabular}{@{}llrrr@{}}
\toprule
\headrow \hd{Модел} & \hd{Критерий за праг} & \hd{Precision} & \hd{Recall} & \hd{Аларми / 10\,000}\\
\midrule
Random Forest & Precision $\geq$ 90\% & 90,0\% & 100,0\% & 1530\\
\rowcolor{rowalt}Random Forest & Precision $\geq$ 95\% & 95,0\% & 100,0\% & 1450\\
Random Forest & Precision $\geq$ 99\% & 99,0\% & 100,0\% & 1392\\
\rowcolor{rowalt}Random Forest & Recall $\geq$ 90\% & 100,0\% & 97,1\% & 1338\\
Random Forest & Recall $\geq$ 95\% & 100,0\% & 97,1\% & 1338\\
\rowcolor{rowalt}SVM (RBF) & Precision $\geq$ 90\% & 90,0\% & 100,0\% & 1531\\
SVM (RBF) & Precision $\geq$ 95\% & 95,0\% & 99,8\% & 1448\\
\rowcolor{rowalt}SVM (RBF) & Precision $\geq$ 99\% & 99,0\% & 99,3\% & 1382\\
SVM (RBF) & Recall $\geq$ 90\% & 100,0\% & 90,0\% & 1241\\
\rowcolor{rowalt}SVM (RBF) & Recall $\geq$ 95\% & 100,0\% & 95,0\% & 1310\\
XGBoost & Precision $\geq$ 90\% & 90,0\% & 100,0\% & 1531\\
\rowcolor{rowalt}XGBoost & Precision $\geq$ 95\% & 95,4\% & 100,0\% & 1444\\
XGBoost & Precision $\geq$ 99\% & 99,0\% & 100,0\% & 1392\\
\rowcolor{rowalt}XGBoost & Recall $\geq$ 90\% & 100,0\% & 90,0\% & 1240\\
XGBoost & Recall $\geq$ 95\% & 100,0\% & 95,0\% & 1309\\
\bottomrule
\end{tabular}}%

\vspace{2pt}
\interprettab{При критерий по Precision се избира прагът с максимален Recall, а при критерий по Recall се избира най-високият праг, който го постига. При RF и XGBoost Precision 99\% се постига с практически пълен Recall (\val{op-RF-Bot-99-recall}). Затова не е нужен компромис между малко фалшиви аларми и много открити атаки.}
\end{table}

Анализът A7 проверява дали модел за една атака открива другите две. Всеки модел открива почти 100\% от своята атака и практически 0\% от другите. Моделите са научили конкретния отпечатък на атаката, а не общо понятие за „необичаен“ трафик, затова практическата IDS се нуждае от отделен модел за всяка категория атаки.

\section{Проверки за приложимост (A10--A14)}\label{sec:r-rob}

Почти перфектните резултати повдигат въпроса дали моделите научават поведението на атаката, или само отпечатъка на конкретния инструмент и случайното разделяне на записите. Петте проверки в този раздел отговарят на този въпрос (\cref{tab:experiments}).

\paragraph{Признаци-отпечатъци (A10).} Началният размер на TCP прозореца (\texttt{Init Fwd Win Byts}, \texttt{Init Bwd Win Byts}) и минималният размер на сегмента (\texttt{Fwd Seg Size Min}) зависят от операционната система и инструмента на изпращача. Без тези три признака точността при случайно разделяне спада най-много с \val{rob-fp-maxdrop} процентни пункта (\cref{fig:rob-split}). Моделите не разчитат само на отпечатъка: атаката се разпознава и по размера на пакетите, броя им и интервалите между тях.

\paragraph{Хронологично разделяне (A11).} Моделът се обучава върху по-ранните 70\% от потоците и се тества върху по-късните 30\%. Границата е в \val{rob-cut-DDoS} при DDoS, \val{rob-cut-BruteForce} при BruteForce и \val{rob-cut-Bot} при Bot. При DDoS и BruteForce всички модели остават над \val{rob-chrono-min-DDoSBF}. При Bot Random Forest и XGBoost спадат до \val{rob-chronological-all-RF-Bot-acc} и \val{rob-chronological-all-XGB-Bot-acc}, защото откриват само \val{rob-chronological-all-XGB-Bot-sens} от по-късните атаки. След 14:42 Bot започва да обменя потоци с различна структура (5 пакета напред и 2 обратно, отговори от около 326 байта), които не се срещат в ранните данни. SVM се справя с тях по-добре: \val{rob-chronological-all-SVM-Bot-acc}, с \val{rob-chronological-all-SVM-Bot-sens} открити атаки.

\begin{figure}[H]
\centering
\includegraphics[width=0.92\textwidth]{24_robustness_splits.png}
\caption{Точност при случайно и хронологично разделяне, с всички признаци и без признаците-отпечатъци (A10, A11, A14)}\label{fig:rob-split}
\vspace{2pt}
\interpret{Всяка клетка е Accuracy, а под нея е делът открити атаки. Премахването на признаците-отпечатъци почти не променя резултата. Хронологичното разделяне е по-трудно, но при DDoS и BruteForce моделите остават над 99,5\%. Единствените червени клетки са Random Forest и XGBoost при Bot, които пропускат около 10\% от по-късните атаки. Комбинираният детектор (последният ред на всяка атака) възстановява откриването до около 100\%.}
\end{figure}

\paragraph{Комбиниран детектор (A14).} Дърветата са най-точни върху познатите модели на атаката, а гладката граница на SVM улавя и новите. Затова е проверен детектор, който обявява атака, ако я открие XGBoost или SVM. При хронологичното разделяне за Bot той постига \val{rob-chronological-all-XGBorSVM-Bot-acc} с \val{rob-chronological-all-XGBorSVM-Bot-sens} открити атаки, при \val{rob-chronological-all-XGBorSVM-Bot-fpr} фалшиви аларми. В основната постановка (същите тестови записи като в \cref{sec:r-op}) цената е повече фалшиви аларми: при Bot Accuracy е \val{hyb-balanced-Bot-acc} (\val{hyb-balanced-Bot-fp} фалшиви и \val{hyb-balanced-Bot-fn} пропуснати), а при естествено съотношение Precision е \val{hyb-natural-Bot-prec}. Комбинацията е подходяща, когато се очакват нови разновидности на познатата атака; когато атаката е добре позната, XGBoost сам е по-прецизен.

\paragraph{Пренос между подвидовете (A12).} Модел, обучен само върху единия инструмент, почти не открива другия (\cref{fig:rob-sub}). HOIC е HTTP наводнение по TCP, а LOIC-UDP е UDP наводнение, затова двете атаки нямат общи признаци. При BruteForce Random Forest открива всички FTP записи, обучен само върху SSH, но без признаците-отпечатъци откриването пада до \val{sub-SSH-FTP-RF-no_fp}. Общият признак е началният TCP прозорец на инструмента Patator, а не поведението на атаката. Изводът за практиката е, че обучаващите данни трябва да съдържат всеки използван инструмент, както е в основния експеримент.

\begin{figure}[H]
\centering
\includegraphics[width=0.85\textwidth]{26_subtype_transfer.png}
\caption{Дял открити атаки от подвид, който не е участвал в обучението (A12)}\label{fig:rob-sub}
\vspace{2pt}
\interpret{Почти всички клетки са близо до 0\%. Изключение е Random Forest при SSH $\to$ FTP с всички признаци, но този пренос се дължи на общия отпечатък на инструмента и изчезва без него. Записите на FTP са само 53 уникални потока, затова редът FTP $\to$ SSH е показан само за пълнота.}
\end{figure}

\paragraph{Нормален трафик от непознат ден (A13).} Нормалният трафик на 21.02.2018 е необичайно еднороден: \val{benign-2102-win65535} от потоците имат начален TCP прозорец 65\,535, а \val{benign-2102-tcp} са по TCP (в другите дни TCP е около \val{benign-other-tcp}, останалото е UDP). Ако DDoS моделът е обучен само с този нормален трафик, той вдига фалшиви аларми за \val{rob-1day-RF-DDoS} (RF) до \val{rob-1day-XGB-DDoS} (XGBoost) от нормалния трафик в другите дни. С нормален трафик от два дни фалшивите аларми в третия, непознат ден падат до \val{rob-2day-RF-DDoS}--\val{rob-2day-SVM-DDoS} (\cref{fig:rob-days}). Затова основният експеримент използва нормален трафик от всички дни. Обратният случай показва границата на подхода: модел за Bot, обучен без нормалния трафик от 21.02, обявява голяма част от него за атака (Random Forest \val{rob-2day-RF-Bot}, XGBoost \val{rob-2day-XGB-Bot}), докато SVM остава на \val{rob-2day-SVM-Bot}.

\begin{figure}[H]
\centering
\includegraphics[width=\textwidth]{25_benign_days.png}
\caption{Фалшиви аларми върху нормален трафик от ден, който не е участвал в обучението (A13)}\label{fig:rob-days}
\vspace{2pt}
\interpret{Светлите стълбове са модели с нормален трафик само от деня на атаката, тъмните са с нормален трафик от два дни. Показан е по-лошият от двата непознати дни; скалата е логаритмична. При DDoS разнообразният нормален трафик намалява фалшивите аларми над сто пъти. При Bot дърветата правят много фалшиви аларми върху необичайния нормален трафик от 21.02, а SVM остава под 0,1\%.}
\end{figure}

Всички допълнителни фигури и таблици (корелационни матрици, разпределения на признаците, чувствителност към хиперпараметрите, стабилност по фолдове, време за обучение, важност на признаците, разновидности на атаките, модели без времеви признаци, пълните резултати от настройката) се генерират в \texttt{results/} на хранилището~\cite{repo}.

% ============================== 5. АНАЛИЗ ==============================
\chapter{Анализ на експерименталните резултати}\label{ch:analysis}

\section{Отговори на изследователските въпроси}

\begin{keybox}[В1. Точност на трите алгоритъма за трите атаки]
И деветте модела постигат Accuracy над 99,5\%. За DDoS тя е \val{RF-DDoS-Accuracy} (RF), \val{SVM-DDoS-Accuracy} (SVM) и \val{XGB-DDoS-Accuracy} (XGBoost). За BruteForce е \val{RF-BruteForce-Accuracy}, \val{SVM-BruteForce-Accuracy} и \val{XGB-BruteForce-Accuracy}, а за Bot \val{RF-Bot-Accuracy}, \val{SVM-Bot-Accuracy} и \val{XGB-Bot-Accuracy}. Най-трудна е атаката Bot, а най-слаб е SVM. Двата ансамбъла от дървета са практически безгрешни при всички атаки. По учебния материал Accuracy е подходяща тук, защото тестът е балансиран~\cite{borovska2022}. Същото показват и F1, и Specificity.
\end{keybox}

\begin{keybox}[В2. Ефект на редукцията на размерността]
Двуетапната редукция свива вектора от \val{n-features} до \val{n-selected} признака (с 74\%), без да губи качество: разликите в F1 спрямо пълния вектор са в рамките на $\pm0{,}03$ пункта (E1). Подборът на признаци е поне толкова добър, колкото PCA, с по-малко измерения, и запазва интерпретируемите признаци (E5). Редукцията ускорява класификацията със SVM \val{SVM-DDoS-speedup} пъти при DDoS. Анализът A4 показва, че пет признака стигат за DDoS и BruteForce, а за Bot са нужни около 20.
\end{keybox}

\begin{keybox}[В3. Статистическа значимост и стабилност. Преобучение]
Разсейването по фолдове и bootstrap интервалите са много малки (полуширина на F1 под $0{,}05$ пункта). Тестът на McNemar показва значима разлика в \val{mcn-n-sig} от \val{mcn-n} двойки. Практически значима е само разликата на SVM при Bot: \val{mcn-Bot-RF-SVM-b} записа, които RF класифицира вярно, а SVM грешно, срещу \val{mcn-Bot-RF-SVM-c} обратни ($p={}$\val{mcn-Bot-RF-SVM-p}). Преобучение няма: разликата между точността при обучение и при тест е най-много \val{gap-max} пункта, а оценката от кръстосаната валидация съвпада с тестовата (\cref{tab:overfit}).
\end{keybox}

\begin{keybox}[В4. Качество при реално съотношение и избор на праг]
При естествения дял на атаките (\val{nat-share-Bot}--\val{nat-share-DDoS}) Precision остава между \val{SVM-Bot-nat-Precision} (SVM при Bot) и 100\%. Проблемът на базовата честота (\cref{eq:baserate}) не се проявява, защото делът на фалшивите аларми е от порядъка на $10^{-4}$--$10^{-3}$, много по-малък от дела на атаките. При RF и XGBoost Precision 99\% се постига с практически пълен Recall, така че изборът на праг не изисква компромис.
\end{keybox}

\begin{keybox}[В5. Носители на информацията и от какво зависи точността]
BruteForce се разпознава по структурата на кратките сесии за вход: \val{ks-top5names-BruteForce} (средно $D=\val{ks-top5-BruteForce}$). DDoS се разпознава по размера на отговорите на сървъра и по TCP прозореца: \val{ks-top5names-DDoS}. Bot се разпознава по размера на пакетите и по интервалите между тях: \val{ks-top5names-Bot}, а пермутационната важност извежда и флага RST. И при трите атаки признаците на отделния поток носят отпечатъка на атаката. При Port Scan в CICIoT2023 не е така, затова там точността остава около 85\%. Постижимата точност се определя преди всичко от данните, а не от алгоритъма.
\end{keybox}

\begin{keybox}[В6. Приложимост извън случайното разделяне]
Точността не идва само от отпечатъка на инструмента: без признаците на TCP стека тя спада най-много с \val{rob-fp-maxdrop} пункта (A10). При хронологично разделяне DDoS и BruteForce остават над \val{rob-chrono-min-DDoSBF}, но при Bot Random Forest и XGBoost пропускат около 10\% от по-късните атаки, докато SVM открива \val{rob-chronological-all-SVM-Bot-sens} (A11). Комбинираният детектор XGBoost или SVM възстановява откриването до \val{rob-chronological-all-XGBorSVM-Bot-sens} (A14). Моделите не се пренасят между инструменти с различен протокол (A12), а фалшивите аларми зависят силно от това дали нормалният трафик при обучението е разнообразен (A13).
\end{keybox}

\section{Защо трите атаки се откриват толкова точно}
Резултатите се обясняват с начина, по който всяка атака променя отделния мрежов поток. Опитът за вход при BruteForce е кратка сесия с еднаква структура: еднакъв брой пакети, еднакъв размер на заглавките и на полезните данни, еднакви TCP опции. Нормалният потребител рядко създава хиляди такива еднакви сесии. Затова признак като \texttt{Fwd Seg Size Min} сам разделя класовете. DDoS с HOIC изпраща HTTP заявки с типичен размер, а сървърът отговаря с типичен размер. Флагове и размер на TCP прозореца, характерни за инструмента, се срещат рядко в нормалния трафик. Botnet Ares поддържа периодична връзка с команден сървър: малки, еднакви по размер заявки и отговори през регулярни интервали, а броят на пакетите с флаг RST е сред най-важните признаци за моделите.

Сканирането на портове в CICIoT2023 е обратният случай. Отделната проба към един порт е нормална кратка TCP сесия. Аномалията се вижда само при агрегиране по източник, когато един адрес се свързва с много различни портове за кратко време. Признак от този тип липсва в записите. Това сравнение е основният методологичен извод на изследването: преди да се търси по-добър алгоритъм, трябва да се провери дали признаците на записа изобщо съдържат информация за атаката.

\section{Сравнение на алгоритмите}
\begin{table}[H]
\centering\small
\caption{Обобщено сравнение на Random Forest, SVM и XGBoost в изследването}\label{tab:algo-compare}
\begin{tabularx}{\textwidth}{@{}l L L L@{}}
\toprule
\textsf{\textbf{Критерий}} & \textsf{\textbf{Random Forest}} & \textsf{\textbf{SVM (RBF)}} & \textsf{\textbf{XGBoost}}\\
\midrule
Качество при DDoS и BruteForce & практически безгрешен & над 99,9\% & практически безгрешен\\
Качество при Bot & \val{RF-Bot-Accuracy} & \val{SVM-Bot-Accuracy} & \val{XGB-Bot-Accuracy}\\
Мащабируемост на обучението & почти линейна, паралелизира се & $O(n^2)$--$O(n^3)$, нужна е подизвадка & почти линейна, хистограмен алгоритъм\\
Скорост в експлоатация (Bot) & \val{RF-Bot-throughput} записа/s & \val{SVM-Bot-throughput} записа/s & \val{XGB-Bot-throughput} записа/s\\
Нужда от мащабиране & не & задължително (\cref{eq:zscore}) & не\\
Защита от преобучение & bagging, листа с мин. брой записи & марж, параметър $C$ & регуляризация $\Omega$, subsample\\
Обяснимост & MDI и пермутационна важност & само пермутационна важност & gain и пермутационна важност\\
\bottomrule
\end{tabularx}
\end{table}

За практическа IDS в център за данни XGBoost е най-добрият избор. Той е поне толкова точен, колкото Random Forest, класифицира около пет пъти по-бързо от него и от 2,6 до 17 пъти по-бързо от SVM. Освен това дава обяснима важност на признаците, която тук се използва и за подбора им. Random Forest е почти равностоен и по-прост за настройка. SVM е най-слабият избор при случайно разделяне: по-ниска точност при Bot, по-бавна класификация и необходимост от подизвадка при обучението. Проверките за приложимост обаче показват и силна страна на SVM: гладката граница го прави по-устойчив при промяна на данните. Той открива новите потоци на Bot, които дърветата пропускат (A11), и не вдига масово фалшиви аларми върху необичаен нормален трафик (A13). Затова за Bot препоръчваме комбинирания детектор XGBoost или SVM, когато се очакват нови разновидности на атаката. Част от изоставането на SVM при Bot се дължи на по-малката обучаваща извадка (\num{50000} срещу \val{train-rows-Bot} записа), но кривата на обучение показва, че и при равен брой записи SVM е под ансамблите (\texttt{results/learning\_curves.csv}).

\section{Редукция на размерността и практическа приложимост}
Подборът на признаци е по-подходящ от PCA за IDS по три причини. Първо, той отчита етикета и запазва признаците, които различават атаката, докато PCA запазва дисперсията, която може да е безполезна. Второ, запазените признаци са оригинални величини (размер на пакет, флаг, TCP прозорец), които операторът разбира. Трето, IDS трябва да изчислява само \val{n-selected} от \val{n-features} признака, което намалява натоварването на сензора. Филтърът по дисперсия е полезен винаги, защото CICFlowMeter-V3 записва някои признаци винаги с нула~\cite{guyon2003}.

Експериментът E2 показва, че при атаки, които са 14--28\% от трафика, и при много малък дял фалшиви аларми, моделите остават полезни и при естествено съотношение. В реален център за данни делът на атаките обикновено е много по-малък. Тогава дори FPR от порядъка на $10^{-4}$ може да даде повече фалшиви, отколкото истински аларми (\cref{eq:baserate}). Практически мерки са изборът на праг по бюджет от аларми (\cref{tab:ops}), агрегирането на алармите по източник във времеви прозорец, отделен модел за всяка категория атаки, тъй като A7 показва, че моделите не се пренасят между атаките, и периодично преобучение с нови данни~\cite{axelsson2000,sommer2010}.

\section{Заплахи за валидността и насоки за развитие}
\paragraph{Вътрешна валидност.} Дубликатите и противоречивите вектори са премахнати преди разделянето, а цялата обработка се обучава само върху обучаващите данни. Признаците \texttt{Timestamp} и \texttt{Dst Port} са изключени, за да не може моделът да разпознава атаката по момента или по порта. Основният тест използва случайно разделяне на записите от едни и същи дни. Затова са добавени хронологично разделяне (A11) и тест на нормален трафик от непознат ден (A13), които са по-близо до реалната експлоатация~\cite{arp2022}. Те потвърждават резултатите за DDoS и BruteForce, а за Bot показват спад при ансамблите, който се компенсира с комбинирания детектор.

\paragraph{Външна валидност и валидност на конструкта.} Атаките в CSE-CIC-IDS2018 са проведени с конкретни инструменти (HOIC, LOIC, Patator, Ares). Част от разделящите признаци, като началния размер на TCP прозореца и минималния размер на сегмента, отразяват мрежовия стек на атакуващите машини, а не самата атака. Друг инструмент или друга операционна система могат да променят тези стойности. Проверката A10 показва, че без тези признаци точността почти не се променя, но A12 показва, че модел, обучен върху един инструмент, не открива друг инструмент с различен протокол. Известни са и грешки в етикетите и в извличането на потоците в наборите на CIC~\cite{engelen2021,liu2022}. Тук те са ограничени с премахването на противоречивите вектори, но не могат да се изключат напълно. Затова резултатите трябва да се четат като горна граница за откриването на тези инструменти, а не като гаранция за всяка атака от същия тип.

\paragraph{Насоки за развитие.} Най-важните посоки са: тест върху атаки от същия тип, проведени в други дни или с други инструменти; тест върху друг реален набор (например CIC-IDS2017 или CICIoT2023) за оценка на преносимостта; многокласова класификация на всички 14 атаки в CSE-CIC-IDS2018; обучение на целия набор (около 16 милиона потока) с разпределени инструменти за големи данни~\cite{leevy2020}; и нови признаци, агрегирани по източник, които биха позволили откриването и на Port Scan~\cite{neto2023,sharafaldin2018}.

% ============================== ЗАКЛЮЧЕНИЕ ==============================
\chapter*{Заключение}
\addcontentsline{toc}{chapter}{Заключение}
\markboth{Заключение}{}

В изследването са изградени, обучени и експериментално оценени девет модела за откриване на прониквания в мрежовия трафик на център за данни: Random Forest, SVM с RBF ядро и XGBoost за DDoS, за Brute Force срещу FTP и SSH (BruteForce) и за Botnet (Bot). Всички етапи използват само реални записи от CSE-CIC-IDS2018, уловени в облачна инфраструктура на AWS: \val{rows-read} изтеглени и \val{rows-clean} след почистване. Векторът на признаците е редуциран в два етапа от \val{n-features} до \val{n-selected} признака. Моделите са оценени по TP, TN, FP, FN, Accuracy, Precision, Recall (Sensitivity), F1 и Specificity с тълкуване „на 100 случая“, допълнени с ROC-AUC, доверителни интервали, тест на McNemar и явна проверка за преобучение.

Всички девет модела постигат Accuracy над 99,5\%. При DDoS и BruteForce Random Forest и XGBoost са практически безгрешни. При Bot, който е най-трудната от трите атаки, XGBoost постига \val{XGB-Bot-Accuracy}, Random Forest \val{RF-Bot-Accuracy}, а SVM \val{SVM-Bot-Accuracy}. Разликите между точността при обучение, при кръстосана валидация и при тест са под \val{gap-max} пункта, т.е. моделите не са преобучени. За това допринасят регуляризацията на алгоритмите и правилото „една стандартна грешка“ при избора на хиперпараметри.

Предварителното изследване върху IoT набора CICIoT2023 показа, че Port Scan там не може да се открие с точност над около 85\%, независимо от алгоритъма и от методите за подобряване. Причината е, че признаците на отделния поток не съдържат отпечатъка на Port Scan. Сравнението на двата набора е основният извод: качеството на ML-базираната IDS се определя преди всичко от това дали представянето на трафика носи информация за атаката, а едва след това от избора на алгоритъм.

Сред алгоритмите XGBoost е най-подходящ за практическа IDS: той е най-точен, класифицира над милион записа в секунда при DDoS и дава обяснима важност на признаците. Random Forest е почти равностоен. SVM изостава по точност при Bot и по скорост и изисква подизвадка при обучение. Подборът на признаци по важност е по-подходящ от PCA: запазва същото качество с по-малко измерения и с интерпретируеми признаци.

Почти перфектните резултати са проверени с пет допълнителни анализа за приложимост. Без признаците, които отразяват TCP стека на атакуващия инструмент, точността почти не се променя. При хронологично разделяне DDoS и BruteForce остават над \val{rob-chrono-min-DDoSBF}. При Bot ансамблите от дървета пропускат около 10\% от по-късните атаки, а комбинираният детектор XGBoost или SVM ги открива почти всички (\val{rob-chronological-all-XGBorSVM-Bot-sens}). Анализите показват и две условия за практическото използване: обучаващите данни трябва да съдържат всеки инструмент на атаката и разнообразен нормален трафик от различни периоди, иначе фалшивите аларми върху непознат нормален трафик могат да достигнат десетки проценти.

Програмната реализация е декларативна и напълно възпроизводима. Атаките и алгоритмите са записи в регистри, всички числа в доклада се генерират от резултатите, а коректността на конвейера се проверява с автоматични тестове върху реални данни. Затова изследването може да се разшири с нови атаки, алгоритми, признаци и набори от данни с минимални промени.

\appendix
% ============================== ПРИЛОЖЕНИЕ ==============================
\chapter{Възпроизвеждане на изследването}\label{app:repro}
\begin{enumerate}
  \item Изтегляне на трите дневни CSV файла на CSE-CIC-IDS2018 от публичното хранилище в AWS (без регистрация)~\cite{aws2018} в \texttt{data/raw/} и на архива на коригирания CIC-IDS2017~\cite{liu2022web}, от който трите дневни файла се разархивират в \texttt{data/raw\_cicids2017/} (вж. \texttt{data/README.md}).
  \item Инсталиране на зависимостите: \texttt{pip install -r requirements.txt}.
  \item Изпълнение на тестовете: \texttt{python -m pytest tests}.
  \item Пълен експеримент: \texttt{python run\_all.py} (около \val{run-minutes} минути за основния експеримент и още \val{rob-minutes} за проверките за приложимост на машината от \cref{tab:env}).
  \item Доклад: \texttt{python report/build\_report.py -{}-no-pdf} (таблици и стойности), след това \texttt{python report/build\_docx.py} (Word и PDF).
  \item Презентация: първо \path{python presentation/render_formulas.py},\\ след това \path{node presentation/build_presentation.js}.
\end{enumerate}

\begin{table}[H]
\centering\small
\caption{Структура на хранилището}\label{tab:repo}
\begin{tabularx}{\textwidth}{@{}lL@{}}
\toprule
\textsf{\textbf{Път}} & \textsf{\textbf{Съдържание}}\\
\midrule
\texttt{src/config.py}, \texttt{src/registry.py} & параметри на експеримента; регистри на атаки и алгоритми\\
\texttt{src/preprocess.py} & зареждане, почистване, бинарни задачи, \texttt{FeatureFilter}, \texttt{ImportanceSelector}\\
\texttt{src/train.py}, \texttt{src/evaluate.py} & конвейер, GridSearchCV с правилото 1-SE, оперативна фаза, метрики\\
\texttt{src/analysis.py}, \texttt{src/stats.py} & разширени анализи A1--A9, bootstrap, McNemar\\
\texttt{src/robustness.py} & проверки за приложимост A10--A14\\
\texttt{src/experiment.py}, \texttt{run\_all.py} & оркестрация и команден ред\\
\texttt{src/plots.py}, \texttt{src/plots\_ext.py}, \texttt{src/diagram.py} & всички фигури\\
\texttt{tests/} & автоматични тестове върху реални записи\\
\texttt{results/} & метрики (CSV/JSON), журнал \texttt{run\_log.txt}, фигури\\
\texttt{report/}, \texttt{presentation/} & генератори на доклада (LaTeX) и презентацията\\
\bottomrule
\end{tabularx}
\end{table}

% източниците са след приложението, но запазват номер 6 като в заданието
\makeatletter\renewcommand{\@chapapp}{\chaptername}\makeatother
\renewcommand{\thechapter}{\arabic{chapter}}\setcounter{chapter}{5}
% ============================== 6. ИНФОРМАЦИОННИ ИЗТОЧНИЦИ ==============================
\chapter{Информационни източници}\label{ch:refs}
\begingroup
\makeatletter
\renewcommand{\chapter}{\@ifstar\@gobble\@gobble}
\makeatother
\small
\begin{thebibliography}{99}\setlength{\itemsep}{3pt}
\bibitem{antonakakis2017} Antonakakis, M. et al. Understanding the Mirai Botnet. В: \emph{26th USENIX Security Symposium}, с.~1093--1110, 2017.
\bibitem{sommer2010} Sommer, R., Paxson, V. Outside the Closed World: On Using Machine Learning for Network Intrusion Detection. В: \emph{IEEE Symposium on Security and Privacy}, с.~305--316, 2010.
\bibitem{buczak2016} Buczak, A.~L., Guven, E. A Survey of Data Mining and Machine Learning Methods for Cyber Security Intrusion Detection. \emph{IEEE Communications Surveys \& Tutorials}, 18(2):1153--1176, 2016.
\bibitem{arp2022} Arp, D. et al. Dos and Don'ts of Machine Learning in Computer Security. В: \emph{31st USENIX Security Symposium}, с.~3971--3988, 2022.
\bibitem{mirkovic2004} Mirkovic, J., Reiher, P. A Taxonomy of DDoS Attack and DDoS Defense Mechanisms. \emph{ACM SIGCOMM Computer Communication Review}, 34(2):39--53, 2004.
\bibitem{sharafaldin2018} Sharafaldin, I., Lashkari, A.~H., Ghorbani, A.~A. Toward Generating a New Intrusion Detection Dataset and Intrusion Traffic Characterization. В: \emph{ICISSP}, с.~108--116, 2018.
\bibitem{tavallaee2009} Tavallaee, M., Bagheri, E., Lu, W., Ghorbani, A.~A. A Detailed Analysis of the KDD CUP 99 Data Set. В: \emph{IEEE Symposium on Computational Intelligence for Security and Defense Applications}, 2009.
\bibitem{cse2018web} Canadian Institute for Cybersecurity. A Realistic Cyber Defense Dataset (CSE-CIC-IDS2018). University of New Brunswick. \url{https://www.unb.ca/cic/datasets/ids-2018.html} (посетен на 13.11.2026).
\bibitem{cic2017web} Canadian Institute for Cybersecurity. Intrusion Detection Evaluation Dataset (CIC-IDS2017). University of New Brunswick. \url{https://www.unb.ca/cic/datasets/ids-2017.html} (посетен на 15.11.2026).
\bibitem{neto2023} Neto, E.~C.~P., Dadkhah, S., Ferreira, R., Zohourian, A., Lu, R., Ghorbani, A.~A. CICIoT2023: A Real-Time Dataset and Benchmark for Large-Scale Attacks in IoT Environment. \emph{Sensors}, 23(13):5941, 2023. \url{https://doi.org/10.3390/s23135941}
\bibitem{cic2023web} Canadian Institute for Cybersecurity. CIC IoT Dataset 2023. University of New Brunswick. \url{https://www.unb.ca/cic/datasets/iotdataset-2023.html} (посетен на 06.11.2026).
\bibitem{lyon2009} Lyon, G.~F. \emph{Nmap Network Scanning: The Official Nmap Project Guide to Network Discovery and Security Scanning}. Insecure.Com LLC, 2009.
\bibitem{aws2018} Registry of Open Data on AWS. A Realistic Cyber Defense Dataset (CSE-CIC-IDS2018). \url{https://registry.opendata.aws/cse-cic-ids2018/} (посетен на 13.11.2026).
\bibitem{leevy2020} Leevy, J.~L., Khoshgoftaar, T.~M. A Survey and Analysis of Intrusion Detection Models Based on CSE-CIC-IDS2018 Big Data. \emph{Journal of Big Data}, 7:104, 2020.
\bibitem{engelen2021} Engelen, G., Rimmer, V., Joosen, W. Troubleshooting an Intrusion Detection Dataset: the CICIDS2017 Case Study. В: \emph{IEEE Security and Privacy Workshops}, с.~7--12, 2021.
\bibitem{liu2022} Liu, L., Engelen, G., Lynar, T., Essam, D., Joosen, W. Error Prevalence in NIDS Datasets: A Case Study on CIC-IDS-2017 and CSE-CIC-IDS-2018. В: \emph{IEEE Conference on Communications and Network Security (CNS)}, с.~254--262, 2022.
\bibitem{liu2022web} Liu, L., Engelen, G., Lynar, T., Essam, D., Joosen, W. Improved CIC-IDS-2017 and CSE-CIC-IDS-2018 Datasets. DistriNet, KU Leuven. \url{https://intrusion-detection.distrinet-research.be/CNS2022/} (посетен на 15.11.2026).
\bibitem{borovska2022} Боровска, П. Метрика за оценка на качеството на моделите и правилата за машинно обучение. Учебен материал към курса „Интелигентни системи за сигурност на екосистемите на големите данни и Интернет на нещата“, 2022.
\bibitem{guyon2003} Guyon, I., Elisseeff, A. An Introduction to Variable and Feature Selection. \emph{Journal of Machine Learning Research}, 3:1157--1182, 2003.
\bibitem{jolliffe2016} Jolliffe, I.~T., Cadima, J. Principal Component Analysis: a Review and Recent Developments. \emph{Philosophical Transactions of the Royal Society A}, 374:20150202, 2016.
\bibitem{pedregosa2011} Pedregosa, F. et al. Scikit-learn: Machine Learning in Python. \emph{Journal of Machine Learning Research}, 12:2825--2830, 2011.
\bibitem{hastie2009} Hastie, T., Tibshirani, R., Friedman, J. \emph{The Elements of Statistical Learning}. 2nd ed., Springer, 2009.
\bibitem{breiman1984} Breiman, L., Friedman, J., Olshen, R., Stone, C. \emph{Classification and Regression Trees}. Wadsworth, 1984.
\bibitem{breiman2001} Breiman, L. Random Forests. \emph{Machine Learning}, 45(1):5--32, 2001.
\bibitem{chen2016} Chen, T., Guestrin, C. XGBoost: A Scalable Tree Boosting System. В: \emph{Proceedings of the 22nd ACM SIGKDD International Conference on Knowledge Discovery and Data Mining}, с.~785--794, 2016.
\bibitem{friedman2001} Friedman, J.~H. Greedy Function Approximation: A Gradient Boosting Machine. \emph{Annals of Statistics}, 29(5):1189--1232, 2001.
\bibitem{cortes1995} Cortes, C., Vapnik, V. Support-Vector Networks. \emph{Machine Learning}, 20(3):273--297, 1995.
\bibitem{scholkopf2002} Schölkopf, B., Smola, A.~J. \emph{Learning with Kernels}. MIT Press, 2002.
\bibitem{chang2011} Chang, C.-C., Lin, C.-J. LIBSVM: A Library for Support Vector Machines. \emph{ACM Transactions on Intelligent Systems and Technology}, 2(3):27, 2011.
\bibitem{kohavi1995} Kohavi, R. A Study of Cross-Validation and Bootstrap for Accuracy Estimation and Model Selection. В: \emph{IJCAI}, с.~1137--1145, 1995.
\bibitem{fawcett2006} Fawcett, T. An Introduction to ROC Analysis. \emph{Pattern Recognition Letters}, 27(8):861--874, 2006.
\bibitem{saito2015} Saito, T., Rehmsmeier, M. The Precision-Recall Plot Is More Informative than the ROC Plot When Evaluating Binary Classifiers on Imbalanced Datasets. \emph{PLoS ONE}, 10(3):e0118432, 2015.
\bibitem{axelsson2000} Axelsson, S. The Base-Rate Fallacy and the Difficulty of Intrusion Detection. \emph{ACM Transactions on Information and System Security}, 3(3):186--205, 2000.
\bibitem{efron1993} Efron, B., Tibshirani, R.~J. \emph{An Introduction to the Bootstrap}. Chapman \& Hall, 1993.
\bibitem{mcnemar1947} McNemar, Q. Note on the Sampling Error of the Difference Between Correlated Proportions or Percentages. \emph{Psychometrika}, 12(2):153--157, 1947.
\bibitem{dietterich1998} Dietterich, T.~G. Approximate Statistical Tests for Comparing Supervised Classification Learning Algorithms. \emph{Neural Computation}, 10(7):1895--1923, 1998.
\bibitem{repo} Програмен код и резултати на изследването: хранилище \texttt{ml-ids-cicids2018} (код, конфигурация, журнали и фигури), 2026. \url{https://github.com/Vankata03/ml-ids-cicids2018} (посетен на 15.11.2026).
\bibitem{fisher2019} Fisher, A., Rudin, C., Dominici, F. All Models are Wrong, but Many are Useful: Learning a Variable's Importance by Studying an Entire Class of Prediction Models Simultaneously. \emph{Journal of Machine Learning Research}, 20(177):1--81, 2019.
\bibitem{vandermaaten2008} van der Maaten, L., Hinton, G. Visualizing Data using t-SNE. \emph{Journal of Machine Learning Research}, 9:2579--2605, 2008.
\end{thebibliography}
\endgroup

\end{document}
